% Angles inscrits — Mathématiques 3ème — Cours signé OnBuch+
% Programme officiel MINESEC. Compiler avec : tectonic N12-angles-inscrits.tex
\newcommand{\DOCMATIERE}{Mathématiques}
\newcommand{\DOCNIVEAU}{3ème}
\newcommand{\DOCMODULE}{Mathématiques – classe de 3ème}
\newcommand{\DOCLECON}{N12}
\newcommand{\DOCTITRE}{Angles inscrits}
% OnBuch+ — préambule « cours signés » ENRICHI (compilé avec tectonic / XeLaTeX)
% Système visuel « Pop » d'OnBuch+ : fond crème (jamais blanc), bordures encre
% épaisses, ombres « dures » décalées, titres Archivo Black en capitales.
% Version MATHÉMATIQUES : graphes (pgfplots/tikz), boîtes pédagogiques
% (définition, propriété, méthode, attention, le savais-tu, exercice résolu,
% exercice, TP, bilan), page de garde et signature OnBuch+ ancrée en bas.
%
% Macros attendues dans main.tex AVANT \input{preamble} :
%   \DOCMATIERE  \DOCNIVEAU  \DOCMODULE  \DOCLECON (numéro)  \DOCTITRE
\documentclass[11pt]{article}

\usepackage{fontspec}
\usepackage[french]{babel}
\usepackage[a4paper,margin=2cm,top=3.4cm,bottom=2.8cm,headheight=1.4cm,footskip=1.3cm]{geometry}
\usepackage{amsmath,amssymb,amsfonts}
\usepackage{mathtools} % \xmapsto, \coloneqq... souvent utilisés par les modèles
\usepackage{cancel} % simplifications barrées
% \abs{z} (module, valeur absolue) et \norm{u} (norme), utilisés par les modèles
\providecommand{\abs}{}\let\abs\relax\DeclarePairedDelimiter{\abs}{\lvert}{\rvert}
\providecommand{\norm}{}\let\norm\relax\DeclarePairedDelimiter{\norm}{\lVert}{\rVert}
\usepackage[table]{xcolor} % [table] : fournit aussi \rowcolors (alternance de lignes)
\usepackage{pagecolor}
\usepackage{tikz}
\usetikzlibrary{arrows.meta,positioning,calc,shapes.geometric,shapes.misc,shapes.arrows,decorations.pathmorphing,decorations.pathreplacing,decorations.markings,patterns,shadows,mindmap,trees,fit,backgrounds,matrix,intersections,angles,quotes}
\usepackage{pgfplots}
\usepackage{pgf-pie} % diagrammes circulaires \pie (statistiques)
\pgfplotsset{compat=1.18}
\usepgfplotslibrary{fillbetween,groupplots} % aires entre courbes (calcul intégral)
\usepackage{enumitem}
\usepackage{fancyhdr}
% [most] charge toutes les bibliothèques tcolorbox (listings, minted, raster,
% documentation...) et gonfle inutilement la mémoire TeX sur un document long
% (~300 boîtes) : on ne charge que ce qui est réellement utilisé.
\usepackage[skins,breakable]{tcolorbox}
\usepackage{graphicx}
\usepackage{colortbl}
\usepackage{booktabs}
\usepackage{tabularx}
\usepackage{diagbox} % case d'en-tête diagonale des tableaux à double entrée
% Colonnes extensibles centrée / gauche / droite (C, L, R), souvent utilisées par les modèles
\newcolumntype{C}{>{\centering\arraybackslash}X}
\newcolumntype{L}{>{\raggedright\arraybackslash}X}
\newcolumntype{R}{>{\raggedleft\arraybackslash}X}
\usepackage{array}
\usepackage{multicol}
\usepackage{siunitx}
% parse-numbers=false : accepte aussi \SI{2,5 \times 10^{-3}}{...}
\sisetup{locale=FR,output-decimal-marker={,},group-minimum-digits=4,parse-numbers=false}
\DeclareSIUnit\ohm{\ensuremath{\symup{\Omega}}} % U+2126 absent de la police
\usepackage{qrcode}
% \degree : commande fréquemment utilisée par les modèles pour un angle ou une
% température en dehors de \SI{}{} (non fournie par les paquets chargés).
\providecommand{\degree}{^{\circ}}
% \qed (fin de démonstration), utilisé par les modèles sans amsthm
\providecommand{\qed}{\hfill$\square$}
% \barre : double barre des valeurs interdites dans un tableau de variations (tkz-tab)
\providecommand{\barre}{\ensuremath{\|}}
% environnement proof (amsthm non chargé) utilisé par les modèles
\ifdefined\proof\else\newenvironment{proof}[1][Démonstration]{\par\noindent\textit{#1.}\ }{\qed\par}\fi
\usepackage{lastpage}
\usepackage{hyperref}
\hypersetup{hidelinks}

% ============================================================
% Palette « Pop » OnBuch+ — mêmes hex que la classe P (plus_ui.dart)
% ============================================================
\definecolor{poporange}{HTML}{F59321}
\definecolor{popdark}{HTML}{E07A0C}
\definecolor{popcream}{HTML}{FFF6E8}
\definecolor{popink}{HTML}{1C1714}
\definecolor{poppurple}{HTML}{7A5AE0}
\definecolor{popgreen}{HTML}{2E9E6A}
\definecolor{poppink}{HTML}{E0457B}
\definecolor{popblue}{HTML}{2F7DE1}
\definecolor{popgold}{HTML}{C98A12}
\definecolor{popmuted}{HTML}{8A7F76}
\definecolor{popline}{HTML}{EADFD2}
\definecolor{popgray}{HTML}{8A7F76}
\colorlet{popgrey}{popgray}
% Alias tolérants (le modèle nomme parfois les couleurs différemment)
\colorlet{popwhite}{white}
\colorlet{popblack}{popink}
\colorlet{popred}{poppink}
\colorlet{popyellow}{popgold}
% Teintes claires pour fonds de tableaux / zones
\colorlet{popblueL}{popblue!10!white}
\colorlet{poporangeL}{poporange!14!white}
\colorlet{popgreenL}{popgreen!12!white}
\colorlet{poppurpleL}{poppurple!12!white}
\colorlet{poppinkL}{poppink!12!white}
\colorlet{popgoldL}{popgold!14!white}

\pagecolor{popcream}

% ============================================================
% Typographie
% ============================================================
\defaultfontfeatures{Ligatures=TeX}
\setmainfont{PlusJakartaSans-Medium.ttf}[Path=../../../fonts/,BoldFont=PlusJakartaSans-Bold.ttf]
% Maths en sans-serif (Fira Math) avec chiffres et lettres droites de Plus
% Jakarta, pour que les formules se fondent dans le texte.
\usepackage[math-style=ISO,bold-style=ISO]{unicode-math}
\setmathfont{FiraMath-Regular.otf}[Path=../../../fonts/]
\setmathfont{PlusJakartaSans-Medium.ttf}[Path=../../../fonts/,range={up/num,up/{Latin,latin},\mathup}]
\usepackage{icomma} % virgule décimale sans espace en mode math
% Caractères mathématiques Unicode absents des polices de texte (Plus Jakarta,
% Archivo Black) : rendus par leur équivalent mathématique, en texte comme en
% formule — sinon ils s'affichent en carré vide (ex. « Arithmétique dans ℤ »).
\usepackage{newunicodechar}
\newunicodechar{ℝ}{\ensuremath{\mathbb{R}}}
\newunicodechar{ℤ}{\ensuremath{\mathbb{Z}}}
\newunicodechar{ℂ}{\ensuremath{\mathbb{C}}}
\newunicodechar{ℕ}{\ensuremath{\mathbb{N}}}
\newunicodechar{ℚ}{\ensuremath{\mathbb{Q}}}
\newunicodechar{𝔻}{\ensuremath{\mathbb{D}}}
\newunicodechar{α}{\ensuremath{\alpha}}
\newunicodechar{β}{\ensuremath{\beta}}
\newunicodechar{γ}{\ensuremath{\gamma}}
\newunicodechar{δ}{\ensuremath{\delta}}
\newunicodechar{ε}{\ensuremath{\varepsilon}}
\newunicodechar{θ}{\ensuremath{\theta}}
\newunicodechar{λ}{\ensuremath{\lambda}}
\newunicodechar{μ}{\ensuremath{\mu}}
\newunicodechar{σ}{\ensuremath{\sigma}}
\newunicodechar{τ}{\ensuremath{\tau}}
\newunicodechar{φ}{\ensuremath{\varphi}}
\newunicodechar{ψ}{\ensuremath{\psi}}
\newunicodechar{ω}{\ensuremath{\omega}}
\newunicodechar{Σ}{\ensuremath{\Sigma}}
% majuscules produites par \MakeUppercase dans les titres (α-aminés -> Α)
\newunicodechar{Α}{\ensuremath{\alpha}}
\newunicodechar{Β}{\ensuremath{\beta}}
\newunicodechar{Γ}{\ensuremath{\Gamma}}
\newunicodechar{Δ}{\ensuremath{\Delta}}
\newunicodechar{Ε}{\ensuremath{\varepsilon}}
\newunicodechar{Θ}{\ensuremath{\Theta}}
\newunicodechar{Λ}{\ensuremath{\Lambda}}
\newunicodechar{Μ}{\ensuremath{\mu}}
\newunicodechar{Τ}{\ensuremath{\tau}}
\newunicodechar{Φ}{\ensuremath{\Phi}}
\newunicodechar{Ψ}{\ensuremath{\Psi}}
\newunicodechar{Ω}{\ensuremath{\Omega}}
\newunicodechar{↦}{\ensuremath{\mapsto}}
\newunicodechar{∈}{\ensuremath{\in}}
\newunicodechar{∉}{\ensuremath{\notin}}
\newunicodechar{∧}{\ensuremath{\wedge}}
\newunicodechar{⇔}{\ensuremath{\Leftrightarrow}}
\newunicodechar{⟺}{\ensuremath{\Longleftrightarrow}}
\newunicodechar{⇒}{\ensuremath{\Rightarrow}}
\newunicodechar{⟹}{\ensuremath{\Longrightarrow}}
\newunicodechar{∘}{\ensuremath{\circ}}
\newunicodechar{∪}{\ensuremath{\cup}}
\newunicodechar{∩}{\ensuremath{\cap}}
\newunicodechar{≡}{\ensuremath{\equiv}}
\newunicodechar{⊂}{\ensuremath{\subset}}
\newunicodechar{⊥}{\ensuremath{\perp}}
\newunicodechar{∃}{\ensuremath{\exists}}
\newunicodechar{∀}{\ensuremath{\forall}}
\newunicodechar{∛}{\ensuremath{\sqrt[3]{\,}}}
\newunicodechar{∜}{\ensuremath{\sqrt[4]{\,}}}
\newunicodechar{⃗}{\ensuremath{{}^{\to}}}
\newunicodechar{ⁿ}{\ifmmode^{n}\else\textsuperscript{\ensuremath{n}}\fi}
\newunicodechar{ˣ}{\ifmmode^{x}\else\textsuperscript{\ensuremath{x}}\fi}
\newunicodechar{ᵇ}{\ifmmode^{b}\else\textsuperscript{\ensuremath{b}}\fi}
\newunicodechar{ʳ}{\ifmmode^{r}\else\textsuperscript{\ensuremath{r}}\fi}
\newunicodechar{ᵘ}{\ifmmode^{u}\else\textsuperscript{\ensuremath{u}}\fi}
\newunicodechar{ʸ}{\ifmmode^{y}\else\textsuperscript{\ensuremath{y}}\fi}
\newunicodechar{ᵃ}{\ifmmode^{a}\else\textsuperscript{\ensuremath{a}}\fi}
\newunicodechar{ᵉ}{\ifmmode^{e}\else\textsuperscript{\ensuremath{e}}\fi}
\newunicodechar{ᵏ}{\ifmmode^{k}\else\textsuperscript{\ensuremath{k}}\fi}
\newunicodechar{ᵖ}{\ifmmode^{p}\else\textsuperscript{\ensuremath{p}}\fi}
\newunicodechar{ᴺ}{\ifmmode^{N}\else\textsuperscript{\ensuremath{N}}\fi}
\newunicodechar{ᶜ}{\ifmmode^{c}\else\textsuperscript{\ensuremath{c}}\fi}
\newunicodechar{ʰ}{\ifmmode^{h}\else\textsuperscript{\ensuremath{h}}\fi}
\newunicodechar{ᵠ}{\ifmmode^{q}\else\textsuperscript{\ensuremath{q}}\fi}
\newunicodechar{ₙ}{\ifmmode_{n}\else\textsubscript{\ensuremath{n}}\fi}
\newunicodechar{ₐ}{\ifmmode_{a}\else\textsubscript{\ensuremath{a}}\fi}
\newunicodechar{ᵢ}{\ifmmode_{i}\else\textsubscript{\ensuremath{i}}\fi}
\newunicodechar{ᵣ}{\ifmmode_{r}\else\textsubscript{\ensuremath{r}}\fi}
\newunicodechar{ₖ}{\ifmmode_{k}\else\textsubscript{\ensuremath{k}}\fi}
\newunicodechar{ᵦ}{\ifmmode_{\beta}\else\textsubscript{\ensuremath{\beta}}\fi}
\newunicodechar{ₓ}{\ifmmode_{x}\else\textsubscript{\ensuremath{x}}\fi}
\newfontfamily\popdisplay{ArchivoBlack-Regular.ttf}[Path=../../../fonts/]
\newfontfamily\poplogo{SpaceGrotesk-Bold.ttf}[Path=../../../fonts/]
\newfontfamily\popsemi{PlusJakartaSans-SemiBold.ttf}[Path=../../../fonts/]
\newfontfamily\popxbold{PlusJakartaSans-ExtraBold.ttf}[Path=../../../fonts/]
\newfontfamily\popxboldls{PlusJakartaSans-ExtraBold.ttf}[Path=../../../fonts/,LetterSpace=3.0]

\setlist[enumerate]{leftmargin=1.7em,itemsep=5pt,topsep=4pt}
\setlist[itemize]{leftmargin=1.7em,itemsep=4pt,topsep=4pt,label={\color{poporange}\textbullet}}
\setlength{\parindent}{0pt}
\setlength{\parskip}{5pt}
\renewcommand{\arraystretch}{1.25}

% pgfplots : style Pop par défaut
\pgfplotsset{
  every axis/.append style={axis line style={popink,line width=1pt},tick style={popink},
    grid style={popline,line width=0.5pt},label style={font=\small},tick label style={font=\footnotesize}},
  popcurve/.style={poporange,line width=1.8pt,no markers,smooth},
}
% Même style disponible dans un \draw TikZ simple (hors pgfplots)
\tikzset{popcurve/.style={poporange,line width=1.8pt}}

% ============================================================
% Pill / squiggle
% ============================================================
\newcommand{\poppill}[3]{%
  \tikz[baseline=-0.55ex]\node[draw=popink,line width=1.2pt,rounded corners=2.6pt,%
    fill=#2,inner xsep=6.5pt,inner ysep=3pt]{%
    \popxboldls\fontsize{7}{7}\selectfont\color{#3}\MakeUppercase{#1}};%
}
\newcommand{\popsquiggle}[1]{%
  \tikz[baseline=-0.4ex]{
    \draw[#1,line width=2.6pt,line cap=round]
      plot [smooth] coordinates {(0,0.09) (0.34,0.01) (0.68,0.21) (1.02,0.01) (1.36,0.10)};
  }%
}
% Mot-clé surligné (marqueur orange sous le texte)
\newcommand{\cle}[1]{{\popxbold\color{popdark}#1}}

% ============================================================
% En-tête / pied
% ============================================================
\pagestyle{fancy}
\fancyhf{}
\renewcommand{\headrulewidth}{0pt}
\renewcommand{\footrulewidth}{0pt}
\lhead{\raisebox{-0.15ex}{\poplogo\fontsize{13}{13}\selectfont\color{popink}OnBuch\color{poporange}+}}
\rhead{\poppill{\DOCMATIERE\ \textbullet\ \DOCNIVEAU\ \textbullet\ Leçon \DOCLECON}{white}{popink}}
\lfoot{\popxbold\fontsize{7.6}{7.6}\selectfont\color{popmuted}\MakeUppercase{Cours signé OnBuch+}\ \textbullet\ {\popsemi\DOCTITRE}}
\rfoot{\tikz[baseline=-0.6ex]\node[rounded corners=8.5pt,fill=popink,minimum height=17pt,minimum width=17pt,inner xsep=5pt,inner sep=0]{\popxbold\fontsize{8}{8}\selectfont\color{popcream}\thepage\,/\,\pageref*{LastPage}};}

% ============================================================
% Titres
% ============================================================
\newcommand{\chaptitre}[2]{%
  \begin{tcolorbox}[enhanced,colback=poporange,colframe=popink,boxrule=2.6pt,arc=11pt,
      drop shadow=popink,left=15pt,right=15pt,top=12pt,bottom=11pt,before skip=3pt,after skip=18pt]
    \poppill{#2}{popink}{popcream}\par\vspace{9pt}
    {\raggedright\hyphenpenalty=10000\exhyphenpenalty=10000\popdisplay\fontsize{22}{24}\selectfont\color{popcream}\MakeUppercase{#1}\par}
  \end{tcolorbox}
}
% Section numérotée : pastille orange avec le numéro + titre Archivo
\newcounter{popsec}
\newcommand{\coursec}[1]{%
  \stepcounter{popsec}\setcounter{popsub}{0}%
  \par\vspace{16pt}\noindent
  \begin{minipage}{\linewidth}
  \tikz[baseline=-0.7ex]\node[draw=popink,line width=1.6pt,fill=poporange,rounded corners=4pt,minimum size=22pt,inner sep=0]{\popdisplay\fontsize{12}{12}\selectfont\color{popcream}\thepopsec};%
  \hspace{8pt}{\popdisplay\fontsize{15}{17}\selectfont\color{popink}\MakeUppercase{#1}}\par
  \vspace{4pt}\noindent{\color{poporange}\rule{36pt}{3.6pt}}
  \end{minipage}\par\nopagebreak\vspace{8pt}
}
\newcounter{popsub}
\newcommand{\courssub}[1]{%
  \stepcounter{popsub}%
  \par\vspace{9pt}\noindent{\popxbold\fontsize{11.5}{13}\selectfont\color{popink}{\color{poporange}\thepopsec.\thepopsub}\hspace{6pt}#1}\par\nopagebreak\vspace{4pt}
}

% ============================================================
% Boîtes pédagogiques — même recette : fond blanc, bordure épaisse colorée,
% ombre dure encre, badge plein. Titre optionnel [#1].
% ============================================================
\tcbset{popbase/.style={breakable,enhanced,colback=white,boxrule=2.3pt,arc=9pt,
  shadow={3pt}{-3pt}{0pt}{popink},left=14pt,right=14pt,top=11pt,bottom=11pt,before skip=11pt,after skip=15pt,
  fontupper=\popsemi\fontsize{10.6}{14.5}\selectfont\color{popink},
  fontlower=\popsemi\fontsize{10.6}{14.5}\selectfont\color{popink}}}
% Coche dessinée (le glyphe ✓ n'existe pas dans les polices du document)
\newcommand{\popcheck}[1]{\tikz[baseline=-0.5ex]\draw[#1,line width=1.6pt,line cap=round,line join=round](-0.13,0.0)--(-0.04,-0.09)--(0.13,0.1);}
\newcommand{\popboxhead}[3]{\poppill{#1}{#2}{white}\ifx\relax#3\relax\else\hspace{6pt}{\popxbold\fontsize{10.6}{12}\selectfont\color{popink}#3}\fi\par\vspace{7pt}}

\newtcolorbox{aretenir}[1][]{popbase,colframe=popgreen,before upper={\popboxhead{À retenir}{popgreen}{#1}}}
\newtcolorbox{exemplebox}[1][]{popbase,colframe=poporange,before upper={\popboxhead{Exemple}{poporange}{#1}}}
\newtcolorbox{definition}[1][]{popbase,colframe=popblue,before upper={\popboxhead{Définition}{popblue}{#1}}}
\newtcolorbox{propriete}[1][]{popbase,colframe=poppurple,before upper={\popboxhead{Propriété}{poppurple}{#1}}}
\newtcolorbox{methode}[1][]{popbase,colframe=popgold,before upper={\popboxhead{Méthode}{popgold}{#1}}}
\newtcolorbox{attention}[1][]{popbase,colframe=poppink,before upper={\popboxhead{Attention}{poppink}{#1}}}
\newtcolorbox{experience}[1][]{popbase,colframe=popblue,colback=popblueL,before upper={\popboxhead{Expérience}{popblue}{#1}}}
% « Le savais-tu ? » : carte violette pleine (contexte, culture, Cameroun)
\newtcolorbox{savaistu}[1][]{popbase,colback=poppurple,colframe=popink,
  fontupper=\popsemi\fontsize{10.4}{14.2}\selectfont\color{white},
  before upper={\poppill{Le savais-tu\,?}{white}{poppurple}\ifx\relax#1\relax\else\hspace{6pt}{\popxbold\color{white}#1}\fi\par\vspace{7pt}}}
% Exercice résolu : énoncé en haut, \tcblower puis la solution sur fond crème
\newcounter{popexo}\newcounter{popexr}
\newtcolorbox{exoresolu}[1][]{popbase,colframe=poporange,colbacklower=popcream,
  segmentation style={popink,line width=1.2pt,dashed},
  before upper={\stepcounter{popexr}\popboxhead{Exercice résolu \thepopexr}{poporange}{#1}},
  before lower={\poppill{Solution}{popink}{popcream}\par\vspace{6pt}}}
% Exercice (non résolu en place) — la correction est dans la partie Corrigés
\newtcolorbox{exercice}[1][]{popbase,colframe=popink,
  before upper={\stepcounter{popexo}\popboxhead{Exercice \thepopexo}{popink}{#1}}}
% Corrigé
\newtcolorbox{corrige}[1][]{popbase,colframe=popgreen,colback=popgreenL,
  before upper={\popboxhead{Corrigé}{popgreen}{#1}}}
% Objectifs / prérequis / bilan
\newtcolorbox{objectifs}{popbase,colframe=popink,colback=poporangeL,
  before upper={\popboxhead{Objectifs de la leçon}{popink}{}}}
\newtcolorbox{prerequis}{popbase,colframe=popink,colback=popblueL,
  before upper={\popboxhead{Prérequis}{popblue}{}}}
\newtcolorbox{bilan}{popbase,colframe=popink,colback=white,boxrule=2.8pt,
  before upper={\popboxhead{Fiche bilan}{poporange}{L'essentiel à connaître pour l'examen}}}
% Figure encadrée avec légende
\newtcolorbox{popfigure}[1][]{enhanced,breakable=false,colback=white,colframe=popline,boxrule=1.4pt,arc=8pt,
  left=10pt,right=10pt,top=10pt,bottom=8pt,before skip=10pt,after skip=12pt,halign=center,
  lower separated=false,halign lower=center,
  fontlower=\popsemi\fontsize{9}{11.5}\selectfont\color{popmuted},
  #1}
\newcommand{\legende}[1]{\par\vspace{6pt}{\popsemi\fontsize{9}{11.5}\selectfont\color{popmuted}#1}}

% ============================================================
% Signature OnBuch+
% ============================================================
\newcommand{\poplogoinline}[1]{{\poplogo\fontsize{#1}{#1}\selectfont\color{popink}OnBuch\color{poporange}+}}
\newcommand{\signatureonbuch}{%
  \begin{tcolorbox}[enhanced,colback=popink,colframe=popink,boxrule=2.2pt,arc=10pt,
      drop shadow=poporange,left=14pt,right=14pt,top=10pt,bottom=10pt,before skip=0pt,after skip=0pt]
    \tikz[baseline=-0.7ex]\node[circle,fill=popgreen,draw=popcream,line width=1.3pt,minimum size=22pt,inner sep=0]{\popcheck{white}};%
    \hspace{9pt}\begin{minipage}[c]{0.62\linewidth}
      {\popxbold\fontsize{12}{13}\selectfont\color{popcream}Signé\ \textbullet\ {\poplogo OnBuch\color{poporange}+}}\par\vspace{2pt}
      {\raggedright\popsemi\fontsize{8}{10.5}\selectfont\color{popline}\MakeUppercase{Équipe pédagogique OnBuch+}\\\MakeUppercase{Conforme au programme officiel MINESEC}\par}
    \end{minipage}\hfill
    {\poplogo\fontsize{22}{22}\selectfont\color{popcream}OnBuch\color{poporange}+}
  \end{tcolorbox}%
}

% ============================================================
% Page de garde
% #1 titre · #2 matière · #3 niveau · #4 module · #5 n° leçon · #6 durée ·
% #7 description / objectifs (texte ou itemize)
% ============================================================
\newcommand{\pagegarde}[7]{%
  \thispagestyle{empty}
  \newgeometry{a4paper,margin=1.7cm,top=1.6cm,bottom=1.4cm}%
  \begin{tikzpicture}[remember picture,overlay]
    \fill[poporange!18] ([xshift=-3.2cm,yshift=-3.6cm]current page.north east) circle (4.6cm);
    \fill[poppurple!14] ([xshift=2.4cm,yshift=7.6cm]current page.south west) circle (3.8cm);
  \end{tikzpicture}%
  {\poplogo\fontsize{30}{30}\selectfont\color{popink}OnBuch\color{poporange}+}\hfill
  \raisebox{0.6ex}{\poppill{Cours signé \textbullet\ Premium}{popink}{popcream}}\par
  \vspace{1pt}\popsquiggle{poppurple}\par
  \vspace{8pt}
  \begin{tcolorbox}[enhanced,colback=poporange,colframe=popink,boxrule=2.8pt,arc=13pt,
      drop shadow=popink,left=18pt,right=18pt,top=12pt,bottom=13pt,after skip=10pt]
    \poppill{#2 \textbullet\ #3}{popink}{popcream}\hspace{5pt}\poppill{#4}{white}{popink}\par\vspace{8pt}
    {\popdisplay\fontsize{14}{14}\selectfont\color{popink}LEÇON #5}\par\vspace{4pt}
    {\raggedright\hyphenpenalty=10000\exhyphenpenalty=10000\popdisplay\fontsize{25}{27}\selectfont\color{popcream}\MakeUppercase{#1}\par}
    \vspace{8pt}
    {\popxbold\fontsize{9.5}{9.5}\selectfont\color{popink}\MakeUppercase{Durée conseillée : #6}}
  \end{tcolorbox}
  \begin{tcolorbox}[enhanced,colback=white,colframe=popink,boxrule=2.2pt,arc=10pt,
      drop shadow=popink,left=16pt,right=16pt,top=10pt,bottom=10pt,after skip=8pt]
    \poppill{Dans cette leçon}{poppurple}{white}\par\vspace{5pt}
    \popsemi\fontsize{9.3}{12.6}\selectfont\color{popink}#7
  \end{tcolorbox}
  \noindent\begin{minipage}[t]{0.487\linewidth}
    \begin{tcolorbox}[enhanced,colback=popgreen,colframe=popink,boxrule=2.2pt,arc=10pt,
        drop shadow=popink,left=11pt,right=11pt,top=10pt,bottom=10pt,height=3.1cm,valign=center]
      \tcbox[colback=white,colframe=popink,boxrule=1.3pt,arc=5pt,boxsep=0pt,nobeforeafter,
        left=3pt,right=3pt,top=3pt,bottom=3pt,valign=center]{\qrcode[height=1.9cm]{https://play.google.com/store/apps/details?id=com.luvvix.onbuch&hl=fr}}%
      \hspace{8pt}\begin{minipage}[c]{0.5\linewidth}
        {\popdisplay\fontsize{11}{12.5}\selectfont\color{white}\MakeUppercase{Télécharge OnBuch}}\par\vspace{4pt}
        {\raggedright\popsemi\fontsize{8.4}{11}\selectfont\color{white}Gratuit \textbullet\ Google Play\\Tuteur IA, annales, concours, résultats\par}
      \end{minipage}
    \end{tcolorbox}
  \end{minipage}\hfill
  \begin{minipage}[t]{0.487\linewidth}
    \begin{tcolorbox}[enhanced,colback=poppurple,colframe=popink,boxrule=2.2pt,arc=10pt,
        drop shadow=popink,left=11pt,right=11pt,top=10pt,bottom=10pt,height=3.1cm,valign=center]
      \tikz[baseline=-0.7ex]\node[circle,fill=white,draw=popink,line width=1.3pt,minimum size=1.6cm,inner sep=0]{\popdisplay\fontsize{24}{24}\selectfont\color{poppurple}+};%
      \hspace{8pt}\begin{minipage}[c]{0.6\linewidth}
        {\popdisplay\fontsize{11}{12.5}\selectfont\color{white}\MakeUppercase{Passe à OnBuch+}}\par\vspace{4pt}
        {\raggedright\popsemi\fontsize{8.4}{11}\selectfont\color{white}Tous les cours signés, Tuteur IA illimité, fiches PDF — onglet OnBuch+ de l'app\par}
      \end{minipage}
    \end{tcolorbox}
  \end{minipage}
  \vspace{8pt}
  \popcontenu
  \vfill
  \signatureonbuch
  \restoregeometry
  \newpage
}

% Rangée « Ce cours contient » (page de garde)
\newcommand{\poptile}[3]{% #1 couleur #2 gros texte #3 légende
  \begin{tcolorbox}[enhanced,colback=#1,colframe=popink,boxrule=2pt,arc=9pt,shadow={2.5pt}{-2.5pt}{0pt}{popink},
      left=6pt,right=6pt,top=9pt,bottom=9pt,halign=center,valign=center,height=1.75cm,width=\linewidth,nobeforeafter]
    {\popdisplay\fontsize{12.5}{13}\selectfont\color{white}\MakeUppercase{#2}}\par\vspace{3pt}
    {\centering\popsemi\fontsize{7.8}{9.5}\selectfont\color{white}#3\par}
  \end{tcolorbox}}
\newcommand{\popcontenu}{%
  \par\noindent{\popxboldls\fontsize{7.5}{7.5}\selectfont\color{popmuted}CE COURS SIGNÉ CONTIENT}\par\vspace{7pt}
  \noindent\begin{minipage}[t]{0.235\linewidth}\poptile{popblue}{Cours}{complet, illustré, pas à pas}\end{minipage}\hfill
  \begin{minipage}[t]{0.235\linewidth}\poptile{poporange}{Méthodes}{et exercices résolus}\end{minipage}\hfill
  \begin{minipage}[t]{0.235\linewidth}\poptile{poppink}{Exercices}{type examen + corrigés}\end{minipage}\hfill
  \begin{minipage}[t]{0.235\linewidth}\poptile{popgreen}{Fiche bilan}{l'essentiel à retenir}\end{minipage}\par}

% Sommaire
\newtcolorbox{sommaire}{enhanced,colback=white,colframe=popink,boxrule=2.2pt,arc=10pt,
  drop shadow=popink,left=18pt,right=18pt,top=13pt,bottom=13pt,before skip=6pt,after skip=20pt,
  before upper={\poppill{Sommaire}{poppurple}{white}\par\vspace{9pt}\popsemi\fontsize{10.6}{14.5}\selectfont\color{popink}}}

% Signature de fin de cours — poussée tout en bas de la dernière page
\newcommand{\findecours}{%
  \par\vfill\nopagebreak
  \begin{center}\popsquiggle{poporange}\end{center}\vspace{4pt}
  \signatureonbuch
}
\AtBeginDocument{\def\llbracket{\mathopen{[\mkern-3.5mu[}}\def\rrbracket{\mathclose{]\mkern-3.5mu]}}} % intervalles d'entiers (glyphe ⟦ absent de la police)
% Glyphes absents de Fira Math : redéfinis à partir de glyphes disponibles
\AtBeginDocument{%
  \def\setminus{\mathbin{\backslash}}\let\smallsetminus\setminus
  \def\vdots{\mathord{\vbox{\baselineskip3.2pt\lineskiplimit0pt\kern4pt\hbox{.}\hbox{.}\hbox{.}}}}%
  \def\ddots{\mathinner{\mkern1mu\raise6.5pt\hbox{.}\mkern2mu\raise3.6pt\hbox{.}\mkern2mu\raise0.7pt\hbox{.}\mkern1mu}}%
  \def\triangle{\mathord{\scalebox{0.9}{$\symup{\Delta}$}}}%
  \def\nabla{\mathord{\raisebox{\depth}{\scalebox{1}[-1]{$\symup{\Delta}$}}}}%
  \def\checkmark{\popcheck{popdark}}%
  \def\star{\mathbin{\ast}}\def\bigstar{\mathord{\ast}}%
  \def\diamond{\mathbin{\scalebox{0.6}{$\Diamond$}}}%
  \def\bigcup{\mathop{\mathchoice{\raisebox{-0.3em}{\scalebox{1.7}{$\cup$}}}{\scalebox{1.25}{$\cup$}}{\cup}{\cup}}\displaylimits}%
  \def\bigcap{\mathop{\mathchoice{\raisebox{-0.3em}{\scalebox{1.7}{$\cap$}}}{\scalebox{1.25}{$\cap$}}{\cap}{\cap}}\displaylimits}%
  \def\lesseqqgtr{\mathrel{\lessgtr}}\def\bigoplus{\mathop{\oplus}\displaylimits}%
}
\providecommand{\ssi}{si et seulement si} % abréviation fréquente
\AtBeginDocument{\providecommand{\wideparen}{\overparen}} % arc de cercle

\begin{document}

\pagegarde{Angles inscrits}{Mathématiques}{3ème}{Mathématiques – classe de 3ème}{N12}{2 séances (fiche de progression harmonisée)}{Cette leçon introduit les angles inscrits dans un cercle et les angles au centre associés, puis établit les deux théorèmes fondamentaux qui les relient. Elle fournit des outils puissants pour calculer des angles, démontrer des alignements, des cocyclicités et résoudre des problèmes concrets de la vie courante.
\vspace{4pt}
\begin{multicols}{2}\small
\begin{itemize}
\item À la fin de cette leçon, je sais reconnaître un angle inscrit et un angle au centre dans un cercle
\item À la fin de cette leçon, je sais identifier l'arc intercepté par un angle inscrit ou un angle au centre
\item À la fin de cette leçon, je sais énoncer et appliquer le théorème de l'angle inscrit et de l'angle au centre associés
\item À la fin de cette leçon, je sais démontrer que deux angles inscrits interceptant le même arc sont égaux
\item À la fin de cette leçon, je sais calculer des mesures d'angles dans des figures comportant un cercle
\item À la fin de cette leçon, je sais utiliser le fait qu'un triangle inscrit dans un cercle dont un côté est un diamètre est rectangle
\end{itemize}\end{multicols}}

\begin{sommaire}
\begin{multicols}{2}
\begin{enumerate}
\item Angles inscrits et angles au centre : définitions
\item Théorème de l'angle inscrit et de l'angle au centre
\item Angles inscrits interceptant le même arc
\item Triangle inscrit dans un demi-cercle
\item Calculs d'angles et démonstrations
\item Applications à la vie courante
\item Activité d'intégration
\item Méthodes et exercices résolus
\item Exercices
\item Corrigés des exercices
\item Fiche bilan
\end{enumerate}
\end{multicols}
\end{sommaire}

\begin{objectifs}
\begin{itemize}
\item À la fin de cette leçon, je sais reconnaître un angle inscrit et un angle au centre dans un cercle
\item À la fin de cette leçon, je sais identifier l'arc intercepté par un angle inscrit ou un angle au centre
\item À la fin de cette leçon, je sais énoncer et appliquer le théorème de l'angle inscrit et de l'angle au centre associés
\item À la fin de cette leçon, je sais démontrer que deux angles inscrits interceptant le même arc sont égaux
\item À la fin de cette leçon, je sais calculer des mesures d'angles dans des figures comportant un cercle
\item À la fin de cette leçon, je sais utiliser le fait qu'un triangle inscrit dans un cercle dont un côté est un diamètre est rectangle
\item À la fin de cette leçon, je sais rédiger et justifier une démarche complète de démonstration géométrique
\item À la fin de cette leçon, je sais appliquer ces notions à des situations concrètes de la vie courante (angles de vision, tir au but, orientation)
\end{itemize}
\end{objectifs}

\begin{prerequis}
\begin{itemize}
\item Notion de cercle : centre, rayon, diamètre, corde, arc de cercle (6ème-5ème)
\item Mesure des angles, angles au centre (4ème)
\item Propriétés des triangles isocèles et somme des angles d'un triangle (5ème)
\item Angles et parallélisme : angles alternes-internes, correspondants (4ème)
\item Notion de triangle rectangle et réciproque du théorème de Pythagore (4ème)
\end{itemize}
\end{prerequis}

\newpage

\coursec{Angles inscrits et angles au centre : définitions}

Au stade Ahmadou Ahidjo de Yaoundé, un soir de match, deux supporters sont assis à des places différentes des tribunes. Tous les deux regardent le même poteau de corner et le drapeau de l'arbitre assistant. Ont-ils le même angle de vision ? Un technicien chargé d'installer des projecteurs autour du terrain, supposé circulaire, se pose exactement la même question : l'angle sous lequel on voit une portion du terrain dépend-il de la position de l'observateur sur le cercle ? Et que dire de l'angle de vision du speaker, placé lui au centre du stade ?

Pour répondre rigoureusement à ces questions, il faut d'abord apprendre à nommer et à reconnaître deux types d'angles liés à un cercle : les \cle{angles inscrits} et les \cle{angles au centre}. C'est l'objet de cette première section. La problématique est la suivante : \emph{comment définir, reconnaître et associer correctement un angle inscrit et un angle au centre dans un cercle ?}

\courssub{Rappels : cercle, corde et arc de cercle}

Avant de parler d'angles, remettons en place le décor : le cercle et ses parties.

\begin{definition}[Cercle]
Un \cle{cercle} de centre $O$ et de rayon $r$ est l'ensemble des points du plan situés à la distance $r$ du point $O$. On le note souvent $(\mathcal{C})$ ou $\mathcal{C}(O\,;\,r)$.
\end{definition}

\begin{definition}[Corde et arc de cercle]
Soit $(\mathcal{C})$ un cercle de centre $O$, et $A$ et $B$ deux points de ce cercle.
\begin{itemize}
\item Le segment $[AB]$ s'appelle une \cle{corde} du cercle. Si la corde passe par le centre $O$, c'est un \cle{diamètre}.
\item La portion de cercle comprise entre $A$ et $B$ s'appelle un \cle{arc de cercle}, noté $\overset{\frown}{AB}$.
\end{itemize}
\end{definition}

\begin{attention}[Deux arcs pour deux points]
Deux points $A$ et $B$ d'un cercle délimitent en réalité \textbf{deux} arcs : le \cle{petit arc} $\overset{\frown}{AB}$ (le plus court, appelé aussi arc mineur) et le \cle{grand arc} $\overset{\frown}{AB}$ (le plus long, appelé aussi arc majeur). Quand il y a un doute, on ajoute un point intermédiaire pour préciser : par exemple $\overset{\frown}{AMB}$ désigne l'arc d'extrémités $A$ et $B$ qui passe par $M$.
\end{attention}

\textbf{Mesure d'un arc.} Un arc de cercle se mesure en degrés : le cercle entier mesure $360°$, un demi-cercle $180°$, un quart de cercle $90°$. La mesure d'un arc est par définition la mesure de l'angle au centre qui le « voit » : si $\widehat{AOB} = 80°$, alors l'arc $\overset{\frown}{AB}$ mesure $80°$.

\courssub{Angle inscrit dans un cercle}

\textbf{Intuition.} Revenons au stade. Un supporter est assis quelque part sur la tribune circulaire : ses yeux sont \emph{sur} le cercle. Il regarde deux objets situés eux aussi sur le cercle (le poteau de corner $A$ et le drapeau $B$). Ses deux rayons visuels partent de sa position $M$ et recoupent le cercle en $A$ et $B$. L'angle $\widehat{AMB}$ ainsi formé est exactement ce que les mathématiciens appellent un angle inscrit : « inscrit » parce qu'il est écrit \emph{à l'intérieur} du cercle, le sommet posé sur le cercle lui-même.

\begin{definition}[Angle inscrit]
Un \cle{angle inscrit} dans un cercle est un angle dont le \textbf{sommet appartient au cercle} et dont les \textbf{deux côtés recoupent le cercle} en deux autres points.

On dit que l'angle $\widehat{AMB}$ est inscrit dans le cercle $(\mathcal{C})$ lorsque les trois points $A$, $M$ et $B$ appartiennent à $(\mathcal{C})$.
\end{definition}

\begin{exemplebox}[Reconnaître un angle inscrit]
Sur un cercle $(\mathcal{C})$, on place quatre points $A$, $B$, $M$, $N$.
\begin{itemize}
\item $\widehat{AMB}$ est un angle inscrit : son sommet $M$ est sur le cercle, et ses côtés $[MA)$ et $[MB)$ recoupent le cercle en $A$ et $B$.
\item $\widehat{ANB}$ est aussi un angle inscrit, de sommet $N$.
\item En revanche, si $O$ est le centre du cercle, $\widehat{AOB}$ n'est \textbf{pas} un angle inscrit, car $O$ n'est pas sur le cercle (sauf cas dégénéré sans intérêt).
\end{itemize}
\end{exemplebox}

\courssub{Angle au centre}

\textbf{Intuition.} Le speaker du stade, lui, est installé au centre exact du terrain circulaire. Lui aussi regarde le poteau $A$ et le drapeau $B$. Son angle de vision $\widehat{AOB}$ a son sommet au \emph{centre} du cercle : c'est un angle au centre.

\begin{definition}[Angle au centre]
Un \cle{angle au centre} est un angle dont le \textbf{sommet est le centre} du cercle.

Si $O$ est le centre de $(\mathcal{C})$ et si $A$ et $B$ sont deux points de $(\mathcal{C})$, alors $\widehat{AOB}$ est un angle au centre.
\end{definition}

\begin{aretenir}[Le critère qui ne trompe pas]
Pour reconnaître la nature d'un angle lié à un cercle, regarde \textbf{uniquement la position du sommet} :
\begin{itemize}
\item sommet \emph{sur} le cercle $\Rightarrow$ angle \textbf{inscrit} ;
\item sommet \emph{au centre} du cercle $\Rightarrow$ angle \textbf{au centre}.
\end{itemize}
\end{aretenir}

\courssub{Arc intercepté par un angle}

\textbf{Intuition.} Quand le supporter en $M$ regarde le poteau $A$ et le drapeau $B$, son regard « encadre » une portion du bord du terrain : l'arc de cercle situé \emph{en face de lui}, entre $A$ et $B$. Cette portion, c'est l'arc intercepté. Point essentiel : ce n'est \textbf{pas} l'arc qui passe derrière le supporter, c'est celui qui est \emph{devant} lui, à l'intérieur de son angle de vision.

\begin{definition}[Arc intercepté]
Soit un angle dont les côtés coupent un cercle en deux points $A$ et $B$. L'\cle{arc intercepté} par cet angle est l'arc $\overset{\frown}{AB}$ situé \textbf{à l'intérieur de l'angle}, c'est-à-dire l'arc qui \textbf{ne contient pas le sommet} de l'angle.
\end{definition}

Ainsi, pour l'angle inscrit $\widehat{AMB}$, l'arc intercepté est l'arc $\overset{\frown}{AB}$ qui ne passe pas par $M$. L'autre arc, celui qui contient $M$, n'est \emph{jamais} l'arc intercepté : c'est l'erreur la plus fréquente chez les élèves.

\begin{attention}[Piège : l'arc intercepté n'est pas l'arc du sommet]
Face à l'angle inscrit $\widehat{AMB}$, beaucoup d'élèves désignent par erreur l'arc qui contient $M$ comme « l'arc intercepté ». C'est faux ! Retiens l'image suivante : l'angle est comme une bouche ouverte en $M$ ; l'arc intercepté est le morceau que la bouche va croquer, c'est-à-dire l'arc \textbf{en face}, celui qui \textbf{ne contient pas} $M$. Si l'arc intercepté est petit (arc mineur), l'arc contenant $M$ est grand (arc majeur), et réciproquement.
\end{attention}

\courssub{Angles inscrit et angle au centre associés}

\textbf{Intuition.} Le supporter en $M$ et le speaker en $O$ regardent tous les deux \emph{la même portion} du bord du terrain : le même arc $\overset{\frown}{AB}$. On dit alors que leurs angles de vision sont \emph{associés}. Cette notion d'association est la clé de toute la leçon, car c'est elle qui permettra, dans la section suivante, de comparer leurs mesures.

\begin{definition}[Angles associés]
Un angle inscrit et un angle au centre sont dits \cle{associés} lorsqu'ils \textbf{interceptent le même arc}.

Autrement dit : l'angle inscrit $\widehat{AMB}$ et l'angle au centre $\widehat{AOB}$ sont associés s'ils « regardent » tous les deux le même arc $\overset{\frown}{AB}$ (celui qui ne contient pas $M$).
\end{definition}

\begin{exemplebox}[Exemple chiffré]
Sur un cercle de centre $O$, l'arc $\overset{\frown}{AB}$ mesure $80°$. Alors, par définition de la mesure d'un arc, l'angle au centre associé mesure :
\[ \widehat{AOB} = 80°. \]
Tout angle inscrit $\widehat{AMB}$ associé à ce même arc (c'est-à-dire dont le sommet $M$ est sur le grand arc) intercepte lui aussi l'arc de $80°$. Nous verrons à la section suivante que sa mesure vaut alors la moitié : $\widehat{AMB} = 40°$.
\end{exemplebox}

\begin{popfigure}
\begin{center}
\begin{tikzpicture}[scale=1.6]
\coordinate (O) at (0,0);
\draw[popline, thick] (O) circle (2);
\coordinate (A) at (210:2);
\coordinate (B) at (-30:2);
\coordinate (M) at (100:2);
% arc intercepté colorié (arc mineur AB ne contenant pas M)
\draw[poporange, very thick] (210:2) arc (210:330:2);
% angle au centre
\draw[popblue, thick] (O) -- (A);
\draw[popblue, thick] (O) -- (B);
% angle inscrit
\draw[popgreen, thick] (M) -- (A);
\draw[popgreen, thick] (M) -- (B);
% points
\fill[popdark] (O) circle (0.045);
\fill[popdark] (A) circle (0.045);
\fill[popdark] (B) circle (0.045);
\fill[popdark] (M) circle (0.045);
% labels
\node[below left, popdark] at (O) {$O$};
\node[below left, popdark] at (A) {$A$};
\node[below right, popdark] at (B) {$B$};
\node[above, popdark] at (M) {$M$};
% marques d'angles
\draw[popblue] (210:0.55) arc (210:330:0.55);
\node[popblue] at (270:0.85) {$\widehat{AOB}$};
\draw[popgreen] ($(M)+(240:0.5)$) arc (240:300:0.5);
\node[popgreen] at ($(M)+(-90:0.85)$) {$\widehat{AMB}$};
\node[poporange] at (270:2.45) {arc intercepté $\overset{\frown}{AB}$};
\end{tikzpicture}
\end{center}
\legende{L'angle inscrit $\widehat{AMB}$ (vert) et l'angle au centre $\widehat{AOB}$ (bleu) interceptent le même arc $\overset{\frown}{AB}$ (orange) : ce sont des angles associés.}
\end{popfigure}

\courssub{Lecture de figures : identifier angles inscrits, angles au centre et arcs}

Dans les exercices du BEPC, les figures contiennent souvent plusieurs points sur un même cercle. Il faut savoir y lire, sans se tromper, les angles inscrits, les angles au centre et les arcs interceptés. Voici la démarche du professeur.

\begin{methode}[Identifier les angles et leurs arcs sur une figure]
\begin{enumerate}
\item Repérer le \textbf{centre} $O$ du cercle (souvent nommé ou marqué).
\item Pour chaque angle de la figure, regarder la \textbf{position du sommet} : sur le cercle (angle inscrit) ou au centre (angle au centre).
\item Pour chaque angle, prolonger mentalement ses deux côtés jusqu'au cercle : ils déterminent deux points, par exemple $A$ et $B$.
\item L'\textbf{arc intercepté} est l'arc $\overset{\frown}{AB}$ situé \emph{en face} de l'angle, c'est-à-dire celui qui \textbf{ne contient pas} le sommet.
\item Deux angles sont \textbf{associés} (ou interceptent le même arc) si leurs arcs interceptés coïncident.
\end{enumerate}
\end{methode}

\begin{popfigure}
\begin{center}
\begin{tikzpicture}[scale=1.6]
\coordinate (O) at (0,0);
\draw[popline, thick] (O) circle (2);
\coordinate (A) at (210:2);
\coordinate (B) at (-30:2);
\coordinate (M) at (80:2);
\coordinate (N) at (140:2);
% arc intercepté colorié
\draw[poporange, very thick] (210:2) arc (210:330:2);
% angles inscrits
\draw[popgreen, thick] (M) -- (A);
\draw[popgreen, thick] (M) -- (B);
\draw[poppurple, thick] (N) -- (A);
\draw[poppurple, thick] (N) -- (B);
% points
\fill[popdark] (O) circle (0.045);
\fill[popdark] (A) circle (0.045);
\fill[popdark] (B) circle (0.045);
\fill[popdark] (M) circle (0.045);
\fill[popdark] (N) circle (0.045);
% labels
\node[below, popdark] at (O) {$O$};
\node[below left, popdark] at (A) {$A$};
\node[below right, popdark] at (B) {$B$};
\node[above right, popdark] at (M) {$M$};
\node[above left, popdark] at (N) {$N$};
\node[poporange] at (270:2.45) {arc $\overset{\frown}{AB}$};
\end{tikzpicture}
\end{center}
\legende{Les angles inscrits $\widehat{AMB}$ (vert) et $\widehat{ANB}$ (violet) ont des sommets différents mais interceptent le même arc $\overset{\frown}{AB}$ (orange).}
\end{popfigure}

\begin{exoresolu}[Lecture d'une figure]
Sur un cercle $(\mathcal{C})$ de centre $O$, on place quatre points $A$, $B$, $M$ et $N$ dans cet ordre, $M$ et $N$ étant sur le même arc délimité par $A$ et $B$.
\begin{enumerate}
\item Citer deux angles inscrits interceptant l'arc $\overset{\frown}{AB}$ qui ne contient ni $M$ ni $N$.
\item Citer l'angle au centre associé à ces angles inscrits.
\item L'angle $\widehat{MON}$ est-il un angle inscrit ? Justifier.
\end{enumerate}
\tcblower
\textbf{Solution détaillée.}
\begin{enumerate}
\item Les angles $\widehat{AMB}$ et $\widehat{ANB}$ ont leurs sommets $M$ et $N$ sur le cercle, et leurs côtés recoupent le cercle en $A$ et $B$ : ce sont des angles inscrits. Comme $M$ et $N$ sont sur le même arc, l'arc situé à l'intérieur de chacun de ces angles est l'arc $\overset{\frown}{AB}$ qui ne contient ni $M$ ni $N$. Ils interceptent donc le même arc.
\item L'angle au centre associé est $\widehat{AOB}$ : son sommet est le centre $O$ et il intercepte le même arc $\overset{\frown}{AB}$. Les angles $\widehat{AMB}$ et $\widehat{AOB}$ sont donc associés, de même que $\widehat{ANB}$ et $\widehat{AOB}$.
\item Non. L'angle $\widehat{MON}$ a bien ses côtés qui coupent le cercle en $M$ et $N$, mais son sommet $O$ est le \emph{centre} du cercle, pas un point du cercle. C'est donc un \textbf{angle au centre}, et non un angle inscrit.
\end{enumerate}
\end{exoresolu}

\begin{attention}[Erreurs fréquentes à éviter]
\begin{itemize}
\item \textbf{Confondre angle inscrit et angle au centre} : vérifie toujours la position du sommet avant de conclure. Un angle dont le sommet est $O$ est un angle au centre, même si ses côtés ressemblent à ceux d'un angle inscrit.
\item \textbf{Se tromper d'arc} : l'arc intercepté par $\widehat{AMB}$ est celui qui ne contient \emph{pas} $M$. Sur une figure où les points sont proches, trace mentalement l'intérieur de l'angle pour être sûr.
\item \textbf{Associer des angles qui ne regardent pas le même arc} : deux angles ne sont associés que si leurs arcs interceptés sont exactement les mêmes (mêmes extrémités, même côté).
\end{itemize}
\end{attention}

\begin{savaistu}[Pourquoi cette leçon est-elle si utile ?]
Les angles inscrits sont partout autour de nous. Les techniciens qui éclairent le stade Ahmadou Ahidjo règlent la position des projecteurs grâce aux angles sous lesquels les portions du terrain sont vues. Les footballeurs, eux, savent intuitivement que l'angle de tir face aux poteaux dépend de leur position : c'est exactement un problème d'angle inscrit ! Historiquement, c'est le mathématicien grec \textbf{Euclide} qui, vers 300 av. J.-C., a énoncé dans ses \emph{Éléments} les propriétés des angles inscrits, et le théorème que nous démontrerons à la section suivante est parfois attribué à \textbf{Thalès de Milet}. Ces résultats, vieux de plus de deux mille ans, servent encore aujourd'hui en navigation, en astronomie et en génie civil.
\end{savaistu}

\begin{aretenir}[Bilan de la section]
\begin{itemize}
\item Un \cle{angle inscrit} a son sommet \textbf{sur} le cercle et ses côtés recoupent le cercle.
\item Un \cle{angle au centre} a son sommet \textbf{au centre} du cercle.
\item L'\cle{arc intercepté} par un angle est l'arc situé \textbf{à l'intérieur} de l'angle, qui \textbf{ne contient pas} son sommet.
\item Un angle inscrit et un angle au centre sont \cle{associés} lorsqu'ils interceptent \textbf{le même arc}.
\item La mesure de l'angle au centre est égale à la mesure de son arc intercepté : arc de $80°$ $\Rightarrow$ $\widehat{AOB} = 80°$.
\end{itemize}
\end{aretenir}

Nous savons maintenant définir et reconnaître ces angles. Mais revenons à nos deux supporters du stade : leurs angles de vision $\widehat{AMB}$ et $\widehat{ANB}$ sont-ils égaux ? Et comment se comparent-ils à l'angle $\widehat{AOB}$ du speaker ? C'est le \textbf{théorème de l'angle inscrit et de l'angle au centre}, objet de la section suivante, qui nous donnera la réponse.

\coursec{Théorème de l'angle inscrit et de l'angle au centre}

Dans la section précédente, nous avons appris à reconnaître un \cle{angle inscrit} (sommet sur le cercle, côtés qui coupent le cercle) et un \cle{angle au centre} (sommet au centre du cercle), et à identifier l'\cle{arc intercepté} par chacun d'eux. Nous avons remarqué qu'un même arc $\wideparen{AB}$ peut être « vu » de deux façons : depuis le centre $O$ (angle au centre $\widehat{AOB}$) ou depuis un point $M$ du cercle (angle inscrit $\widehat{AMB}$). Une question naturelle se pose alors : \emph{existe-t-il une relation entre ces deux mesures ?} C'est l'objet de cette section, et la réponse est l'un des plus beaux théorèmes de la géométrie du cercle.

\courssub{Conjecture expérimentale : mesurer pour deviner}

\begin{experience}[À toi de mesurer !]
Réalise l'activité suivante avec ton rapporteur et ton compas.
\begin{enumerate}
\item Trace un cercle de centre $O$ et de rayon $4$ cm. Place deux points $A$ et $B$ sur le cercle tels que $\widehat{AOB} = 100°$.
\item Place un point $M$ sur le \emph{grand} arc $\wideparen{AB}$ (celui qui ne contient pas l'angle au centre mesuré). Mesure $\widehat{AMB}$ au rapporteur.
\item Déplace $M$ en un autre point $M'$ du même arc, puis en un point $M''$. Mesure à nouveau $\widehat{AM'B}$ et $\widehat{AM''B}$.
\item Recommence avec $\widehat{AOB} = 80°$, puis avec $\widehat{AOB} = 130°$.
\end{enumerate}
\textbf{Constat :} quelle que soit la position de $M$ sur le même arc, l'angle inscrit garde la même mesure, et cette mesure vaut \emph{la moitié} de l'angle au centre : $50°$ pour $100°$, $40°$ pour $80°$, $65°$ pour $130°$.
\end{experience}

Cette conjecture, vérifiée expérimentalement, va maintenant être démontrée rigoureusement. C'est une démarche essentielle en mathématiques : \textbf{mesurer ne suffit pas, il faut prouver}.

\courssub{Énoncé du théorème}

\begin{propriete}[Théorème de l'angle inscrit et de l'angle au centre]
Dans un cercle, si un \cle{angle au centre} et un \cle{angle inscrit} interceptent le \textbf{même arc}, alors la mesure de l'angle au centre est le \textbf{double} de la mesure de l'angle inscrit :
\[ \widehat{AOB} = 2 \times \widehat{AMB}. \]
De façon équivalente : un angle inscrit mesure la \textbf{moitié} de l'angle au centre associé :
\[ \widehat{AMB} = \frac{1}{2}\, \widehat{AOB}. \]
\emph{Cette démonstration est exigible : elle est développée en classe et tu dois savoir la refaire.}
\end{propriete}

\begin{attention}[Condition indispensable : le même arc]
Le théorème ne s'applique \textbf{que si les deux angles interceptent le même arc} $\wideparen{AB}$. Avant d'écrire $\widehat{AOB} = 2\,\widehat{AMB}$, vérifie toujours :
\begin{itemize}
\item que $O$ est bien le \textbf{centre} du cercle ;
\item que $M$ est bien \textbf{sur le cercle} ;
\item que $M$ se trouve sur l'arc $\wideparen{AB}$ \textbf{opposé} à celui « vu » par l'angle au centre considéré.
\end{itemize}
Si $M$ est placé sur l'autre arc, la relation change : on ne peut plus appliquer directement le théorème.
\end{attention}

\courssub{Démonstration guidée du théorème}

L'idée de la démonstration est d'étudier trois cas selon la position du centre $O$ par rapport à l'angle inscrit $\widehat{AMB}$.

\begin{popfigure}
\begin{center}
\begin{tikzpicture}[scale=0.9]
\coordinate (O) at (0,0);
% Cas 1 : O sur un côté
\begin{scope}
\draw[popline] (0,0) circle (1.4);
\fill (0,0) circle (1pt) node[below left=-1pt] {$O$};
\coordinate (A) at (140:1.4);
\coordinate (B) at (220:1.4);
\coordinate (M) at (0:1.4);
\draw[popblue,thick] (M)--(A);
\draw[popblue,thick] (M)--(B);
\draw[poporange,thick] (O)--(A);
\draw[poporange,thick] (O)--(B);
\draw[popdark,dashed] (O)--(M);
\fill (A) circle (1pt) node[above left] {$A$};
\fill (B) circle (1pt) node[below left] {$B$};
\fill (M) circle (1pt) node[right] {$M$};
\node at (0,-1.9) {\textbf{Cas 1} : $O$ sur un côté};
\end{scope}
% Cas 2 : O à l'intérieur
\begin{scope}[xshift=4.6cm]
\draw[popline] (0,0) circle (1.4);
\fill (0,0) circle (1pt) node[below=1pt] {$O$};
\coordinate (A) at (150:1.4);
\coordinate (B) at (250:1.4);
\coordinate (M) at (30:1.4);
\draw[popblue,thick] (M)--(A);
\draw[popblue,thick] (M)--(B);
\draw[poporange,thick] (O)--(A);
\draw[poporange,thick] (O)--(B);
\draw[popdark,dashed] (M)--(O)--($(M)!2!(O)$);
\fill (A) circle (1pt) node[above left] {$A$};
\fill (B) circle (1pt) node[below left] {$B$};
\fill (M) circle (1pt) node[above right] {$M$};
\node at (0,-1.9) {\textbf{Cas 2} : $O$ à l'intérieur};
\end{scope}
% Cas 3 : O à l'extérieur
\begin{scope}[xshift=9.2cm]
\draw[popline] (0,0) circle (1.4);
\fill (0,0) circle (1pt) node[right=1pt] {$O$};
\coordinate (A) at (120:1.4);
\coordinate (B) at (160:1.4);
\coordinate (M) at (200:1.4);
\draw[popblue,thick] (M)--(A);
\draw[popblue,thick] (M)--(B);
\draw[poporange,thick] (O)--(A);
\draw[poporange,thick] (O)--(B);
\draw[popdark,dashed] (M)--(O)--($(M)!2!(O)$);
\fill (A) circle (1pt) node[above right] {$A$};
\fill (B) circle (1pt) node[above left] {$B$};
\fill (M) circle (1pt) node[below left] {$M$};
\node at (0,-1.9) {\textbf{Cas 3} : $O$ à l'extérieur};
\end{scope}
\end{tikzpicture}
\end{center}
\legende{Les trois cas de la démonstration selon la position du centre $O$ par rapport à l'angle inscrit $\widehat{AMB}$ (en bleu). L'angle au centre $\widehat{AOB}$ est en orange ; le diamètre auxiliaire est en pointillés.}
\end{popfigure}

\begin{experience}[Démonstration — Cas 1 : le centre $O$ est sur un côté de l'angle inscrit]
Supposons que $O$ appartient au segment $[MB]$ : alors $[MB]$ est un \textbf{diamètre}.
\begin{enumerate}
\item Dans le triangle $OMA$, on a $OA = OM$ (rayons du cercle), donc le triangle $OMA$ est \cle{isocèle} en $O$.
\item Par conséquent, les angles à la base sont égaux : $\widehat{OMA} = \widehat{OAM}$. Notons $x$ cette mesure commune.
\item La somme des angles du triangle $OMA$ vaut $180°$, donc :
\[ \widehat{MOA} = 180° - 2x. \]
\item Or $[MB]$ est un diamètre, donc $\widehat{MOA}$ et $\widehat{AOB}$ sont \textbf{supplémentaires} :
\[ \widehat{AOB} = 180° - \widehat{MOA} = 180° - (180° - 2x) = 2x. \]
\item Comme $x = \widehat{OMA} = \widehat{BMA} = \widehat{AMB}$, on obtient bien :
\[ \widehat{AOB} = 2\,\widehat{AMB}. \]
\end{enumerate}
\emph{Remarque :} à l'étape 4, on a en fait utilisé la propriété de l'\textbf{angle extérieur} d'un triangle : l'angle extérieur $\widehat{AOB}$ est égal à la somme des deux angles intérieurs non adjacents $\widehat{OMA} + \widehat{OAM} = 2x$.
\end{experience}

\begin{experience}[Démonstration — Cas 2 : le centre $O$ est à l'intérieur de l'angle inscrit]
L'astuce consiste à tracer le \textbf{diamètre $[MD]$} passant par $M$ et $O$. Ce diamètre partage l'angle inscrit en deux angles, et l'angle au centre en deux angles :
\begin{enumerate}
\item $\widehat{AMB} = \widehat{AMD} + \widehat{DMB}$ et $\widehat{AOB} = \widehat{AOD} + \widehat{DOB}$.
\item Pour l'angle $\widehat{AMD}$, le centre $O$ est sur le côté $[MD]$ : on est ramené au \textbf{cas 1}, donc $\widehat{AOD} = 2\,\widehat{AMD}$.
\item De même, pour $\widehat{DMB}$ : $\widehat{DOB} = 2\,\widehat{DMB}$.
\item En additionnant membre à membre :
\begin{align*}
\widehat{AOB} &= \widehat{AOD} + \widehat{DOB} \\
&= 2\,\widehat{AMD} + 2\,\widehat{DMB} \\
&= 2\left(\widehat{AMD} + \widehat{DMB}\right) \\
&= 2\,\widehat{AMB}.
\end{align*}
\end{enumerate}
\end{experience}

\begin{experience}[Démonstration — Cas 3 : le centre $O$ est à l'extérieur de l'angle inscrit]
On trace encore le diamètre $[MD]$, mais cette fois l'angle $\widehat{AMB}$ s'obtient par \textbf{soustraction} :
\begin{enumerate}
\item On a $\widehat{AMB} = \widehat{DMB} - \widehat{DMA}$ et $\widehat{AOB} = \widehat{DOB} - \widehat{DOA}$.
\item Le cas 1 appliqué deux fois donne : $\widehat{DOB} = 2\,\widehat{DMB}$ et $\widehat{DOA} = 2\,\widehat{DMA}$.
\item En soustrayant membre à membre :
\begin{align*}
\widehat{AOB} &= \widehat{DOB} - \widehat{DOA} \\
&= 2\,\widehat{DMB} - 2\,\widehat{DMA} \\
&= 2\left(\widehat{DMB} - \widehat{DMA}\right) \\
&= 2\,\widehat{AMB}.
\end{align*}
\end{enumerate}
\textbf{Conclusion :} dans les trois cas possibles, $\widehat{AOB} = 2\,\widehat{AMB}$. Le théorème est démontré.
\end{experience}

\courssub{Reformulation : l'angle inscrit et l'arc intercepté}

\begin{aretenir}[Trois façons de dire la même chose]
Si $\widehat{AMB}$ est un angle inscrit et $\widehat{AOB}$ l'angle au centre qui intercepte le même arc $\wideparen{AB}$, alors :
\begin{itemize}
\item l'angle au centre est le \textbf{double} de l'angle inscrit : $\widehat{AOB} = 2\,\widehat{AMB}$ ;
\item l'angle inscrit est la \textbf{moitié} de l'angle au centre : $\widehat{AMB} = \dfrac{\widehat{AOB}}{2}$ ;
\item la mesure de l'angle inscrit est la \textbf{moitié de la mesure de l'arc intercepté} $\wideparen{AB}$ (car la mesure d'un arc est, par définition, celle de l'angle au centre qui l'intercepte).
\end{itemize}
\end{aretenir}

\courssub{Applications directes : calculs d'angles}

\begin{exemplebox}[De l'angle au centre vers l'angle inscrit]
Sur un cercle de centre $O$, on donne $\widehat{AOB} = 120°$. Alors \textbf{tout} point $M$ du cercle situé sur le grand arc $\wideparen{AB}$ vérifie :
\[ \widehat{AMB} = \frac{\widehat{AOB}}{2} = \frac{120°}{2} = 60°. \]
Peu importe où l'on place $M$ sur cet arc : l'angle inscrit vaut toujours $60°$.
\end{exemplebox}

\begin{popfigure}
\begin{center}
\begin{tikzpicture}[scale=1.5]
\draw[popline] (0,0) circle (1.6);
\fill (0,0) circle (1.2pt) node[below right=-1pt] {$O$};
\coordinate (A) at (130:1.6);
\coordinate (B) at (230:1.6);
\coordinate (M) at (0:1.6);
\draw[poporange,very thick] (O)--(A);
\draw[poporange,very thick] (O)--(B);
\draw[popblue,very thick] (M)--(A);
\draw[popblue,very thick] (M)--(B);
\fill (A) circle (1.2pt) node[above left] {$A$};
\fill (B) circle (1.2pt) node[below left] {$B$};
\fill (M) circle (1.2pt) node[right] {$M$};
% Arc angle au centre
\draw[poporange,thick] (130:0.55) arc (130:230:0.55);
\node[poporange] at (180:0.85) {$100°$};
% Arc angle inscrit
\draw[popblue,thick] ($(M)+(155:0.5)$) arc (155:205:0.5);
\node[popblue] at ($(M)+(180:0.8)$) {$50°$};
\end{tikzpicture}
\end{center}
\legende{L'angle au centre $\widehat{AOB} = 100°$ (orange) est le double de l'angle inscrit $\widehat{AMB} = 50°$ (bleu) : ils interceptent le même arc $\wideparen{AB}$.}
\end{popfigure}

\begin{methode}[Calculer un angle avec le théorème de l'angle inscrit]
\begin{enumerate}
\item \textbf{Repérer l'arc intercepté} par l'angle cherché (ou connu) : identifier ses deux extrémités, par exemple $A$ et $B$.
\item \textbf{Chercher l'angle associé} qui intercepte le même arc : l'angle au centre $\widehat{AOB}$ si l'on connaît un angle inscrit, ou un angle inscrit $\widehat{AMB}$ si l'on connaît l'angle au centre.
\item \textbf{Vérifier la condition} : les deux angles interceptent bien le \emph{même} arc (le sommet $M$ est sur l'arc opposé).
\item \textbf{Appliquer le facteur 2} : multiplier par 2 (inscrit $\to$ centre) ou diviser par 2 (centre $\to$ inscrit).
\item \textbf{Conclure} par une phrase qui cite le théorème.
\end{enumerate}
\end{methode}

\begin{exoresolu}[De l'angle inscrit vers l'angle au centre]
Sur un cercle de centre $O$, les points $A$, $B$ et $M$ sont tels que $\widehat{AMB} = 35°$, avec $M$ sur le grand arc $\wideparen{AB}$. Détermine la mesure de l'angle au centre $\widehat{AOB}$ associé.
\tcblower
\textbf{Solution.}
\begin{itemize}
\item \textbf{Étape 1 :} l'angle inscrit $\widehat{AMB}$ intercepte l'arc $\wideparen{AB}$ (le petit arc, ne contenant pas $M$).
\item \textbf{Étape 2 :} l'angle au centre qui intercepte ce même arc est $\widehat{AOB}$.
\item \textbf{Étape 3 :} les deux angles interceptent bien le même arc $\wideparen{AB}$ : le théorème s'applique.
\item \textbf{Étape 4 :} d'après le théorème de l'angle inscrit :
\[ \widehat{AOB} = 2 \times \widehat{AMB} = 2 \times 35° = 70°. \]
\end{itemize}
\textbf{Conclusion :} l'angle au centre $\widehat{AOB}$ mesure $70°$.
\end{exoresolu}

\begin{exoresolu}[Sens inverse : retrouver l'angle inscrit]
Un secteur angulaire de $140°$ est tracé au centre $O$ d'un cercle ; ses côtés coupent le cercle en $P$ et $Q$. Un point $R$ est placé sur le grand arc $\wideparen{PQ}$. Calcule $\widehat{PRQ}$.
\tcblower
\textbf{Solution.} L'angle au centre est $\widehat{POQ} = 140°$. L'angle inscrit $\widehat{PRQ}$ intercepte le même arc $\wideparen{PQ}$ que $\widehat{POQ}$. D'après le théorème de l'angle inscrit :
\[ \widehat{PRQ} = \frac{\widehat{POQ}}{2} = \frac{140°}{2} = 70°. \]
\textbf{Conclusion :} $\widehat{PRQ} = 70°$.
\end{exoresolu}

\begin{attention}[Pièges fréquents]
\begin{itemize}
\item \textbf{Mauvais arc :} si $M$ est sur le \emph{petit} arc $\wideparen{AB}$, alors $\widehat{AMB}$ n'intercepte pas le même arc que l'angle au centre $\widehat{AOB}$ « direct ». Le théorème ne s'applique pas tel quel : dessine toujours une figure avant de calculer.
\item \textbf{Facteur 2 inversé :} retiens le sens logique : le centre « voit plus grand » que le point du cercle. C'est donc l'angle \emph{au centre} qui est le plus grand : $\widehat{AOB} = 2\,\widehat{AMB}$, et non l'inverse.
\item \textbf{Rédaction :} au BEPC, cite explicitement le théorème (« d'après le théorème de l'angle inscrit ») et précise l'arc commun intercepté. Une réponse numérique sans justification ne rapporte pas tous les points.
\end{itemize}
\end{attention}

\begin{exercice}[À ton tour]
\begin{enumerate}
\item Sur un cercle de centre $O$, $\widehat{AOB} = 86°$. Calcule la mesure de tout angle inscrit interceptant le même arc $\wideparen{AB}$.
\item Sur un cercle de centre $O$, un angle inscrit $\widehat{CMD}$ mesure $27°$. Calcule l'angle au centre $\widehat{COD}$ associé.
\item Refais la démonstration du cas 1 en rédigeant toutes les justifications (triangle isocèle, somme des angles, angles supplémentaires).
\end{enumerate}
\end{exercice}

\begin{savaistu}[Un théorème vieux de presque 4000 ans]
La relation entre l'angle inscrit et l'angle au centre était déjà connue des \textbf{Babyloniens}, qui l'utilisaient dans leurs tablettes d'astronomie pour prévoir les positions des astres. Elle fut démontrée rigoureusement par \textbf{Euclide} vers 300 av. J.-C. : c'est la \emph{proposition 20 du livre III} de ses célèbres \emph{Éléments}, l'un des livres les plus lus de toute l'histoire de l'humanité. Aujourd'hui encore, ce théorème sert en astronomie, en navigation et dans les systèmes de positionnement : lorsqu'un géomètre camerounais mesure des angles depuis différents points du terrain pour cartographier une parcelle à Yaoundé ou à Garoua, il utilise — sans toujours le savoir — ce résultat vieux de plus de deux mille ans !
\end{savaistu}

\begin{aretenir}[L'essentiel de la section]
\begin{itemize}
\item \textbf{Théorème de l'angle inscrit :} dans un cercle, un angle au centre mesure le double de tout angle inscrit qui intercepte le même arc : $\widehat{AOB} = 2\,\widehat{AMB}$.
\item La démonstration repose sur le \textbf{triangle isocèle} $OMA$ (deux rayons) et se traite en trois cas selon la position de $O$ ; l'astuce des cas 2 et 3 est de tracer le \textbf{diamètre} passant par $M$.
\item Un angle inscrit mesure la \textbf{moitié de l'arc} qu'il intercepte.
\item Condition d'application : vérifier que les deux angles interceptent \textbf{le même arc}.
\end{itemize}
Dans la section suivante, nous verrons une conséquence immédiate et très utile : deux angles inscrits qui interceptent le même arc sont \textbf{égaux}.
\end{aretenir}

\coursec{Angles inscrits interceptant le même arc}

Dans la section précédente, nous avons établi un résultat fondamental : dans un cercle, un \cle{angle inscrit} mesure la moitié de l'\cle{angle au centre} qui intercepte le même arc. Autrement dit, si $\widehat{AMB}$ et $\widehat{AOB}$ interceptent le même arc $\wideparen{AB}$, alors $\widehat{AMB} = \dfrac{1}{2}\widehat{AOB}$.

Faisons maintenant une remarque en apparence anodine, mais aux conséquences spectaculaires : l'angle au centre $\widehat{AOB}$ est \emph{unique} pour un arc donné, alors qu'on peut placer le sommet $M$ d'un angle inscrit \emph{n'importe où} sur l'autre arc. Que se passe-t-il si l'on choisit deux sommets différents $M$ et $N$ ? C'est toute la richesse de cette section.

\courssub{Une observation qui change tout : l'angle de vision constant}

\begin{experience}[Activité de découverte]
Trace un cercle de centre $O$ et place deux points $A$ et $B$ sur ce cercle. Ces deux points délimitent deux arcs. Choisis ensuite trois points $M$, $N$ et $P$ sur \textbf{l'un de ces deux arcs} (par exemple l'arc qui ne contient pas le centre $O$ si $A$ et $B$ ne sont pas diamétralement opposés).
\begin{enumerate}
\item Mesure au rapporteur les angles $\widehat{AMB}$, $\widehat{ANB}$ et $\widehat{APB}$.
\item Compare les trois mesures obtenues. Que constates-tu ?
\item Mesure l'angle au centre $\widehat{AOB}$ (celui qui intercepte l'arc \emph{opposé} à $M$, $N$, $P$) et compare-le à chacun des angles inscrits.
\end{enumerate}
Tu dois constater que les trois angles inscrits sont \textbf{égaux}, et que chacun vaut la moitié de l'angle au centre $\widehat{AOB}$.
\end{experience}

Cette expérience révèle un phénomène remarquable : tant que le sommet $M$ se déplace \emph{sur le même arc} (délimité par $A$ et $B$), l'angle $\widehat{AMB}$ ne change pas. On dit que la corde $[AB]$ est \emph{vue sous un angle constant} depuis tous les points de cet arc. C'est exactement ce qui se passe dans un stade : deux supporters assis sur la même tribune en arc de cercle voient les deux poteaux de but sous le même angle.

\begin{popfigure}
\begin{center}
\begin{tikzpicture}[scale=1.6]
\coordinate (O) at (0,0);
\draw[popline, thick] (O) circle (2);
\coordinate (A) at (210:2);
\coordinate (B) at (-30:2);
\coordinate (M) at (90:2);
\coordinate (N) at (55:2);
\coordinate (P) at (125:2);
\draw[popblue, thick] (A)--(M)--(B);
\draw[poporange, thick] (A)--(N)--(B);
\draw[poppurple, thick] (A)--(P)--(B);
\draw[popdark, thick] (A)--(B);
\foreach \p/\pos/\col in {A/below left/popdark, B/below right/popdark, M/above/popblue, N/above right/poporange, P/above left/poppurple}
  \fill[\col] (\p) circle (1.3pt) node[\pos, \col] {\footnotesize $\p$};
\fill[popdark] (O) circle (1.3pt) node[below] {\footnotesize $O$};
% Arc intercepted (the bottom one)
\draw[popgreen, very thick, decoration={markings, mark=at position 0.5 with {\arrow{>}}}, postaction={decorate}] (210:2) arc (210:-30:2);
\node[popgreen, font=\footnotesize] at (0,-1.5) {arc intercept\'e};
\end{tikzpicture}
\end{center}
\legende{Trois angles inscrits $\widehat{AMB}$, $\widehat{ANB}$ et $\widehat{APB}$ ont leurs sommets sur le m\^eme arc (le grand arc $AB$ en haut) : ils interceptent tous le m\^eme arc (le petit arc $AB$ en bas) et sont donc \'egaux.}
\end{popfigure}

\courssub{La propri\'et\'e fondamentale}

\begin{propriete}[Angles inscrits interceptant le m\^eme arc]
Dans un cercle, deux \cle{angles inscrits} qui interceptent le \cle{m\^eme arc} ont la \cle{m\^eme mesure}.

Autrement dit : si $M$ et $N$ sont deux points d'un m\^eme arc $\wideparen{AB}$ d'un cercle (c'est-\`a-dire du m\^eme c\^ot\'e de la corde $[AB]$), alors les angles $\widehat{AMB}$ et $\widehat{ANB}$ interceptent l'autre arc $\wideparen{AB}$ et l'on a :
\[ \widehat{AMB} = \widehat{ANB}. \]
\end{propriete}

\begin{experience}[D\'emonstration de la propri\'et\'e]
Cette d\'emonstration est une cons\'equence directe du th\'eor\`eme de l'angle inscrit et de l'angle au centre. Suivons le raisonnement pas \`a pas.

\textbf{Hypoth\`eses.} Soit un cercle de centre $O$, et soient $M$ et $N$ deux points situ\'es sur le m\^eme arc d\'elimit\'e par $A$ et $B$. Les angles $\widehat{AMB}$ et $\widehat{ANB}$ interceptent donc le m\^eme arc (l'arc oppos\'e).

\textbf{\'Etape 1.} L'angle $\widehat{AMB}$ est un angle inscrit qui intercepte l'arc $\wideparen{AB}$ (l'arc ne contenant pas $M$). D'apr\`es le th\'eor\`eme de l'angle inscrit :
\[ \widehat{AMB} = \frac{1}{2}\,\widehat{AOB}. \]

\textbf{\'Etape 2.} De m\^eme, l'angle $\widehat{ANB}$ est un angle inscrit qui intercepte \emph{le m\^eme} arc $\wideparen{AB}$ (l'arc ne contenant pas $N$). D'apr\`es le m\^eme th\'eor\`eme :
\[ \widehat{ANB} = \frac{1}{2}\,\widehat{AOB}. \]

\textbf{\'Etape 3.} Les deux angles sont \'egaux \`a la m\^eme quantit\'e $\dfrac{1}{2}\widehat{AOB}$. Par transitivit\'e de l'\'egalit\'e :
\[ \widehat{AMB} = \widehat{ANB}. \]
La propri\'et\'e est d\'emontr\'ee. $\blacksquare$
\end{experience}

\begin{aretenir}[L'id\'ee \`a retenir]
L'angle au centre $\widehat{AOB}$ associ\'e \`a un arc est \emph{fixe}. Tous les angles inscrits qui interceptent cet arc valent sa moiti\'e : ils sont donc \emph{tous \'egaux entre eux}. Retiens l'image : « depuis un m\^eme arc, on voit une corde toujours sous le m\^eme angle ».
\end{aretenir}

\begin{exemplebox}[Exemple chiffr\'e]
Sur un cercle de centre $O$, les points $A$, $B$, $M$ et $N$ sont plac\'es de sorte que $M$ et $N$ appartiennent au m\^eme arc d\'elimit\'e par $A$ et $B$. On donne $\widehat{AMB} = 35^\circ$.

Comme $\widehat{ANB}$ est un angle inscrit qui intercepte le m\^eme arc que $\widehat{AMB}$, on a :
\[ \widehat{ANB} = \widehat{AMB} = 35^\circ. \]
V\'erification avec l'angle au centre : $\widehat{AOB} = 2 \times 35^\circ = 70^\circ$, et l'on a bien $\widehat{ANB} = \dfrac{70^\circ}{2} = 35^\circ$.
\end{exemplebox}

\courssub{Cas particulier : arcs de m\^eme mesure}

La propri\'et\'e se g\'en\'eralise naturellement : dans un m\^eme cercle, des arcs de \emph{m\^eme mesure} sont intercept\'es par des angles inscrits \emph{\'egaux}, m\^eme si les angles n'interceptent pas le m\^eme arc.

\begin{propriete}[Angles inscrits et arcs de m\^eme mesure]
Dans un m\^eme cercle, deux angles inscrits qui interceptent des arcs de m\^eme mesure sont \'egaux.
\end{propriete}

En effet, si les arcs $\wideparen{AB}$ et $\wideparen{CD}$ ont la m\^eme mesure, alors les angles au centre associ\'es sont \'egaux : $\widehat{AOB} = \widehat{COD}$, et leurs moiti\'es aussi. Ce cas particulier est tr\`es utile d\`es qu'une figure comporte des cordes de m\^eme longueur, car dans un m\^eme cercle, des cordes \'egales sous-tendent des arcs \'egaux.

\courssub{Interpr\'etation concr\`ete : l'angle de vision constant}

\begin{popfigure}
\begin{center}
\begin{tikzpicture}[scale=1.5]
\coordinate (O) at (0,0);
\draw[popline, thick] (O) circle (2);
\coordinate (A) at (200:2);
\coordinate (B) at (-20:2);
\coordinate (M) at (80:2);
\coordinate (N) at (110:2);
\draw[popblue, thick] (A)--(M)--(B);
\draw[poporange, thick] (A)--(N)--(B);
\draw[popdark, very thick] (A)--(B);
\foreach \p/\pos/\col in {A/below left/popdark, B/below right/popdark, M/above right/popblue, N/above left/poporange}
  \fill[\col] (\p) circle (1.3pt) node[\pos, \col] {\footnotesize $\p$};
\node[popdark, font=\footnotesize] at (0,-1.55) {poteaux $A$ et $B$};
\node[popblue, font=\footnotesize] at (1.35,1.75) {supporter $M$};
\node[poporange, font=\footnotesize] at (-1.35,1.75) {supporter $N$};
\fill[popdark] (O) circle (1.2pt) node[below] {\footnotesize $O$};
\end{tikzpicture}
\end{center}
\legende{Au stade, deux supporters $M$ et $N$ plac\'es sur la m\^eme tribune en arc de cercle voient les poteaux $A$ et $B$ sous le m\^eme angle : $\widehat{AMB} = \widehat{ANB}$.}
\end{popfigure}

Imagine le stade Ahmadou Ahidjo de Yaound\'e : la tribune forme un arc de cercle, et les deux poteaux du but sont les points $A$ et $B$. Un supporter assis au point $M$ voit le but sous l'angle $\widehat{AMB}$ ; un autre supporter, assis au point $N$ de la m\^eme tribune, le voit sous l'angle $\widehat{ANB}$. La propri\'et\'e des angles inscrits garantit que ces deux angles sont \emph{\'egaux} : tous les spectateurs de la m\^eme tribune circulaire ont exactement le m\^eme angle de vision sur le but. C'est une des raisons pour lesquelles les stades et les th\'e\^atres antiques sont construits en arc de cercle !

\courssub{Application : d\'emontrer que deux angles sont \'egaux}

\begin{methode}[Prouver que deux angles sont \'egaux gr\^ace au cercle]
Pour montrer que deux angles $\widehat{AMB}$ et $\widehat{ANB}$ sont \'egaux dans une figure comportant un cercle :
\begin{enumerate}
\item V\'erifier que les quatre points $A$, $M$, $N$, $B$ appartiennent au \textbf{m\^eme cercle} (les angles sont bien inscrits).
\item V\'erifier que $M$ et $N$ sont situ\'es \textbf{du m\^eme c\^ot\'e} de la corde $[AB]$, c'est-\`a-dire sur le m\^eme arc d\'elimit\'e par $A$ et $B$.
\item Identifier l'arc commun intercept\'e : ici l'arc $\wideparen{AB}$ ne contenant pas $M$ et $N$.
\item Conclure par la propri\'et\'e : $\widehat{AMB} = \widehat{ANB}$.
\end{enumerate}
\end{methode}

\begin{exoresolu}[Montrer que deux angles sont \'egaux]
\textbf{\'Enonc\'e.} Soit un cercle de centre $O$ passant par quatre points $A$, $B$, $C$ et $D$, plac\'es dans cet ordre sur le cercle. On donne $\widehat{ACB} = 28^\circ$. D\'eterminer la mesure de l'angle $\widehat{ADB}$, en justifiant.
\tcblower
\textbf{Solution d\'etaill\'ee.}

\textbf{\'Etape 1.} Les points $A$, $B$, $C$, $D$ sont sur le m\^eme cercle : les angles $\widehat{ACB}$ et $\widehat{ADB}$ sont des angles inscrits.

\textbf{\'Etape 2.} Les points $C$ et $D$ sont cons\'ecutifs sur le cercle (ordre $A, B, C, D$), donc situ\'es du m\^eme c\^ot\'e de la corde $[AB]$ : les deux angles interceptent le m\^eme arc $\wideparen{AB}$ (l'arc ne contenant ni $C$ ni $D$).

\textbf{\'Etape 3.} D'apr\`es la propri\'et\'e des angles inscrits interceptant le m\^eme arc :
\[ \widehat{ADB} = \widehat{ACB} = 28^\circ. \]

\textbf{Conclusion.} L'angle $\widehat{ADB}$ mesure $28^\circ$.
\end{exoresolu}

\begin{exoresolu}[Avec des arcs de m\^eme mesure]
\textbf{\'Enonc\'e.} Sur un cercle de centre $O$, on place quatre points $A$, $B$, $C$, $D$ tels que les cordes $[AB]$ et $[CD]$ ont la m\^eme longueur. On donne $\widehat{AMB} = 40^\circ$, o\`u $M$ est un point du cercle. Que vaut l'angle $\widehat{CND}$, o\`u $N$ est un point du cercle situ\'e de fa\c{c}on que $\widehat{CND}$ intercepte l'arc $\wideparen{CD}$ ?
\tcblower
\textbf{Solution d\'etaill\'ee.}

Dans un m\^eme cercle, deux cordes de m\^eme longueur sous-tendent des arcs de m\^eme mesure : les arcs $\wideparen{AB}$ et $\wideparen{CD}$ ont donc la m\^eme mesure.

Les angles inscrits $\widehat{AMB}$ et $\widehat{CND}$ interceptent des arcs de m\^eme mesure : ils sont \'egaux.
\[ \widehat{CND} = \widehat{AMB} = 40^\circ. \]
\end{exoresolu}

\courssub{Le pi\`ege \`a \'eviter : les sommets de part et d'autre de la corde}

\begin{attention}[M\^eme corde ne signifie pas m\^eme arc !]
Si les sommets $M$ et $N$ sont situ\'es \textbf{de part et d'autre} de la corde $[AB]$, alors les angles $\widehat{AMB}$ et $\widehat{ANB}$ n'interceptent \textbf{pas le m\^eme arc} : l'un intercepte l'arc $\wideparen{AB}$ « d'en face » de $M$, l'autre l'arc « d'en face » de $N$. Ces deux arcs sont \emph{compl\'ementaires} : r\'eunis, ils forment le cercle entier.

Dans ce cas, les angles ne sont pas \'egaux mais \textbf{suppl\'ementaires} :
\[ \widehat{AMB} + \widehat{ANB} = 180^\circ. \]
En effet, si $\widehat{AMB} = \dfrac{1}{2}\widehat{AOB}$ (angle au centre interceptant un arc) et $\widehat{ANB} = \dfrac{1}{2}\widehat{AOB'}$ (angle au centre interceptant l'autre arc), alors la somme vaut la moiti\'e de $360^\circ$, c'est-\`a-dire $180^\circ$.

Avant de conclure \`a l'\'egalit\'e de deux angles inscrits, v\'erifie toujours que leurs sommets sont \textbf{du m\^eme c\^ot\'e} de la corde !
\end{attention}

Ce r\'esultat sur les angles suppl\'ementaires est une premi\`ere approche d'une notion que tu approfondiras au lyc\'ee : les \emph{quadrilat\`eres inscriptibles}, dont les angles oppos\'es sont suppl\'ementaires.

\courssub{Vers la notion de points cocycliques}

\begin{definition}[Points cocycliques]
Des points sont dits \cle{cocycliques} s'ils appartiennent tous \`a un \cle{m\^eme cercle}.
\end{definition}

La propri\'et\'e des angles inscrits fonctionne dans les deux sens, et c'est un outil de d\'emonstration tr\`es puissant :
\begin{itemize}
\item \textbf{Sens direct} : si quatre points $A$, $B$, $M$, $N$ sont cocycliques avec $M$ et $N$ du m\^eme c\^ot\'e de $[AB]$, alors $\widehat{AMB} = \widehat{ANB}$.
\item \textbf{Sens r\'eciproque (admis / d\'emontrable par l'absurde)} : si deux points $M$ et $N$ situ\'es du m\^eme c\^ot\'e d'un segment $[AB]$ v\'erifient $\widehat{AMB} = \widehat{ANB}$, alors les quatre points $A$, $B$, $M$, $N$ sont cocycliques.
\end{itemize}

Autrement dit, l'\'egalit\'e de deux angles de vision est une \emph{signature} de la cocyclic\'it\'e : elle permet de prouver qu'un cercle passe par quatre points, ce qui d\'ebloque ensuite toutes les propri\'et\'es du cercle (angles au centre, triangle inscrit dans un demi-cercle, etc.). Tu utiliseras cette strat\'egie dans les d\'emonstrations de la section 5.

\begin{savaistu}[Pourquoi tous les rempla\c{c}ants voient-ils le match de la m\^eme fa\c{c}on ?]
Sur un banc de touche dispos\'e en arc de cercle — comme dans certains stades modernes ou les th\'e\^atres grecs antiques — tous les rempla\c{c}ants voient l'action sous le m\^eme angle, car leurs positions sont des points d'un m\^eme arc et la sc\`ene joue le r\^ole de la corde. Les architectes grecs appliquaient d\'ej\`a, il y a plus de deux mille ans, cette propri\'et\'e des angles inscrits pour garantir \`a chaque spectateur la m\^eme qualit\'e de vision. Les math\'ematiques de la classe de 3\`eme sont litt\'eralement inscrites dans la pierre des amphith\'e\^atres !
\end{savaistu}

\begin{exercice}[\`A toi de jouer]
Sur un cercle de centre $O$, on place cinq points $A$, $B$, $M$, $N$, $P$ dans cet ordre. On donne $\widehat{APB} = 52^\circ$.
\begin{enumerate}
\item D\'eterminer $\widehat{AMB}$ et $\widehat{ANB}$ en justifiant.
\item On place un point $Q$ sur l'arc $\wideparen{AB}$ ne contenant pas $M$, $N$, $P$. D\'eterminer $\widehat{AQB}$.
\end{enumerate}
\end{exercice}

\begin{corrige}[Exercice : \`A toi de jouer]
\begin{enumerate}
\item Les angles $\widehat{AMB}$, $\widehat{ANB}$ et $\widehat{APB}$ sont inscrits dans le m\^eme cercle et leurs sommets $M$, $N$, $P$ sont sur le m\^eme arc par rapport \`a la corde $[AB]$ (l'ordre $A, B, M, N, P$ place $M, N, P$ sur l'arc $AB$ ne contenant pas l'autre arc) : ils interceptent le m\^eme arc $\wideparen{AB}$ (l'arc oppos\'e). Donc $\widehat{AMB} = \widehat{ANB} = \widehat{APB} = 52^\circ$.
\item Le point $Q$ est de l'autre c\^ot\'e de la corde $[AB]$ (sur l'arc ne contenant pas $M, N, P$) : $\widehat{AQB}$ et $\widehat{APB}$ interceptent des arcs compl\'ementaires, ils sont suppl\'ementaires : $\widehat{AQB} = 180^\circ - 52^\circ = 128^\circ$.
\end{enumerate}
\end{corrige}

\coursec{Triangle inscrit dans un demi-cercle}

Dans la section précédente, nous avons découvert que deux angles inscrits qui interceptent le même arc ont la même mesure. Nous allons maintenant étudier le cas le plus célèbre et le plus utile de tous : celui où l'arc intercepté est un \cle{demi-cercle}, c'est-à-dire un arc de $180°$. Le résultat que nous allons établir est l'un des plus anciens théorèmes de l'histoire des mathématiques, et tu vas voir qu'il sert encore aujourd'hui, y compris sur les chantiers de construction au Cameroun.

\courssub{Une observation pour commencer}

Prends une feuille, trace un cercle de centre $O$, puis un diamètre $[AB]$ de ce cercle. Choisis maintenant un point $M$ sur le cercle, différent de $A$ et de $B$, et trace le triangle $AMB$. Mesure l'angle $\widehat{AMB}$ avec ton rapporteur. Recommence avec un autre point $M'$ sur le cercle. Que constates-tu ?

Quelle que soit la position du point $M$ sur le cercle, l'angle $\widehat{AMB}$ mesure toujours $90°$ : le triangle $AMB$ semble toujours \cle{rectangle en} $M$. Ce n'est pas un hasard, et nous allons le démontrer.

\begin{popfigure}
\begin{center}
\begin{tikzpicture}[scale=1.1]
  \coordinate (O) at (0,0);
  \coordinate (A) at (-2.5,0);
  \coordinate (B) at (2.5,0);
  \coordinate (M) at (1.5,2);
  \draw[popline,thick] (O) circle (2.5);
  \draw[popblue,thick] (A) -- (B);
  \draw[poporange,thick] (A) -- (M) -- (B);
  \fill[popdark] (O) circle (1.5pt) node[below=2pt] {$O$};
  \fill[popdark] (A) circle (1.5pt) node[left=2pt] {$A$};
  \fill[popdark] (B) circle (1.5pt) node[right=2pt] {$B$};
  \fill[popdark] (M) circle (1.5pt) node[above=2pt] {$M$};
  % codage angle droit en M
  \draw[popgreen,thick] ($(M)!0.28cm!(A)$) -- ($(M)!0.28cm!(A) + (M)!0.28cm!(B) - (M)$) -- ($(M)!0.28cm!(B)$);
  \node[popblue] at (0,-0.35) {diamètre};
\end{tikzpicture}
\end{center}
\legende{Le point $M$ est sur le cercle de diamètre $[AB]$ : l'angle $\widehat{AMB}$ est droit.}
\end{popfigure}

\courssub{La propriété fondamentale}

\begin{propriete}[Triangle inscrit dans un demi-cercle]
Si un triangle est inscrit dans un cercle et si l'un de ses côtés est un \cle{diamètre} de ce cercle, alors ce triangle est \cle{rectangle}, et l'angle droit est l'angle opposé au diamètre.

Autrement dit : si $[AB]$ est un diamètre d'un cercle $\mathcal{C}$ et si $M$ est un point de $\mathcal{C}$ distinct de $A$ et $B$, alors le triangle $AMB$ est rectangle en $M$, c'est-à-dire $\widehat{AMB} = 90°$.
\end{propriete}

Cette propriété n'est rien d'autre qu'un cas particulier du théorème de l'angle inscrit étudié dans cette leçon. La démonstration est courte et tu dois savoir la rédiger.

\begin{experience}[Démonstration guidée de la propriété]
On considère un cercle de centre $O$, un diamètre $[AB]$ et un point $M$ du cercle distinct de $A$ et $B$.
\begin{enumerate}
  \item \textbf{Repérer l'angle au centre associé.} L'angle inscrit $\widehat{AMB}$ intercepte l'arc $\overset{\frown}{AB}$ qui ne contient pas $M$. L'angle au centre qui intercepte ce même arc est $\widehat{AOB}$.
  \item \textbf{Mesurer l'angle au centre.} Comme $[AB]$ est un diamètre, les points $A$, $O$ et $B$ sont alignés, avec $O$ entre $A$ et $B$. L'angle $\widehat{AOB}$ est donc un angle \cle{plat} : $\widehat{AOB} = 180°$.
  \item \textbf{Appliquer le théorème de l'angle inscrit.} L'angle inscrit vaut la moitié de l'angle au centre associé :
  \[ \widehat{AMB} = \frac{1}{2}\,\widehat{AOB} = \frac{1}{2}\times 180° = 90°. \]
  \item \textbf{Conclure.} Le triangle $AMB$ est rectangle en $M$. \hfill$\blacksquare$
\end{enumerate}
\end{experience}

\begin{attention}[Bien vérifier que le côté est un diamètre]
La propriété ne s'applique que si le côté du triangle est réellement un \textbf{diamètre} du cercle, c'est-à-dire s'il \textbf{passe par le centre}. Une simple corde, même très longue, ne suffit pas. Dans un exercice, cherche dans l'énoncé une phrase du type « $[AB]$ est un diamètre du cercle » ou « $O$ est le milieu de $[AB]$ et $O$ est le centre du cercle ». Si le segment ne passe pas par le centre, l'angle inscrit n'a aucune raison d'être droit.
\end{attention}

\courssub{La propriété réciproque}

Nous venons de voir : « diamètre $\Rightarrow$ angle droit ». La réciproque est également vraie, et elle est tout aussi précieuse.

\begin{propriete}[Réciproque]
Si un triangle est \cle{rectangle}, alors son cercle circonscrit a pour \cle{diamètre} son hypoténuse.

Autrement dit : si le triangle $AMB$ est rectangle en $M$, alors le cercle circonscrit au triangle $AMB$ est le cercle de diamètre $[AB]$.
\end{propriete}

\textbf{Justification.} Soit $O$ le milieu de l'hypoténuse $[AB]$. Construisons le symétrique $M'$ du point $M$ par rapport à $O$. Le quadrilatère $AMBM'$ a ses diagonales $[AB]$ et $[MM']$ qui se coupent en leur milieu $O$ : c'est un parallélogramme. Comme il possède un angle droit (en $M$), c'est un \cle{rectangle}. Or les diagonales d'un rectangle ont la même longueur : $AB = MM'$, donc $OA = OB = OM = OM'$. Le point $M$ est donc sur le cercle de centre $O$ et de rayon $OA$, c'est-à-dire le cercle de diamètre $[AB]$. \hfill$\blacksquare$

\begin{aretenir}[Le lien avec la propriété vue en 4ème]
Tu as appris en classe de 4ème que « le centre du cercle circonscrit à un triangle rectangle est le milieu de son hypoténuse ». C'est exactement la même propriété, formulée autrement :
\begin{itemize}
  \item l'hypoténuse est un \textbf{diamètre} du cercle circonscrit ;
  \item le milieu de l'hypoténuse est le \textbf{centre} du cercle circonscrit ;
  \item le rayon du cercle circonscrit vaut la \textbf{moitié de l'hypoténuse}.
\end{itemize}
Retiens aussi la conséquence pratique : dans un triangle rectangle, la médiane issue de l'angle droit mesure la moitié de l'hypoténuse (car $OM = OA = \dfrac{AB}{2}$).
\end{aretenir}

\begin{methode}[Pour montrer qu'un angle est droit ou qu'un triangle est rectangle]
Pour démontrer qu'un triangle $AMB$ est rectangle en $M$ dans une configuration avec un cercle :
\begin{enumerate}
  \item Vérifier que $A$, $M$ et $B$ sont sur un même cercle (souvent donné par l'énoncé).
  \item Vérifier que $[AB]$ est un \textbf{diamètre} de ce cercle (il passe par le centre).
  \item Conclure avec la phrase de justification : « le triangle $AMB$ est inscrit dans un cercle dont $[AB]$ est un diamètre, donc $AMB$ est rectangle en $M$ ».
\end{enumerate}
Réciproquement, pour trouver le cercle circonscrit à un triangle rectangle : son centre est le milieu de l'hypoténuse et son rayon est la moitié de l'hypoténuse.
\end{methode}

\courssub{Exemple chiffré et exercice résolu}

\begin{exemplebox}[Calcul d'une longueur grâce à l'angle droit]
Soit $\mathcal{C}$ un cercle de diamètre $[AB]$ avec $AB = 10$ cm. On place sur $\mathcal{C}$ un point $M$ tel que $AM = 6$ cm. Calculons $MB$.

\textbf{Étape 1 : justifier l'angle droit.} Le triangle $AMB$ est inscrit dans le cercle $\mathcal{C}$ et $[AB]$ est un diamètre de $\mathcal{C}$. Donc $AMB$ est rectangle en $M$.

\textbf{Étape 2 : appliquer le théorème de Pythagore.} Dans le triangle $AMB$ rectangle en $M$ :
\begin{align*}
AB^2 &= AM^2 + MB^2\\
10^2 &= 6^2 + MB^2\\
MB^2 &= 100 - 36 = 64\\
MB &= \sqrt{64} = 8 \text{ cm}.
\end{align*}
Le triangle $AMB$ est le célèbre triangle « $6$-$8$-$10$ », multiple du triangle $3$-$4$-$5$.
\end{exemplebox}

\begin{exoresolu}[Démontrer qu'un triangle est rectangle]
\textbf{Énoncé.} Soit $\mathcal{C}$ un cercle de centre $O$ et de rayon $5$ cm. $[PQ]$ est un diamètre de $\mathcal{C}$ et $R$ est un point de $\mathcal{C}$ tel que $PR = 8$ cm.
\begin{enumerate}
  \item Démontrer que le triangle $PQR$ est rectangle en $R$.
  \item Calculer $QR$.
  \item Calculer la mesure de l'angle $\widehat{RPQ}$, arrondie au degré.
\end{enumerate}
\tcblower
\textbf{Solution.}
\begin{enumerate}
  \item Le triangle $PQR$ est inscrit dans le cercle $\mathcal{C}$, et $[PQ]$ est un diamètre de $\mathcal{C}$. Or, si un triangle est inscrit dans un cercle et a pour côté un diamètre de ce cercle, alors il est rectangle. Donc $PQR$ est rectangle en $R$.
  \item Le diamètre vaut $PQ = 2 \times 5 = 10$ cm. D'après le théorème de Pythagore dans $PQR$ rectangle en $R$ :
  \[ QR^2 = PQ^2 - PR^2 = 10^2 - 8^2 = 100 - 64 = 36, \quad \text{donc } QR = 6 \text{ cm}. \]
  \item Dans le triangle rectangle en $R$, $\cos\widehat{RPQ} = \dfrac{PR}{PQ} = \dfrac{8}{10} = 0{,}8$. À la calculatrice : $\widehat{RPQ} \approx 37°$.
\end{enumerate}
\end{exoresolu}

\begin{attention}[Ne pas confondre propriété et réciproque]
Les deux phrases suivantes sont différentes, et il faut choisir la bonne selon la question :
\begin{itemize}
  \item « $[AB]$ est un diamètre, donc le triangle est rectangle » : c'est la propriété \textbf{directe}, elle sert à prouver un \emph{angle droit}.
  \item « Le triangle est rectangle, donc son hypoténuse est un diamètre du cercle circonscrit » : c'est la \textbf{réciproque}, elle sert à trouver le \emph{cercle circonscrit} (centre et rayon).
\end{itemize}
Dans ta rédaction au BEPC, cite toujours les hypothèses avant de conclure : « Comme $M$ appartient au cercle de diamètre $[AB]$, alors... ».
\end{attention}

\courssub{Application : construire un angle droit sans équerre}

La propriété réciproque fournit un moyen élégant de construire une perpendiculaire à une droite en un point, avec seulement un compas et une règle, sans équerre ni rapporteur.

\begin{methode}[Construire une perpendiculaire à une droite en un point $A$]
Pour tracer la perpendiculaire en $A$ à une droite $(d)$ :
\begin{enumerate}
  \item Placer un point $O$ qui n'est pas sur $(d)$, et tracer le cercle de centre $O$ passant par $A$.
  \item Ce cercle recoupe $(d)$ en un deuxième point $B$.
  \item Tracer le diamètre $[BC]$ : tracer la droite $(BO)$, qui recoupe le cercle en $C$. Alors $[BC]$ est un diamètre.
  \item Tracer la droite $(AC)$. Le triangle $ABC$ est inscrit dans le cercle et $[BC]$ est un diamètre, donc il est rectangle en $A$ : la droite $(AC)$ est perpendiculaire à $(d)$ en $A$.
\end{enumerate}
\end{methode}

\begin{popfigure}
\begin{center}
\begin{tikzpicture}[scale=1.05]
  \coordinate (A) at (-1.2,0);
  \coordinate (B) at (2.2,0);
  \coordinate (O) at (0.5,1.1);
  % C = 2O - B
  \coordinate (C) at (-1.2,2.2);
  \draw[popline,thick] (O) circle (2.42);
  \draw[popdark,thick] (-2.2,0) -- (3.2,0) node[right] {$(d)$};
  \draw[popblue,thick] (B) -- (C);
  \draw[poporange,very thick] (A) -- (C);
  \draw[popmuted] (A) -- (B);
  \fill[popdark] (A) circle (1.5pt) node[below left=1pt] {$A$};
  \fill[popdark] (B) circle (1.5pt) node[below right=1pt] {$B$};
  \fill[popdark] (C) circle (1.5pt) node[above left=1pt] {$C$};
  \fill[popdark] (O) circle (1.5pt) node[right=2pt] {$O$};
  \draw[popgreen,thick] ($(A)+(0,0.25)$) -- ($(A)+(-0.25,0.25)$) -- ($(A)+(-0.25,0)$);
  \node[popblue] at (1.4,1.6) {diamètre};
\end{tikzpicture}
\end{center}
\legende{$[BC]$ est un diamètre et $A$ est sur le cercle : l'angle $\widehat{BAC}$ est droit, donc $(AC) \perp (d)$.}
\end{popfigure}

\begin{savaistu}[Thalès de Milet et la corde des maçons]
Ce résultat est traditionnellement attribué à \cle{Thalès de Milet}, mathématicien grec du VI\textsuperscript{e} siècle avant J.-C., l'un des sept sages de la Grèce antique. La légende raconte qu'il émerveilla les Égyptiens en mesurant la hauteur des pyramides grâce à leur ombre. Aujourd'hui encore, les maçons camerounais utilisent l'idée de l'angle droit inscrit dans un demi-cercle — souvent sous la forme du triangle $3$-$4$-$5$ — avec une simple corde nouée pour tracer des angles droits parfaits sur les chantiers de Yaoundé, Douala ou Bafoussam, et vérifier que les murs d'une maison sont bien « d'équerre ». Les mathématiques que tu apprends en 3ème construisent littéralement les maisons !
\end{savaistu}

\begin{exercice}[À toi de jouer]
Soit $\mathcal{C}$ un cercle de centre $I$ et de diamètre $[EF]$ tel que $EF = 13$ cm. $G$ est un point de $\mathcal{C}$ tel que $EG = 5$ cm.
\begin{enumerate}
  \item Faire une figure.
  \item Démontrer que le triangle $EFG$ est rectangle, et préciser en quel point.
  \item Calculer $FG$.
  \item Calculer la distance $IG$ et justifier ta réponse.
\end{enumerate}
\end{exercice}

\begin{corrige}[Exercice : triangle inscrit dans un demi-cercle]
\begin{enumerate}
  \item Figure : cercle de centre $I$, $[EF]$ passant par $I$, $G$ sur le cercle, triangle $EFG$ tracé.
  \item Le triangle $EFG$ est inscrit dans le cercle $\mathcal{C}$ et $[EF]$ est un diamètre de $\mathcal{C}$. Or tout triangle inscrit dans un cercle dont un côté est un diamètre est rectangle. Donc $EFG$ est rectangle en $G$.
  \item D'après le théorème de Pythagore : $FG^2 = EF^2 - EG^2 = 13^2 - 5^2 = 169 - 25 = 144$, donc $FG = 12$ cm.
  \item $G$ est sur le cercle de centre $I$, donc $IG$ est un rayon : $IG = \dfrac{EF}{2} = \dfrac{13}{2} = 6{,}5$ cm. On retrouve la propriété : la médiane issue de l'angle droit vaut la moitié de l'hypoténuse.
\end{enumerate}
\end{corrige}

\coursec{Calculs d'angles et démonstrations}

Dans les sections précédentes, nous avons établi les résultats fondamentaux : le théorème de l'angle inscrit, l'égalité des angles inscrits interceptant le même arc, et la propriété remarquable du triangle inscrit dans un demi-cercle. Il est temps maintenant d'apprendre à \emph{utiliser} ces outils ensemble pour résoudre des problèmes complets, exactement comme au BEPC. Un exercice type te donne un cercle, quelques points, une mesure d'angle, et te demande de calculer d'autres angles \emph{en justifiant chaque étape}. Cette section t'apprend la stratégie complète, pas à pas.

\courssub{La stratégie générale : lire une figure avec un cercle}

Face à une figure comportant un cercle, le réflexe du bon élève n'est pas de se précipiter sur les calculs, mais d'\emph{observer} d'abord. Trois questions doivent guider ton regard :

\begin{enumerate}
\item \textbf{Où est le centre ?} Le point $O$ (souvent nommé ainsi) permet d'identifier les angles au centre et les rayons. Rappelle-toi : deux rayons forment toujours un triangle isocèle !
\item \textbf{Quels angles sont inscrits ?} Un angle inscrit a son sommet \emph{sur} le cercle et ses côtés coupent le cercle en deux autres points.
\item \textbf{Quels arcs sont interceptés ?} Chaque angle inscrit « regarde » un arc (celui qui ne contient pas son sommet). Deux angles qui regardent le même arc sont égaux ; un angle inscrit vaut la moitié de l'angle au centre qui regarde le même arc.
\end{enumerate}

\begin{methode}[Pour calculer un angle dans un cercle]
\begin{enumerate}
\item \textbf{Repérer l'arc intercepté} par l'angle cherché (ou par un angle connu de la figure).
\item \textbf{Chercher l'angle au centre ou un angle inscrit associé} à ce même arc.
\item \textbf{Appliquer le théorème} : angle au centre $= 2 \times$ angle inscrit, ou angles inscrits interceptant le même arc sont égaux.
\item \textbf{Conclure en justifiant} : écrire la propriété utilisée et vérifier que ses hypothèses sont remplies (sommet sur le cercle, même arc).
\end{enumerate}
Si l'angle cherché n'est pas directement lié à un arc, utilise les propriétés des triangles (somme des angles égale à $180°$, triangle isocèle) ou du parallélisme (angles alternes-internes, correspondants) comme \emph{relais} entre les étapes.
\end{methode}

\courssub{Combiner avec les propriétés des triangles}

Le théorème de l'angle inscrit ne suffit pas toujours seul. Heureusement, tu connais déjà deux outils puissants :

\begin{itemize}
\item la \cle{somme des angles d'un triangle} vaut $180°$ ;
\item un triangle \cle{isocèle} a deux angles à la base égaux. Or, dès que deux côtés d'un triangle sont des \emph{rayons} du même cercle, le triangle est isocèle !
\end{itemize}

\begin{exemplebox}[Triangle formé par deux rayons]
Soit un cercle de centre $O$, et $A$, $B$ deux points du cercle tels que $\widehat{AOB} = 70°$. Le triangle $OAB$ est isocèle en $O$ car $OA = OB$ (rayons). Donc :
\[ \widehat{OAB} = \widehat{OBA} = \frac{180° - 70°}{2} = 55°. \]
De plus, pour tout point $M$ du cercle situé sur l'arc majeur $AB$ (celui qui ne contient pas $O$), l'angle inscrit $\widehat{AMB}$ intercepte le même arc $AB$ que l'angle au centre $\widehat{AOB}$, donc :
\[ \widehat{AMB} = \frac{1}{2}\widehat{AOB} = \frac{70°}{2} = 35°. \]
\end{exemplebox}

\begin{exemplebox}[Enchaînement avec un triangle isocèle]
Toujours dans le cercle de centre $O$ avec $\widehat{AOB} = 70°$, soit $M$ un point du cercle tel que $\widehat{AMB} = 35°$. Supposons que le triangle $AMN$ soit isocèle en $M$ avec $\widehat{AMN} = 35°$. Alors les angles à la base sont égaux : $\widehat{MAN} = \widehat{MNA}$. Comme la somme des angles du triangle vaut $180°$ :
\[ \widehat{MAN} = \frac{180° - \widehat{AMN}}{2} = \frac{180° - 35°}{2} = 72{,}5°. \]
\end{exemplebox}

\begin{attention}[Le sommet doit être SUR le cercle]
L'erreur la plus fréquente au BEPC est d'appliquer le théorème de l'angle inscrit à un angle dont le sommet n'est \textbf{pas} sur le cercle. Par exemple, si deux cordes $[AB]$ et $[CD]$ se coupent en un point $I$ \emph{intérieur} au cercle, l'angle $\widehat{AIC}$ n'est \textbf{pas} un angle inscrit : on ne peut pas écrire $\widehat{AIC} = \frac{1}{2}\widehat{AOC}$. Avant d'appliquer le théorème, vérifie toujours : le sommet de l'angle est-il un point du cercle ? Ses deux côtés recoupent-ils le cercle en deux autres points distincts du sommet ?
\end{attention}

\courssub{Combiner avec le parallélisme}

Quand la figure comporte des droites parallèles (souvent une corde parallèle à une autre droite, ou une tangente parallèle à une corde), les angles \cle{alternes-internes} et \cle{correspondants} deviennent des relais précieux : ils permettent de « transporter » une mesure d'angle d'un endroit de la figure à un autre, là où le théorème de l'angle inscrit pourra s'appliquer.

\begin{exemplebox}[Parallélisme et angle inscrit]
Soit un cercle de centre $O$, $A$, $B$, $C$ trois points du cercle, et une droite $(d)$ passant par $C$ et parallèle à $(AB)$. On donne $\widehat{BAC} = 40°$. Comme $(d) \parallel (AB)$, les angles alternes-internes formés par la sécante $(AC)$ sont égaux : l'angle entre $(d)$ et $(CA)$ en $C$ vaut $40°$. D'autre part, l'angle inscrit $\widehat{BAC}$ intercepte l'arc $BC$, donc l'angle au centre associé vaut :
\[ \widehat{BOC} = 2 \times \widehat{BAC} = 2 \times 40° = 80°. \]
\end{exemplebox}

\courssub{Rédiger une démonstration : la rigueur avant tout}

Au BEPC, un calcul d'angle non justifié ne rapporte qu'une partie des points. Chaque étape doit citer \emph{précisément} la propriété utilisée et vérifier ses hypothèses.

\begin{methode}[Modèle de rédaction]
Pour montrer que $\widehat{AOB} = 2 \times \widehat{AMB}$, écris :

« Les points $A$, $M$ et $B$ appartiennent au cercle de centre $O$. L'angle $\widehat{AMB}$ est un angle \textbf{inscrit} et l'angle $\widehat{AOB}$ est l'angle \textbf{au centre} associé : ils interceptent le \textbf{même arc} $AB$ (celui ne contenant pas $M$). D'après le théorème de l'angle inscrit, $\widehat{AOB} = 2 \times \widehat{AMB}$. »

De même, pour deux angles inscrits : « Les angles $\widehat{AMB}$ et $\widehat{ANB}$ sont inscrits dans le même cercle et interceptent le même arc $AB$, donc $\widehat{AMB} = \widehat{ANB}$. »
\end{methode}

\begin{attention}[Hypothèses à vérifier avant de citer le théorème]
\begin{itemize}
\item Pour « angle au centre $= 2 \times$ angle inscrit » : les deux angles doivent intercepter le \emph{même} arc. L'angle au centre associé à un angle inscrit est celui qui intercepte l'arc \textbf{ne contenant pas le sommet} de l'angle inscrit. Ne pas confondre avec l'autre angle au centre (l'angle rentrant, $> 180°$).
\item Pour « angles inscrits égaux » : les deux sommets doivent être du \emph{même côté} de la corde $AB$ (ils interceptent alors le même arc). Si les sommets sont de part et d'autre de la corde, les angles sont \textbf{supplémentaires} (leur somme vaut $180°$), pas égaux.
\item Pour « triangle rectangle inscrit » : il faut que l'un des côtés soit un \emph{diamètre}, pas seulement une corde quelconque.
\end{itemize}
\end{attention}

\courssub{Exemple guidé : calculs d'angles en chaîne}

Voici une figure complexe, typique de l'épreuve du BEPC, avec un diamètre, deux cordes et plusieurs angles à déterminer de proche en proche.

\begin{popfigure}
\begin{center}
\begin{tikzpicture}[scale=1.8]
\coordinate (O) at (0,0);
\draw[popblue, thick] (O) circle (2);
% Diamètre AD horizontal
\coordinate (A) at (-2,0); % 180°
\coordinate (D) at (2,0);  % 0°
% Points B, C sur l'arc supérieur (ordre A, B, C, D)
\coordinate (B) at (130:2);
\coordinate (C) at (70:2);
% Diamètre
\draw[popdark, thick] (A) -- (D);
\fill[popink] (O) circle (0.035) node[below=2pt] {$O$};
% Cordes
\draw[poporange, thick] (A) -- (B);
\draw[poporange, thick] (B) -- (C);
\draw[popgreen, thick] (C) -- (D);
\draw[popgreen, thick] (A) -- (C);
% Points
\fill[popink] (A) circle (0.04) node[left] {$A$};
\fill[popink] (B) circle (0.04) node[above left] {$B$};
\fill[popink] (C) circle (0.04) node[above right] {$C$};
\fill[popink] (D) circle (0.04) node[right] {$D$};
% Angles donnés en A
\draw[popblue!60, thick] (A) ++(10:0.6) arc (10:50:0.6);
\node[popblue!60] at ($(A)+(30:0.95)$) {\small $30°$};
\draw[popblue!60, thick] (A) ++(-50:0.6) arc (-50:-10:0.6);
\node[popblue!60] at ($(A)+(-30:0.95)$) {\small $25°$};
% Angles cherchés (codés ?)
\draw[poppurple, thick] (B) ++(190:0.5) arc (190:250:0.5);
\node[poppurple] at ($(B)+(220:0.85)$) {\small $?$};
\draw[poppurple, thick] (C) ++(-30:0.5) arc (-30:30:0.5);
\node[poppurple] at ($(C)+(0:0.85)$) {\small $?$};
% Angle droit en C (indicatif)
\draw[popgreen!50!black, thick] ($(C)+(220:0.3)$) -- ($(C)+(220:0.3)+(310:0.3)$) -- ($(C)+(310:0.3)$);
\end{tikzpicture}
\end{center}
\legende{Figure type examen : $[AD]$ est un diamètre, $(AB)$, $(BC)$, $(CD)$, $(AC)$ sont des cordes. On donne $\widehat{BAC} = 30°$ et $\widehat{CAD} = 25°$ ; on cherche $\widehat{ABC}$ et $\widehat{ACD}$.}
\end{popfigure}

\begin{exoresolu}[Calculs en chaîne sur la figure type examen]
Sur la figure ci-dessus, $[AD]$ est un diamètre du cercle de centre $O$. On donne $\widehat{BAC} = 30°$ et $\widehat{CAD} = 25°$. Calculer $\widehat{ACD}$ puis $\widehat{ABC}$, en justifiant chaque étape.
\tcblower
\textbf{Étape 1 : un triangle rectangle.} Le triangle $ACD$ est inscrit dans le cercle et son côté $[AD]$ est un diamètre. D'après le théorème du triangle inscrit dans un demi-cercle, $ACD$ est rectangle en $C$ : $\widehat{ACD} = 90°$.

\textbf{Étape 2 : angle dans le triangle $ACD$.} La somme des angles du triangle $ACD$ vaut $180°$ :
\[ \widehat{ADC} = 180° - \widehat{ACD} - \widehat{CAD} = 180° - 90° - 25° = 65°. \]

\textbf{Étape 3 : angles inscrits interceptant le même arc.} Les angles $\widehat{ABC}$ et $\widehat{ADC}$ sont inscrits dans le même cercle et interceptent le même arc $AC$ (leurs sommets $B$ et $D$ sont du même côté de la corde $[AC]$). Donc :
\[ \widehat{ABC} = \widehat{ADC} = 65°. \]

\textbf{Vérification} : dans le triangle $ABC$, on a $\widehat{BAC} = 30°$, $\widehat{ABC} = 65°$, d'où $\widehat{ACB} = 180° - 30° - 65° = 85°$, ce qui est cohérent avec la figure ($\widehat{ACD} = \widehat{ACB} + \widehat{BCD} = 90°$).
\end{exoresolu}

\courssub{Exemples de démonstrations classiques}

Au BEPC, on ne te demande pas seulement de calculer : on te demande aussi de \emph{démontrer}. Voici trois modèles à connaître par cœur.

\begin{exoresolu}[Démontrer une égalité d'angles]
Soit un cercle de centre $O$ et quatre points $A$, $B$, $M$, $N$ du cercle, $M$ et $N$ étant du même côté de la corde $[AB]$. Démontrer que $\widehat{AMB} = \widehat{ANB}$.
\tcblower
Les angles $\widehat{AMB}$ et $\widehat{ANB}$ sont deux angles \textbf{inscrits} dans le même cercle (leurs sommets $M$ et $N$ sont sur le cercle). Ils interceptent le \textbf{même arc} $AB$, car $M$ et $N$ sont du même côté de la corde $[AB]$. Or, deux angles inscrits qui interceptent le même arc sont égaux. Donc $\widehat{AMB} = \widehat{ANB}$. \hfill$\blacksquare$
\end{exoresolu}

\begin{exoresolu}[Démontrer que deux droites sont perpendiculaires]
Soit un cercle de centre $O$, $[AB]$ un diamètre et $M$ un point du cercle distinct de $A$ et $B$. Démontrer que $(MA) \perp (MB)$.
\tcblower
Le triangle $AMB$ est inscrit dans le cercle et son côté $[AB]$ est un \textbf{diamètre}. D'après le théorème du triangle inscrit dans un demi-cercle, le triangle $AMB$ est rectangle en $M$, c'est-à-dire $\widehat{AMB} = 90°$. Par conséquent, les droites $(MA)$ et $(MB)$ sont perpendiculaires. \hfill$\blacksquare$

\emph{Variante par le calcul} : l'angle inscrit $\widehat{AMB}$ intercepte l'arc $AB$ ne contenant pas $M$, dont l'angle au centre associé est $\widehat{AOB} = 180°$ (angle plat, car $[AB]$ est un diamètre). Donc $\widehat{AMB} = \dfrac{180°}{2} = 90°$.
\end{exoresolu}

\begin{exoresolu}[Démontrer que quatre points sont cocycliques -- Réciproque]
Soit un triangle $ABC$ rectangle en $A$. On note $O$ le milieu de l'hypoténuse $[BC]$. Démontrer que les points $A$, $B$ et $C$ appartiennent à un même cercle de centre $O$, puis que $OA = OB = OC$.
\tcblower
Le triangle $ABC$ est rectangle en $A$ et $[BC]$ est son hypoténuse. Or, la propriété réciproque du triangle inscrit dans un demi-cercle affirme : si un triangle est rectangle, alors son hypoténuse est un diamètre de son cercle circonscrit. Donc $A$, $B$, $C$ appartiennent au cercle de diamètre $[BC]$, dont le centre est le milieu $O$ de $[BC]$. Les segments $OA$, $OB$, $OC$ sont des rayons de ce cercle, donc $OA = OB = OC$. \hfill$\blacksquare$

C'est d'ailleurs ainsi que l'on démontre la célèbre propriété : \emph{dans un triangle rectangle, le milieu de l'hypoténuse est équidistant des trois sommets}.
\end{exoresolu}

\begin{popfigure}
\begin{center}
\begin{tikzpicture}[scale=1.5]
\coordinate (O) at (0,0);
\draw[popblue, thick] (O) circle (2);
\coordinate (A) at (210:2);
\coordinate (B) at (30:2);
\coordinate (M) at (110:2);
% Diamètre AB
\draw[popdark, thick] (A) -- (B);
\fill[popink] (O) circle (0.035) node[below=2pt] {$O$};
% Triangle inscrit
\draw[poporange, thick] (M) -- (A);
\draw[poporange, thick] (M) -- (B);
% Codage angle droit en M
\draw[popgreen, thick] ($(M)+(250:0.3)$) -- ($(M)+(250:0.3)+(340:0.3)$) -- ($(M)+(340:0.3)$);
% Codage rayons égaux
\draw[popmuted, dashed] (O) -- (M);
\fill[popink] (A) circle (0.04) node[below left] {$A$};
\fill[popink] (B) circle (0.04) node[above right] {$B$};
\fill[popink] (M) circle (0.04) node[above] {$M$};
\node[popgreen] at ($(M)+(-90:0.55)$) {\small $90°$};
\end{tikzpicture}
\end{center}
\legende{Triangle $AMB$ inscrit, $[AB]$ diamètre : l'angle en $M$ est droit. À retrouver pas à pas : $\widehat{AOB} = 180°$ (angle plat), donc $\widehat{AMB} = \frac{1}{2}\widehat{AOB} = 90°$.}
\end{popfigure}

\courssub{Exercice d'entraînement}

\begin{exercice}[Pour t'entraîner]
Soit un cercle de centre $O$. Les points $A$, $B$, $C$, $D$ sont placés sur le cercle dans cet ordre. On donne $\widehat{AOB} = 80°$ et $\widehat{BAC} = 50°$.
\begin{enumerate}
\item Calculer $\widehat{ADB}$ en justifiant.
\item Calculer $\widehat{BOC}$ en justifiant.
\item Le triangle $OAB$ est-il isocèle ? Calculer $\widehat{OAB}$.
\end{enumerate}
\end{exercice}

\begin{corrige}[Exercice d'entraînement]
\begin{enumerate}
\item L'angle $\widehat{ADB}$ est inscrit et intercepte l'arc $AB$, comme l'angle au centre $\widehat{AOB}$. D'après le théorème de l'angle inscrit : $\widehat{ADB} = \dfrac{1}{2}\widehat{AOB} = \dfrac{80°}{2} = 40°$.
\item L'angle inscrit $\widehat{BAC}$ intercepte l'arc $BC$, dont l'angle au centre associé est $\widehat{BOC}$. Donc $\widehat{BOC} = 2 \times \widehat{BAC} = 2 \times 50° = 100°$.
\item $OA = OB$ (rayons du même cercle), donc $OAB$ est isocèle en $O$. Ainsi $\widehat{OAB} = \widehat{OBA} = \dfrac{180° - 80°}{2} = 50°$.
\end{enumerate}
\end{corrige}

\begin{aretenir}[L'essentiel de la section]
\begin{itemize}
\item \textbf{Stratégie} : repérer le centre, les arcs, les angles inscrits ; puis appliquer le théorème en vérifiant les hypothèses.
\item \textbf{Relais utiles} : somme des angles d'un triangle ($180°$), triangles isocèles formés par deux rayons, angles alternes-internes ou correspondants en cas de parallélisme.
\item \textbf{Rédaction type} : « Les angles $\widehat{AMB}$ et $\widehat{AOB}$ sont respectivement inscrit et au centre, associés au même arc $AB$, donc $\widehat{AOB} = 2 \times \widehat{AMB}$. »
\item \textbf{Piège mortel} : ne jamais appliquer le théorème à un angle dont le sommet n'est pas sur le cercle.
\end{itemize}
\end{aretenir}

Tu maîtrises maintenant toute la boîte à outils des angles dans le cercle. Dans la dernière section, nous verrons comment ces notions s'appliquent à des situations concrètes de la vie courante : orientation, architecture, sport et même navigation !

\coursec{Applications à la vie courante}

Dans la section précédente, nous avons appris à calculer des angles et à rédiger des démonstrations rigoureuses grâce au théorème de l'angle inscrit et à la propriété du triangle inscrit dans un demi-cercle. Mais à quoi servent vraiment ces résultats en dehors de la salle de classe ? C'est toute la question de cette dernière section : tu vas découvrir que les angles inscrits sont partout autour de toi — sur un terrain de football, sur le fleuve Wouri, sur un chantier de construction à Douala, dans l'atelier d'un artisan. Le but est de savoir \cle{modéliser} une situation réelle, c'est-à-dire la traduire en langage géométrique, appliquer les théorèmes, puis interpréter la réponse dans le contexte.

\courssub{La démarche de modélisation}

\begin{methode}[Modéliser un problème concret avec les angles inscrits]
Pour résoudre un problème de la vie courante faisant intervenir des angles de vision ou des positions sur un cercle, on procède en quatre étapes :
\begin{enumerate}
\item \textbf{Identifier} les objets assimilables à des \cle{points} et à un \cle{cercle} (le but devient un segment $[AB]$, les positions des spectateurs deviennent des points d'un arc de cercle, etc.) ;
\item \textbf{Traduire} la question en calcul d'angle : de quel angle parle-t-on ? Est-il inscrit ? Au centre ? Quel arc intercepte-t-il ?
\item \textbf{Appliquer} les théorèmes de la leçon : angle inscrit = moitié de l'angle au centre ; angles inscrits interceptant le même arc sont égaux ; triangle inscrit dans un demi-cercle est rectangle ;
\item \textbf{Rédiger} la conclusion dans le contexte, avec une phrase complète et l'unité (degrés) : « le supporter voit la scène sous un angle de $45°$ ».
\end{enumerate}
\end{methode}

\begin{attention}[Vérifier l'hypothèse du cercle]
Dans un problème concret, on ne peut appliquer les propriétés des angles inscrits \textbf{que si les points concernés sont bien sur le même cercle}. Avant d'écrire « les angles sont égaux car ils interceptent le même arc », il faut s'assurer (et écrire !) que les points $A$, $B$, $M$, $N$ sont \cle{cocycliques}, c'est-à-dire situés sur un même cercle. Si les spectateurs d'un stade sont alignés sur une tribune \emph{droite}, leurs angles de vision ne sont \textbf{pas} égaux : le théorème ne s'applique pas. C'est l'erreur la plus fréquente dans ce type de problème.
\end{attention}

\courssub{Angles de vision : le stade et la scène}

Imaginons un concert au stade omnisports de Yaoundé. La scène est assimilée à un segment $[AB]$. Les spectateurs les mieux placés sont répartis sur des tribunes en arc de cercle passant par $A$ et $B$. Que constate-t-on ?

\begin{propriete}[Angle de vision constant sur un arc]
Tous les points $M$ situés sur un même arc de cercle passant par $A$ et $B$ voient le segment $[AB]$ sous le \textbf{même angle} : les angles inscrits $\widehat{AMB}$ qui interceptent le même arc $\overset{\frown}{AB}$ sont égaux.
\end{propriete}

Autrement dit, deux supporters situés à des endroits différents d'une même tribune circulaire ont exactement la \textbf{même qualité de vision} de la scène. C'est pourquoi les architectes dessinent les tribunes des stades et des théâtres en arcs de cercle : cela garantit une équité de vision entre les spectateurs d'un même rang.

\begin{exemplebox}[Projecteurs et supporter]
Deux projecteurs placés de part et d'autre du centre $O$ du stade éclairent la ligne de but $[AB]$. L'angle au centre $\widehat{AOB}$ vaut $90°$. Un supporter $M$ est assis sur la tribune circulaire de centre $O$, sur l'arc opposé. Sous quel angle voit-il la ligne de but ?

L'angle $\widehat{AMB}$ est un angle \cle{inscrit} et $\widehat{AOB}$ est l'angle \cle{au centre} associé : ils interceptent le même arc $\overset{\frown}{AB}$. D'après le théorème de l'angle inscrit :
\[ \widehat{AMB} = \frac{1}{2}\,\widehat{AOB} = \frac{1}{2}\times 90° = 45°. \]
Le supporter voit donc la ligne de but sous un angle de $45°$.
\end{exemplebox}

\courssub{Le problème du tir au but au football}

Voici un problème classique que se posent les joueurs du Canon de Yaoundé ou des Tonnerre de Kalara : un attaquant $M$ s'approche du but $[AB]$. Sous quel angle voit-il le but ? Plus cet angle $\widehat{AMB}$ est grand, plus le tir a de chances d'être cadré.

\begin{popfigure}
\begin{center}
\begin{tikzpicture}[scale=0.85]
% Terrain
\fill[popgreenL] (-6,-1) rectangle (6,6.5);
\draw[white,thick] (-6,-1) rectangle (6,6.5);
\draw[white,thick] (0,-1) -- (0,6.5);
\draw[white,thick] (0,2.75) circle (1.2);
% But AB
\draw[line width=2.5pt,popdark] (-1.5,6.5) -- (1.5,6.5);
\fill[popdark] (-1.5,6.5) circle (0.07) node[above left,black]{$A$};
\fill[popdark] (1.5,6.5) circle (0.07) node[above right,black]{$B$};
% Centre O
\coordinate (O) at (0,5);
\fill[popblue] (O) circle (0.07) node[below,popblue]{$O$};
% Arc de cercle passant par A et B (arc majeur, centre O(0,5), R=sqrt(4.5)=2.121)
\draw[popblue,very thick] (-1.5,6.5) arc (135:405:2.121);
% Points M, N, P sur l'arc (angles 200, 270, 340)
\coordinate (M) at (-1.993,4.275); % 200 deg
\coordinate (N) at (0,2.879);     % 270 deg
\coordinate (P) at (1.993,4.275);  % 340 deg
\fill[poporange] (M) circle (0.09) node[left,popdark]{$M$};
\fill[poporange] (N) circle (0.09) node[below,popdark]{$N$};
\fill[poporange] (P) circle (0.09) node[right,popdark]{$P$};
\draw[poporange] (M) -- (-1.5,6.5) (M) -- (1.5,6.5);
\draw[poporange] (N) -- (-1.5,6.5) (N) -- (1.5,6.5);
\draw[poporange,dashed] (P) -- (-1.5,6.5) (P) -- (1.5,6.5);
% Rayon OA, OB
\draw[popblue,thin,dashed] (O) -- (-1.5,6.5) (O) -- (1.5,6.5);
\end{tikzpicture}
\end{center}
\legende{Les tireurs $M$, $N$ et $P$, placés sur le même arc de cercle passant par les poteaux $A$ et $B$, voient le but sous le même angle : $\widehat{AMB} = \widehat{ANB} = \widehat{APB}$.}
\end{popfigure}

\begin{exoresolu}[Où frapper pour avoir le meilleur angle ?]
Un attaquant peut frapper depuis trois positions $M$, $N$ et $P$ situées sur un même arc de cercle passant par les deux poteaux $A$ et $B$. L'angle au centre associé à l'arc $\overset{\frown}{AB}$ vaut $70°$.
\begin{enumerate}
\item Comparer les angles de tir $\widehat{AMB}$, $\widehat{ANB}$ et $\widehat{APB}$.
\item Calculer leur valeur commune.
\item Un autre joueur $Q$, situé \emph{plus près} du but sur un arc de cercle plus petit, voit le but sous un angle au centre de $110°$. Qui a le meilleur angle de tir ?
\end{enumerate}
\tcblower
\textbf{Solution.}
\begin{enumerate}
\item Les points $A$, $B$, $M$, $N$, $P$ sont sur le même cercle. Les angles $\widehat{AMB}$, $\widehat{ANB}$ et $\widehat{APB}$ sont des angles inscrits qui interceptent le même arc $\overset{\frown}{AB}$ : ils sont donc \textbf{égaux}.
\item D'après le théorème de l'angle inscrit :
\[ \widehat{AMB} = \widehat{ANB} = \widehat{APB} = \frac{1}{2}\times 70° = 35°. \]
\item Le joueur $Q$ voit le but sous l'angle $\frac{1}{2}\times 110° = 55°$. Comme $55° > 35°$, c'est $Q$ qui a le meilleur angle de tir : plus on est proche du but, plus l'arc est petit, plus l'angle de tir est grand.
\end{enumerate}
\textbf{Conclusion dans le contexte :} sur un même arc de cercle, toutes les positions de tir sont équivalentes ; pour améliorer son angle, l'attaquant doit se rapprocher du but (passer sur un arc plus petit).
\end{exoresolu}

\courssub{Orientation et navigation sur le fleuve Wouri}

Un piroguier navigue sur le fleuve Wouri, près de Douala. Pour connaître sa position, il repère deux points fixes sur la rive : un phare $A$ et une antenne $B$. Il mesure l'angle $\widehat{AMB}$ sous lequel il voit ces deux repères depuis son embarcation $M$.

\begin{popfigure}
\begin{center}
\begin{tikzpicture}[scale=0.9]
% Rive
\fill[popgoldL] (-6,2.5) rectangle (6,5);
\draw[popgold,thick] (-6,2.5) -- (6,2.5);
\node[popdark] at (4.5,4.4) {\textbf{Rive du Wouri}};
% Fleuve
\fill[popblueL] (-6,-2.5) rectangle (6,2.5);
\node[popblue!70!black] at (-4.5,-1.8) {\textit{Fleuve Wouri}};
% Repères A, B
\coordinate (A) at (-2.5,2.5);
\coordinate (B) at (2.5,2.5);
\fill[popdark] (A) circle (0.08) node[above]{$A$ (phare)};
\fill[popdark] (B) circle (0.08) node[above]{$B$ (antenne)};
% Cercle : Centre O(0, -1.83), R=5. Arc majeur (côté fleuve) angles 120 à 420.
\coordinate (O) at (0,-1.83);
\fill[popink] (O) circle (0.07) node[left,popink]{$O$};
\draw[popink,very thick] (A) arc (120:420:5);
% Points M, N, K sur l'arc dans la zone visible (y > -2.5)
% M: angle 140 deg -> (-3.83, 1.38)
% N: angle 180 deg -> (-5.00, -1.83)
% K: angle 400 deg (40 deg) -> (3.83, 1.38)
\coordinate (M) at (-3.83,1.38);
\coordinate (N) at (-5.00,-1.83);
\coordinate (K) at (3.83,1.38);
\fill[poporange] (M) circle (0.1) node[left,popdark]{$M$};
\fill[poporange] (N) circle (0.1) node[below,popdark]{$N$};
\fill[poporange] (K) circle (0.1) node[right,popdark]{$K$};
\draw[poporange] (M) -- (A) (M) -- (B);
\draw[poporange,dashed] (N) -- (A) (N) -- (B);
\draw[poporange,dashed] (K) -- (A) (K) -- (B);
\end{tikzpicture}
\end{center}
\legende{Depuis n'importe quelle position $M$, $N$ ou $K$ sur le même arc de cercle, le navigateur voit les repères $A$ et $B$ sous le même angle : c'est le principe du relevé de position.}
\end{popfigure}

\begin{propriete}[Lieu des points de même angle de visée]
L'ensemble des points $M$ d'un même côté de la droite $(AB)$ d'où l'on voit le segment $[AB]$ sous un angle donné $\alpha$ est un \cle{arc de cercle} passant par $A$ et $B$, appelé \textbf{arc capable}.
\end{propriete}

C'est exactement la réciproque de la propriété des angles inscrits : si $\widehat{AMB} = \widehat{ANB}$, alors $M$ et $N$ sont sur le même arc de cercle passant par $A$ et $B$. Le navigateur qui mesure un angle de $60°$ sait donc qu'il se trouve quelque part sur l'arc capable correspondant ; en mesurant un deuxième angle avec un troisième repère, il détermine sa position exacte à l'intersection des deux arcs.

\begin{savaistu}[Le sextant, ancêtre du GPS]
Bien avant les satellites et le GPS, les navigateurs utilisaient le \cle{sextant}, un instrument de mesure d'angles inventé au XVIII\textsuperscript{e} siècle. En visant un astre (le Soleil, l'étoile Polaire) et l'horizon, le marin mesurait un angle, puis reportait cet angle sur des cercles tracés sur sa carte — exactement le principe des arcs capables et des angles inscrits que tu viens d'étudier ! C'est grâce à cette géométrie que les grands voyages maritimes ont été possibles. Aujourd'hui encore, les officiers de marine marchande apprennent à s'en servir : en cas de panne électronique au large des côtes du Cameroun, le sextant reste un instrument de secours fiable.
\end{savaistu}

\courssub{Construction : tracer un angle droit avec une corde}

Sur un chantier de construction à Bonabéri, un maçon doit vérifier qu'un mur est bien perpendiculaire à un autre. Il n'a pas d'équerre professionnelle, mais il a une corde. Comment faire ? Il utilise la propriété du triangle inscrit dans un demi-cercle !

\begin{methode}[Tracer un angle droit avec une corde]
\begin{enumerate}
\item Tendre la corde entre deux piquets $A$ et $B$ : $[AB]$ sera le \cle{diamètre}.
\item Trouver le milieu $O$ de la corde en la pliant en deux.
\item Fixer un troisième piquet $M$ en gardant la corde tendue de $O$ vers $M$, avec $OM = OA$ : le point $M$ est alors sur le cercle de diamètre $[AB]$.
\item L'angle $\widehat{AMB}$ est alors un angle \textbf{droit}, car tout triangle inscrit dans un demi-cercle est rectangle.
\end{enumerate}
\end{methode}

\begin{exemplebox}[Vérification sur chantier]
Le maçon place $A$ et $B$ à \SI{4}{\meter} l'un de l'autre, donc $AB = \SI{4}{\meter}$. Le milieu $O$ est à \SI{2}{\meter} de chaque extrémité. Il tend la corde de $O$ vers $M$ avec $OM = \SI{2}{\meter}$ : $M$ appartient au cercle de diamètre $[AB]$. Le triangle $AMB$ est inscrit dans le demi-cercle de diamètre $[AB]$, donc il est rectangle en $M$ :
\[ \widehat{AMB} = 90°. \]
Le mur tracé le long de $(MA)$ est donc bien perpendiculaire au mur tracé le long de $(MB)$.
\end{exemplebox}

\begin{attention}[La condition du diamètre]
La propriété ne fonctionne que si $[AB]$ est bien un \cle{diamètre} du cercle, c'est-à-dire si $M$ est à égale distance de $O$ que $A$ et $B$. Si la corde n'est pas tendue correctement et que $OM \neq OA$, le point $M$ n'est pas sur le cercle et l'angle $\widehat{AMB}$ n'est pas droit. Sur le chantier, une erreur de quelques centimètres sur la longueur de la corde fausse l'équerre entière !
\end{attention}

\courssub{Artisanat : rosaces et motifs circulaires}

Les artisans camerounais — sculpteurs de l'Ouest, brodeuses du Nord, vanniers de la côte — créent de magnifiques motifs circulaires : rosaces sur les portes sculptées, motifs des chapeaux tissés, décorations des calebasses. Pour diviser un cercle en parts égales, l'artisan utilise (souvent sans le nommer) les angles au centre et les angles inscrits.

\begin{exemplebox}[Une rosace à six branches]
Un sculpteur veut graver une rosace à $6$ branches identiques sur une porte. Il partage le tour complet en $6$ angles au centre égaux :
\[ \frac{360°}{6} = 60°. \]
Chaque arc entre deux branches consécutives mesure donc $60°$. S'il veut ensuite graver une étoile dont chaque pointe $M$ regarde deux branches voisines $A$ et $B$, l'angle inscrit de la pointe vaut :
\[ \widehat{AMB} = \frac{1}{2}\times 60° = 30°. \]
Toutes les pointes de l'étoile ont le même angle, ce qui garantit la régularité et la beauté du motif.
\end{exemplebox}

\courssub{Synthèse : un problème complet de modélisation}

Terminons par un problème qui mobilise toute la démarche, comme on t'en demandera au BEPC.

\begin{exoresolu}[Le feu d'artifice de la fête nationale]
Pour la fête du 20 mai, deux artificiers sont postés en $A$ et $B$, distants de \SI{100}{\meter}, de part et d'autre de l'esplanade. Le public est placé derrière une barrière de sécurité en arc de cercle passant par $A$ et $B$. Le centre $O$ de ce cercle est le poste de commandement, et l'angle $\widehat{AOB}$ vaut $120°$.
\begin{enumerate}
\item Schématiser la situation (identifier points et cercle).
\item Calculer l'angle $\widehat{AMB}$ sous lequel un spectateur $M$ de la barrière voit les deux postes de tir.
\item Un organisateur propose de placer des invités en $I$, à l'intérieur du disque, sur un arc plus petit où l'angle au centre vaut $150°$. Les invités verront-ils mieux le spectacle ? Justifier.
\end{enumerate}
\tcblower
\textbf{Solution.}
\begin{enumerate}
\item Les postes de tir sont les points $A$ et $B$ ; la barrière de sécurité est un arc de cercle de centre $O$ passant par $A$ et $B$ ; chaque spectateur est un point $M$ de cet arc.
\item Les points $A$, $B$ et $M$ sont sur le même cercle de centre $O$. L'angle inscrit $\widehat{AMB}$ et l'angle au centre $\widehat{AOB}$ interceptent le même arc $\overset{\frown}{AB}$. D'après le théorème de l'angle inscrit :
\[ \widehat{AMB} = \frac{1}{2}\,\widehat{AOB} = \frac{1}{2}\times 120° = 60°. \]
Chaque spectateur de la barrière voit donc les deux postes de tir sous un angle de $60°$, quelle que soit sa place sur la barrière.
\item Pour les invités en $I$ : $\widehat{AIB} = \dfrac{1}{2}\times 150° = 75°$. Comme $75° > 60°$, les invités voient le spectacle sous un angle plus grand : ils ont une meilleure vue d'ensemble des deux postes de tir.
\end{enumerate}
\textbf{Conclusion :} la barrière circulaire garantit à tous les spectateurs le même angle de vision de $60°$ ; rapprocher les invités sur un arc plus petit améliore leur angle de vision, qui passe à $75°$.
\end{exoresolu}

\begin{aretenir}[Ce qu'il faut retenir de cette section]
\begin{itemize}
\item Tous les points d'un même arc de cercle voient un segment $[AB]$ sous le \textbf{même angle} : c'est le principe des tribunes de stade, des positions de tir au football et des relevés de navigation.
\item Plus on est \textbf{proche} du segment (arc plus petit), plus l'angle de vision est \textbf{grand}.
\item Tout triangle inscrit dans un demi-cercle est \textbf{rectangle} : c'est le principe de l'équerre de chantier faite avec une corde.
\item Pour modéliser : identifier les points et le cercle, traduire en calcul d'angle, appliquer les théorèmes, conclure dans le contexte — et toujours vérifier que les points sont bien sur le même cercle.
\end{itemize}
\end{aretenir}

\begin{exercice}[À toi de jouer]
Deux villages $A$ et $B$ sont situés sur les deux rives d'un lac circulaire de centre $O$. Un pêcheur en pirogue $M$ sur le lac voit les deux villages sous un angle $\widehat{AMB} = 40°$.
\begin{enumerate}
\item Calculer l'angle au centre $\widehat{AOB}$.
\item Un deuxième pêcheur $N$, sur le même arc que $M$, voit-il les villages sous le même angle ? Justifier.
\item Le diamètre du lac est de \SI{2}{\kilo\meter}. Si le pêcheur $M$ se place de sorte que $[AB]$ soit un diamètre, que vaut l'angle $\widehat{AMB}$ ? Justifier.
\end{enumerate}
\end{exercice}

\begin{corrige}[Exercice 1]
\begin{enumerate}
\item L'angle au centre est le double de l'angle inscrit interceptant le même arc : $\widehat{AOB} = 2\times 40° = 80°$.
\item Oui : $M$ et $N$ sont sur le même cercle, et les angles inscrits $\widehat{AMB}$ et $\widehat{ANB}$ interceptent le même arc $\overset{\frown}{AB}$, donc $\widehat{ANB} = \widehat{AMB} = 40°$.
\item Si $[AB]$ est un diamètre, le triangle $AMB$ est inscrit dans un demi-cercle : il est rectangle en $M$, donc $\widehat{AMB} = 90°$.
\end{enumerate}
\end{corrige}

\newpage

\coursec{Activité d'intégration}

\coursec{Leçon 12 : Angles inscrits}

\courssub{Objectifs de la leçon}

À l'issue de cette leçon, tu seras capable de :

\begin{itemize}
\item définir un \cle{angle inscrit}, un \cle{angle au centre} et identifier l'\cle{arc intercepté} ;
\item énoncer et appliquer le \cle{théorème de l'angle inscrit et de l'angle au centre} : l'angle au centre est le double de l'angle inscrit interceptant le même arc ;
\item énoncer et appliquer la propriété des \cle{angles inscrits interceptant le même arc} : ils sont égaux ;
\item énoncer et appliquer la propriété du \cle{triangle inscrit dans un demi-cercle} : tout triangle inscrit dans un demi-cercle est rectangle (et sa réciproque) ;
\item modéliser une situation concrète (angle de tir, angle de vue) par une figure géométrique, calculer des angles et \cle{rédiger une justification complète} en citant les théorèmes utilisés, comme exigé au BEPC.
\end{itemize}

\courssub{Situation déclenchante : Le meilleur angle de tir au stade Ahmadou Ahidjo}

Lors d'un match des Lions Indomptables au stade Ahmadou Ahidjo de Yaoundé, un commentateur sportif affirme : « L'attaquant aurait dû se rapprocher du but, son \cle{angle de tir} était trop faible ! ». Mais qu'est-ce que l'angle de tir, et comment le rendre le plus grand possible ?

On modélise la situation : les deux poteaux du but sont les points $A$ et $B$. La largeur réglementaire est $AB = 7,32$~m. Un attaquant en $M$ voit le but sous l'angle $\widehat{AMB}$ : c'est l'\textbf{angle de tir}. Plus cet angle est grand, plus la cible est large.

Un résultat géométrique fondamental nous guide : \textbf{tous les points $M$ situés sur un même arc de cercle passant par $A$ et $B$ voient le but sous le même angle}, car ce sont des angles inscrits interceptant le même arc $\overset{\frown}{AB}$. On appelle ces arcs des « lignes d'égal angle de tir ».

Pour comparer les positions, on étudie trois cercles \textbf{différents} passant par $A$ et $B$ (trois « lignes d'égal angle ») :

\begin{itemize}
\item \textbf{Cercle $\mathcal{C}_1$ (éloigné) :} centre $O_1$, angle au centre $\widehat{AO_1B} = 60°$. Un point $M_1$ sur l'arc $\overset{\frown}{AB}$ (ne contenant pas $O_1$).
\item \textbf{Cercle $\mathcal{C}_2$ (intermédiaire) :} centre $O_2$, angle inscrit $\widehat{AM_2B} = 45°$ pour un point $M_2$ sur l'arc.
\item \textbf{Cercle $\mathcal{C}_3$ (proche, diamètre $[AB]$) :} centre $O_3$ (milieu de $[AB]$), angle inscrit $\widehat{AM_3B} = 90°$ pour un point $M_3$ sur le demi-cercle.
\end{itemize}

\begin{center}
\begin{popfigure}
\begin{tikzpicture}[scale=0.9]
% Coordinates
\coordinate (A) at (-3.66, 0);
\coordinate (B) at (3.66, 0);
% Circle 1: Large (Center O1 high up), Central angle 60 -> Radius = AB / (2 sin 30) = 3.66 / 0.5 = 7.32. Center height = R cos 30 = 7.32 * 0.866 = 6.34
\coordinate (O1) at (0, 6.34);
\def\Rone{7.32}
% Circle 2: Medium (Central angle 90 -> Inscribed 45). Radius = AB / (2 sin 45) = 3.66 / 0.707 = 5.17. Center height = R cos 45 = 5.17 * 0.707 = 3.66
\coordinate (O2) at (0, 3.66);
\def\Rtwo{5.17}
% Circle 3: Diameter AB. Center O3(0,0), Radius 3.66.
\coordinate (O3) at (0, 0);
\def\Rthree{3.66}
% Draw Circles (arcs only top part)
\draw[popblue, thick] (O1) circle (\Rone);
\draw[poporange, thick] (O2) circle (\Rtwo);
\draw[poppurple, thick] (O3) circle (\Rthree);
% Segment AB (Goal)
\draw[popdark, very thick] (A) -- (B);
\fill[popdark] (A) circle (2pt) node[below left] {$A$};
\fill[popdark] (B) circle (2pt) node[below right] {$B$};
% Centers
\fill[popblue] (O1) circle (1.5pt) node[left] {$O_1$};
\fill[poporange] (O2) circle (1.5pt) node[left] {$O_2$};
\fill[poppurple] (O3) circle (1.5pt) node[below right] {$O_3$};
% Points M on arcs (approx 90 deg from horizontal for visibility)
% M1 on Circle 1: angle 90 deg from center -> (0, 6.34+7.32) = (0, 13.66) too high. Let's take angle 120 deg from center (30 deg from vertical).
% Actually easier: use polar from centers.
% M1 on C1: angle 90° from center (top) -> (0, 6.34+7.32) = (0, 13.66). 
% Let's place them at x=0 for symmetry (top of arcs).
\coordinate (M1) at (0, 6.34+\Rone); % Top of C1
\coordinate (M2) at (0, 3.66+\Rtwo); % Top of C2
\coordinate (M3) at (0, \Rthree);    % Top of C3 (0, 3.66)
% Lines AM, BM
\draw[popblue, dashed, thin] (M1) -- (A) (M1) -- (B);
\draw[poporange, dashed, thin] (M2) -- (A) (M2) -- (B);
\draw[poppurple, dashed, thin] (M3) -- (A) (M3) -- (B);
% Center lines for C1 (60 deg central)
\draw[popblue, thin] (O1) -- (A) (O1) -- (B);
% Center lines for C2 (90 deg central)
\draw[poporange, thin] (O2) -- (A) (O2) -- (B);
% Labels M
\fill[popblue] (M1) circle (2pt) node[above] {$M_1$};
\fill[poporange] (M2) circle (2pt) node[above] {$M_2$};
\fill[poppurple] (M3) circle (2pt) node[above] {$M_3$};
% Angle marks (optional, keep clean)
% \pic [draw, popblue, angle radius=1.5cm] {angle = A--O1--B}; 
% \pic [draw, poporange, angle radius=1cm] {angle = A--O2--B};
\end{tikzpicture}
\legende{Trois cercles passant par $A$ et $B$ (poteaux). Plus le cercle est petit (joueur proche du but), plus l'angle de tir $\widehat{AMB}$ est grand. Sur chaque cercle, tous les points de l'arc $\overset{\frown}{AB}$ offrent le même angle de tir.}
\end{popfigure}
\end{center}

\textbf{Problème :} En tant que préparateur physique, tu dois établir une fiche conseil. Détermine comment varie l'angle de tir selon la distance au but, vérifie la relation angle au centre / angle inscrit, et rédige une conclusion argumentée pour les attaquants.

\courssub{Activité 1 : Manipulation et découverte (Travail de groupe)}

\begin{enumerate}
\item Sur une feuille, trace le segment $[AB]$ à l'échelle $1/200$ ($AB = 3,66$~cm).
\item Trace la médiatrice de $[AB]$. Choisis un point $O$ sur cette médiatrice (au-dessus de $[AB]$) et trace le cercle de centre $O$ passant par $A$ et $B$.
\item Place trois points $M_1$, $M_2$, $M_3$ sur le \textbf{même arc} $\overset{\frown}{AB}$ (celui qui ne contient pas le point diamétralement opposé à $O$ par rapport à $AB$).
\item Mesure au rapporteur les angles $\widehat{AM_1B}$, $\widehat{AM_2B}$, $\widehat{AM_3B}$. \textbf{Que constates-tu ?} Quel théorème ce constat illustre-t-il ?
\item Trace les rayons $[OA]$ et $[OB]$. Mesure l'angle au centre $\widehat{AOB}$. Compare sa mesure au double de l'angle de tir mesuré à la question 4. \textbf{Quel théorème vérifies-tu ?}
\end{enumerate}

\begin{aretenir}[Bilan de l'activité 1]
\begin{itemize}
\item \textbf{Angle inscrit :} Angle dont le sommet est sur le cercle et dont les côtés coupent le cercle en deux autres points ($A$ et $B$). Il \cle{intercepte} l'arc $\overset{\frown}{AB}$ qui ne contient pas son sommet.
\item \textbf{Angle au centre :} Angle dont le sommet est le centre $O$ du cercle et dont les côtés passent par $A$ et $B$. Il intercepte le même arc $\overset{\frown}{AB}$.
\item \textbf{Théorème 1 (Angles inscrits même arc) :} Dans un même cercle (ou dans des cercles égaux), deux angles inscrits qui interceptent le \textbf{même arc} sont \textbf{égaux}.
\item \textbf{Théorème 2 (Angle inscrit / Angle au centre) :} Dans un cercle, un angle au centre mesure le \textbf{double} d'un angle inscrit qui intercepte le \textbf{même arc} :
\[ \widehat{AOB} = 2 \times \widehat{AMB}. \]
\end{itemize}
\end{aretenir}

\courssub{Activité 2 : Le triangle rectangle inscrit dans un demi-cercle}

\begin{enumerate}
\setcounter{enumi}{5}
\item Sur une nouvelle figure, trace un segment $[AB] = 3,66$~cm. Construis son milieu $O_3$. Trace le cercle de centre $O_3$ et de rayon $O_3A$ (c'est le cercle de \textbf{diamètre} $[AB]$).
\item Place un point $M_3$ sur ce cercle (différent de $A$ et $B$). Trace le triangle $AM_3B$. Mesure l'angle $\widehat{AM_3B}$.
\item Recommence avec un autre point $M'_3$ sur le même demi-cercle. \textbf{Que constates-tu ?}
\end{enumerate}

\begin{propriete}[Triangle inscrit dans un demi-cercle -- Théorème 3]
Si $[AB]$ est un \textbf{diamètre} d'un cercle et $M$ un point quelconque de ce cercle (différent de $A$ et $B$), alors le triangle $AMB$ est \textbf{rectangle en $M$} :
\[ \widehat{AMB} = 90°. \]
\textbf{Réciproque :} Si un triangle $AMB$ est rectangle en $M$, alors son cercle circonscrit a pour diamètre l'hypoténuse $[AB]$ (le centre est le milieu de $[AB]$).
\end{propriete}

\begin{attention}[Pièges à éviter]
\begin{itemize}
\item \textbf{Sommet :} Vérifie toujours la position du sommet : sur le cercle $\rightarrow$ angle inscrit ; au centre $\rightarrow$ angle au centre.
\item \textbf{Facteur 2 :} L'angle au \textbf{centre} est le double de l'angle \textbf{inscrit} ($\widehat{AOB} = 2\widehat{AMB}$), jamais l'inverse.
\item \textbf{Même arc :} L'égalité des angles inscrits ne vaut que pour des points sur le \textbf{même arc} (même côté de la corde $[AB]$). Sur l'arc opposé, l'angle est supplémentaire ($180° - \alpha$).
\item \textbf{Justification BEPC :} Toute affirmation doit être justifiée par le \textbf{nom du théorème} (ex : « D'après le théorème de l'angle inscrit et de l'angle au centre... »).
\end{itemize}
\end{attention}

\courssub{Méthode : Résoudre un problème d'angle de vue / angle de tir}

\begin{methode}[Calculer et comparer des angles inscrits]
\begin{enumerate}
\item \textbf{Modéliser :} Représente la situation par un segment $[AB]$ (largeur but, largeur écran, distance entre deux arbres...) et des points $M$ (positions de l'observateur).
\item \textbf{Identifier le cercle :} Les points $M$ offrant le même angle sont sur un arc de cercle passant par $A$ et $B$. Détermine si on te donne l'angle au centre ou l'angle inscrit.
\item \textbf{Appliquer le bon théorème :}
      \begin{itemize}
      \item Si angle au centre $\widehat{AOB}$ connu $\rightarrow$ Angle inscrit $= \frac{1}{2} \widehat{AOB}$.
      \item Si angle inscrit $\widehat{AMB}$ connu $\rightarrow$ Angle au centre $= 2 \times \widehat{AMB}$.
      \item Si angle inscrit $= 90°$ $\rightarrow$ $[AB]$ est un diamètre (Théorème 3).
      \end{itemize}
\item \textbf{Comparer :} Pour un segment $[AB]$ fixe, plus le cercle passant par $A, B, M$ est \textbf{petit} (son centre $O$ est proche de $[AB]$), plus l'angle inscrit $\widehat{AMB}$ est \textbf{grand}. L'angle maximal ($90°$) est atteint sur le cercle de diamètre $[AB]$.
\item \textbf{Rédiger la conclusion :} Phrase complète, vocabulaire précis, théorème cité.
\end{enumerate}
\end{methode}

\courssub{Exercices résolus : Application au stade Ahmadou Ahidjo}

\begin{exoresolu}[Calculs d'angles de tir]
\textbf{Énoncé.} On reprend les trois positions du début de la leçon.
\begin{enumerate}
\item Pour $M_1$ sur $\mathcal{C}_1$, l'angle au centre est $\widehat{AO_1B} = 60°$. Calculer l'angle de tir $\widehat{AM_1B}$.
\item Pour $M_2$ sur $\mathcal{C}_2$, l'angle de tir est $\widehat{AM_2B} = 45°$. Calculer l'angle au centre $\widehat{AO_2B}$.
\item Pour $M_3$ sur $\mathcal{C}_3$, l'angle de tir est $\widehat{AM_3B} = 90°$. Quelle est la nature du segment $[AB]$ pour ce cercle ? Justifier.
\item Classer les angles de tir par ordre croissant. En déduire la position la plus favorable pour marquer.
\end{enumerate}
\tcblower
\textbf{Solution détaillée.}

\textbf{1. Angle de tir en $M_1$.}
Les angles $\widehat{AM_1B}$ (inscrit) et $\widehat{AO_1B}$ (au centre) interceptent le même arc $\overset{\frown}{AB}$ du cercle $\mathcal{C}_1$.
D'après le \textbf{théorème de l'angle inscrit et de l'angle au centre} :
\[ \widehat{AO_1B} = 2 \times \widehat{AM_1B} \implies \widehat{AM_1B} = \frac{\widehat{AO_1B}}{2} = \frac{60°}{2} = 30°. \]

\textbf{2. Angle au centre pour $M_2$.}
Les angles $\widehat{AM_2B}$ (inscrit) et $\widehat{AO_2B}$ (au centre) interceptent le même arc $\overset{\frown}{AB}$ du cercle $\mathcal{C}_2$.
D'après le \textbf{même théorème} :
\[ \widehat{AO_2B} = 2 \times \widehat{AM_2B} = 2 \times 45° = 90°. \]

\textbf{3. Nature de $[AB]$ pour $M_3$.}
On a $\widehat{AM_3B} = 90°$. Le triangle $AM_3B$ est donc rectangle en $M_3$.
D'après la \textbf{réciproque du théorème du triangle inscrit dans un demi-cercle}, l'hypoténuse $[AB]$ est un \textbf{diamètre} du cercle circonscrit $\mathcal{C}_3$. Le centre $O_3$ est le milieu de $[AB]$.

\textbf{4. Comparaison et conseil.}
Classement : $\widehat{AM_1B} = 30° \;<\; \widehat{AM_2B} = 45° \;<\; \widehat{AM_3B} = 90°$.
Le cercle $\mathcal{C}_1$ (angle centre $60°$) est le plus grand $\rightarrow$ position $M_1$ la plus éloignée.
Le cercle $\mathcal{C}_3$ (diamètre $[AB]$) est le plus petit $\rightarrow$ position $M_3$ la plus proche.

\textbf{Conseil pour la fiche technique :} « Pour maximiser son angle de tir, l'attaquant doit \textbf{se rapprocher de la ligne de but} (à l'intérieur de la surface de réparation). Tous les points d'un même arc de cercle passant par les poteaux offrent le même angle (Théorème 1), mais cet angle augmente quand le cercle diminue (Théorème 2). L'angle maximal théorique de $90°$ s'obtient sur le cercle de diamètre $[AB]$ (Théorème 3), soit au plus près des buts. »
\end{exoresolu}

\begin{exoresolu}[Angle de vue d'un écran géant]
\textbf{Énoncé.} Un écran géant est installé sur le boulevard du 20 Mai à Yaoundé. Sa largeur est $AB = 4$~m. Un spectateur se place en $M$ de sorte que l'angle sous lequel il voit l'écran soit $\widehat{AMB} = 30°$.
\begin{enumerate}
\item Quel est l'ensemble des positions où le spectateur voit l'écran sous un angle de $30°$ ?
\item Calculer l'angle au centre $\widehat{AOB}$ du cercle correspondant.
\item Si le spectateur recule pour se placer sur un cercle où l'angle au centre vaut $40°$, quel sera son nouvel angle de vue ? A-t-il intérêt à reculer pour mieux voir ?
\end{enumerate}
\tcblower
\textbf{Solution.}
\begin{enumerate}
\item L'ensemble des points $M$ tels que $\widehat{AMB} = 30°$ est l'\textbf{arc de cercle} $\overset{\frown}{AB}$ (de part et d'autre de $[AB]$, mais on prend celui devant l'écran) capable de $30°$, c'est-à-dire le lieu des points voyant $[AB]$ sous un angle de $30°$.
\item D'après le théorème angle au centre / angle inscrit : $\widehat{AOB} = 2 \times 30° = 60°$.
\item Nouvel angle au centre $= 40°$. Nouvel angle inscrit $\widehat{AM'B} = \frac{40°}{2} = 20°$.
\item Comme $20° < 30°$, l'angle de vue \textbf{diminue} en reculant. Le spectateur a intérêt à \textbf{s'approcher} (dans la limite du confort visuel) pour voir l'écran sous un angle plus grand.
\end{enumerate}
\end{exoresolu}

\courssub{Exercices d'entraînement (BEPC)}

\begin{exercice}[Angle de tir et placement]
Soit un but $[AB]$ de largeur $7,32$~m. Un joueur se place en $M$ tel que $\widehat{AMB} = 40°$.
\begin{enumerate}
\item Déterminer l'angle au centre $\widehat{AOB}$ du cercle passant par $A, M, B$.
\item Le joueur avance vers le but et se place en $M'$ sur le cercle de diamètre $[AB]$. Quelle est la valeur de $\widehat{AM'B}$ ? Justifier.
\item Comparer les deux angles de tir. Que conseilles-tu au joueur ?
\end{enumerate}
\end{exercice}

\begin{exercice}[Démonstration : Propriété des angles inscrits]
Soit un cercle de centre $O$. $A, B, M, N$ sont quatre points sur ce cercle, situés sur le même arc $\overset{\frown}{AB}$.
Démontrer que $\widehat{AMB} = \widehat{ANB}$.
\textit{Indication : Relier chaque angle inscrit à l'angle au centre $\widehat{AOB}$.}
\end{exercice}

\begin{exercice}[Problème ouvert : Le photographe]
Un photographe veut photographier l'hôtel de ville de Yaoundé (largeur de façade $AB = 20$~m). Il dispose d'un objectif offrant un angle de champ de $50°$.
\begin{enumerate}
\item Sur quel lieu géométrique doit-il se placer pour que la façade remplisse exactement son champ de vision (angle de vue $= 50°$) ?
\item Calculer l'angle au centre du cercle correspondant.
\item S'il recule de sorte que l'angle au centre devienne $30°$, quel sera l'angle de vue ? La façade rentrera-t-elle encore dans le cadre ?
\end{enumerate}
\end{exercice}

\courssub{Bilan et synthèse de la leçon}

\begin{aretenir}[Les 3 piliers des angles inscrits (Programme 3ème)]
\begin{enumerate}
\item \textbf{Égalité des angles inscrits même arc :} $\widehat{AMB} = \widehat{ANB}$ (même arc $\overset{\frown}{AB}$).
\item \textbf{Relation Angle au centre / Angle inscrit :} $\widehat{AOB} = 2 \times \widehat{AMB}$ (même arc $\overset{\frown}{AB}$).
\item \textbf{Triangle dans un demi-cercle :} $[AB]$ diamètre $\iff \widehat{AMB} = 90°$ (triangle rectangle en $M$).
\end{enumerate}
\end{aretenir}

\begin{savaistu}[Mathématiques et sport : La géométrie du but]
Le concept d'« angle de tir » est utilisé par les analystes de performance (data analysts) des grandes équipes (comme les Lions Indomptables ou le Real Madrid). Ils modélisent la surface de but visible par le joueur comme un \textbf{angle solide} en 3D, mais la base 2D reste l'angle inscrit $\widehat{AMB}$.
\begin{itemize}
\item \textbf{Expected Goals (xG) :} Les modèles statistiques de probabilité de marquer intègrent la distance \textbf{et} l'angle de tir. Un tir à 10~m face au but (angle $\approx 50°$) a un xG bien plus élevé qu'un tir à 10~m sur le côté (angle $\approx 15°$).
\item \textbf{Positionnement :} Les entraîneurs placent les attaquants sur des « lignes d'égal angle » (arcs de cercle) pour créer des combinaisons : un joueur attire le gardien, un autre sur le même arc a le même angle mais un décalage latéral.
\end{itemize}
La géométrie du cercle n'est pas qu'un exercice scolaire, c'est un outil de décision pour les champions !
\end{savaistu}

\begin{attention}[Check-list pour le BEPC : Rédaction parfaite]
\begin{itemize}
\item $\square$ \textbf{Citer le théorème} par son nom exact : « D'après le théorème de l'angle inscrit et de l'angle au centre... » ou « D'après la propriété des angles inscrits interceptant le même arc... ».
\item $\square$ \textbf{Préciser l'arc intercepté} : « ... qui interceptent le même arc $\overset{\frown}{AB}$ ».
\item $\square$ \textbf{Conclure par une phrase} : « Donc $\widehat{AMB} = 30°$. » ou « Donc le triangle $AMB$ est rectangle en $M$. »
\item $\square$ \textbf{Unités :} Ne jamais oublier le symbole $°$ pour les degrés.
\end{itemize}
\end{attention}

\newpage

\coursec{Méthodes et exercices résolus}

\coursec{Angles inscrits et angles au centre : définitions}

\begin{definition}[Angle inscrit, angle au centre, arc intercepté]
Soit $(\mathcal{C})$ un cercle de centre $O$.
\begin{itemize}
\item Un \cle{angle inscrit} dans $(\mathcal{C})$ est un angle dont le \textbf{sommet est sur le cercle} et dont les côtés coupent le cercle en deux points distincts.
\item Un \cle{angle au centre} de $(\mathcal{C})$ est un angle dont le \textbf{sommet est le centre $O$} du cercle et dont les côtés coupent le cercle en deux points distincts.
\item Dans les deux cas, les deux points d'intersection des côtés avec le cercle délimitent un \cle{arc intercepté}. On note cet arc $\overset{\frown}{AB}$ (arc $AB$ qui ne contient pas le sommet de l'angle).
\end{itemize}
Deux angles (inscrit et au centre, ou deux inscrits) sont dits \cle{associés} s'ils interceptent le \textbf{même arc}.
\end{definition}

\begin{methode}[Identifier un angle inscrit et l'angle au centre associé]
Pour bien appliquer les théorèmes de cette leçon, il faut d'abord savoir \emph{reconnaître} les angles en jeu sur une figure.
\begin{enumerate}
\item Vérifier que le \cle{sommet} de l'angle est \textbf{sur le cercle} : c'est un \cle{angle inscrit}.
\item Vérifier que le sommet de l'angle est le \textbf{centre} du cercle : c'est un \cle{angle au centre}.
\item Repérer les deux points où les côtés de l'angle coupent le cercle : ils délimitent l'\cle{arc intercepté} (l'arc \emph{opposé} au sommet).
\item Dire qu'un angle inscrit et un angle au centre sont \cle{associés} lorsqu'ils interceptent le \textbf{même arc}.
\end{enumerate}
\textbf{Réflexe :} nomme toujours l'arc intercepté (par exemple $\overset{\frown}{AB}$) avant de citer un théorème : c'est lui qui justifie tout le raisonnement.
\end{methode}

\begin{exoresolu}[Lecture d'une figure]
Sur la figure ci-dessous, $(\mathcal{C})$ est un cercle de centre $O$, et $A$, $B$, $M$ sont trois points de $(\mathcal{C})$.
\begin{enumerate}
\item Citer un angle inscrit et un angle au centre interceptant l'arc $\overset{\frown}{AB}$.
\item Le point $M$ étant sur l'arc $\overset{\frown}{AB}$, l'angle $\widehat{AMB}$ est-il un angle inscrit interceptant $\overset{\frown}{AB}$ ?
\end{enumerate}
\begin{popfigure}
\begin{tikzpicture}[scale=1.1]
\draw[popblue,thick] (0,0) circle (1.6);
\fill (0,0) circle (1.5pt) node[below right]{$O$};
\coordinate (A) at (150:1.6);
\coordinate (B) at (30:1.6);
\coordinate (M) at (-95:1.6);
\fill (A) circle (1.5pt) node[above left]{$A$};
\fill (B) circle (1.5pt) node[above right]{$B$};
\fill (M) circle (1.5pt) node[below]{$M$};
\draw[popink,thick] (A)--(0,0)--(B);
\draw[poporange,thick] (A)--(M)--(B);
\draw[popmuted,dashed] (A) -- (B);
\end{tikzpicture}
\legende{Angle au centre $\widehat{AOB}$ et angle inscrit $\widehat{AMB}$ associés}
\end{popfigure}
\tcblower
\textbf{1.} L'angle $\widehat{AOB}$ a pour sommet le centre $O$ du cercle : c'est un \textbf{angle au centre}. Ses côtés coupent le cercle en $A$ et $B$, donc il intercepte l'arc $\overset{\frown}{AB}$.

L'angle $\widehat{AMB}$ a pour sommet le point $M$ situé \textbf{sur le cercle} : c'est un \textbf{angle inscrit}. Ses côtés $[MA]$ et $[MB]$ coupent le cercle en $A$ et $B$ : il intercepte l'arc $\overset{\frown}{AB}$ (celui qui ne contient pas $M$).

\textbf{2.} Non. Si $M$ est \emph{sur} l'arc $\overset{\frown}{AB}$ considéré, alors l'angle $\widehat{AMB}$ intercepte l'\textbf{autre} arc $\overset{\frown}{AB}$ (le grand arc). Un angle inscrit intercepte toujours l'arc \textbf{opposé à son sommet}.

\textbf{Conclusion :} $\widehat{AOB}$ (angle au centre) et $\widehat{AMB}$ (angle inscrit, avec $M$ hors de l'arc $\overset{\frown}{AB}$) sont associés : ils interceptent le même arc $\overset{\frown}{AB}$.
\end{exoresolu}

\coursec{Théorème de l'angle inscrit et de l'angle au centre}

\begin{propriete}[Théorème de l'angle inscrit et de l'angle au centre]
Dans un cercle, un \cle{angle au centre} mesure le \textbf{double} d'un \cle{angle inscrit} qui intercepte le même arc.
\[ \widehat{AOB} = 2 \times \widehat{AMB} \qquad \text{ou} \qquad \widehat{AMB} = \frac{\widehat{AOB}}{2} \]
\textbf{Condition d'application :} le sommet de l'angle inscrit est sur le cercle, le sommet de l'angle au centre est le centre du cercle, et les deux angles interceptent le même arc (les sommets sont du même côté de la corde).
\end{propriete}

\begin{methode}[Calculer un angle avec le théorème de l'angle au centre]
Pour l'appliquer correctement :
\begin{enumerate}
\item Identifier l'arc commun intercepté par les deux angles.
\item Vérifier que le sommet de l'angle inscrit est bien \textbf{sur le cercle} et que l'angle au centre a bien pour sommet le \textbf{centre}.
\item Écrire la relation : angle au centre $= 2 \times$ angle inscrit, ou angle inscrit $= \dfrac{\text{angle au centre}}{2}$.
\item Conclure par une phrase avec la valeur de l'angle cherché.
\end{enumerate}
\textbf{Réflexe :} si tu connais l'angle inscrit, tu doubles ; si tu connais l'angle au centre, tu divises par $2$.
\end{methode}

\begin{exoresolu}[Calcul direct]
Soit $(\mathcal{C})$ un cercle de centre $O$. $A$, $B$ et $M$ sont trois points de $(\mathcal{C})$ tels que $\widehat{AOB} = 76^\circ$.
\begin{enumerate}
\item Calculer la mesure de l'angle $\widehat{AMB}$, sachant que $M$ n'appartient pas à l'arc $\overset{\frown}{AB}$ intercepté par $\widehat{AOB}$.
\item Le point $N$ est un autre point du cercle, situé sur le même arc que $M$. Que vaut $\widehat{ANB}$ ?
\end{enumerate}
\tcblower
\textbf{1.} L'angle $\widehat{AOB}$ est un \textbf{angle au centre} (sommet $O$, centre du cercle) et l'angle $\widehat{AMB}$ est un \textbf{angle inscrit} (sommet $M$ sur le cercle). Comme $M$ n'est pas sur l'arc $\overset{\frown}{AB}$, les deux angles interceptent le \textbf{même arc} $\overset{\frown}{AB}$.

D'après le \textbf{théorème de l'angle inscrit et de l'angle au centre} :
\[ \widehat{AMB} = \frac{\widehat{AOB}}{2} = \frac{76^\circ}{2} = 38^\circ. \]

\textbf{2.} $N$ est sur le même arc que $M$, donc $\widehat{ANB}$ intercepte le même arc $\overset{\frown}{AB}$ que $\widehat{AMB}$. Deux angles inscrits qui interceptent le même arc ont la même mesure (propriété vue dans la section suivante), donc :
\[ \widehat{ANB} = \widehat{AMB} = 38^\circ. \]

\textbf{Conclusion :} $\widehat{AMB} = \widehat{ANB} = 38^\circ$.
\end{exoresolu}

\begin{experience}[Découverte guidée : le théorème de l'angle au centre]
\textbf{Matériel :} papier calque, règle, rapporteur, compas.
\begin{enumerate}
\item Trace un cercle $(\mathcal{C})$ de centre $O$. Place trois points $A$, $B$, $M$ sur le cercle ($M$ sur le grand arc $AB$).
\item Trace l'angle au centre $\widehat{AOB}$ et l'angle inscrit $\widehat{AMB}$. Mesure-les au rapporteur.
\item Recommence avec un autre point $M'$ sur le même arc. Que observes-tu pour $\widehat{AM'B}$ ?
\item Trace la droite $(OM)$. Elle coupe le cercle en $M$ et $M''$. Utilise le théorème de l'angle au centre dans les triangles isocèles $OAM$ et $OBM$ pour \textbf{démontrer} que $\widehat{AOB} = 2\widehat{AMB}$.
\end{enumerate}
\textbf{Piste de démonstration :} Dans $\triangle OAM$, $OA=OM$ (rayons) donc $\widehat{OAM}=\widehat{OMA}$. L'angle extérieur $\widehat{AOX} = \widehat{OAM}+\widehat{OMA} = 2\widehat{OMA}$. De même $\widehat{BOX}=2\widehat{OMB}$. En additionnant (ou soustrayant selon la figure) : $\widehat{AOB}=2\widehat{AMB}$.
\end{experience}

\coursec{Angles inscrits interceptant le même arc}

\begin{propriete}[Angles inscrits interceptant le même arc]
Deux \cle{angles inscrits} qui interceptent le \textbf{même arc} ont la \textbf{même mesure}.
\textbf{Condition :} les sommets des deux angles sont du \textbf{même côté} de la corde joignant les extrémités de l'arc.
\end{propriete}

\begin{methode}[Prouver que deux angles sont égaux]
\begin{enumerate}
\item Repérer les deux angles inscrits (sommets sur le cercle).
\item Montrer que leurs côtés coupent le cercle aux \textbf{mêmes deux points} (donc même arc intercepté).
\item Vérifier que les deux sommets sont du \textbf{même côté} de la corde (sinon les angles sont supplémentaires, voir section suivante).
\item Conclure à l'égalité des mesures.
\end{enumerate}
\textbf{Cas particulier utile :} des angles inscrits interceptant le même arc servent souvent à démontrer que deux droites sont parallèles (angles alternes-internes égaux) ou que des triangles sont semblables.
\end{methode}

\begin{exoresolu}[Démontrer une égalité d'angles]
$A$, $B$, $C$, $D$ sont quatre points d'un cercle $(\mathcal{C})$, placés dans cet ordre. Les cordes $[AC]$ et $[BD]$ se coupent en $I$.
\begin{enumerate}
\item Démontrer que $\widehat{BAC} = \widehat{BDC}$.
\item En déduire que les triangles $IAB$ et $IDC$ ont des angles deux à deux égaux.
\end{enumerate}
\tcblower
\textbf{1.} L'angle $\widehat{BAC}$ est inscrit dans $(\mathcal{C})$ (sommet $A$ sur le cercle) et ses côtés coupent le cercle en $B$ et $C$ : il intercepte l'arc $\overset{\frown}{BC}$.

L'angle $\widehat{BDC}$ est inscrit dans $(\mathcal{C})$ (sommet $D$ sur le cercle) et ses côtés coupent le cercle en $B$ et $C$ : il intercepte aussi l'arc $\overset{\frown}{BC}$.

Comme $A$ et $D$ sont du même côté de la corde $[BC]$ (points placés dans l'ordre $A$, $B$, $C$, $D$), les deux angles inscrits interceptent le \textbf{même arc} $\overset{\frown}{BC}$. Donc :
\[ \widehat{BAC} = \widehat{BDC}. \]

\textbf{2.} Dans les triangles $IAB$ et $IDC$ :
\begin{itemize}
\item $\widehat{IAB} = \widehat{BAC}$ et $\widehat{IDC} = \widehat{BDC}$, donc $\widehat{IAB} = \widehat{IDC}$ d'après la question 1 ;
\item les angles $\widehat{AIB}$ et $\widehat{DIC}$ sont \textbf{opposés par le sommet}, donc égaux ;
\item les troisièmes angles sont alors égaux : $\widehat{IBA} = \widehat{ICD}$ (la somme des angles d'un triangle vaut $180^\circ$).
\end{itemize}
\textbf{Conclusion :} les triangles $IAB$ et $IDC$ ont leurs trois angles deux à deux égaux (ils sont semblables).
\end{exoresolu}

\coursec{Quadrilatère inscriptible et angles supplémentaires}

\begin{propriete}[Angles opposés d'un quadrilatère inscrit]
Dans un \cle{quadrilatère inscrit} dans un cercle, les \textbf{angles opposés sont supplémentaires} (leur somme vaut $180^\circ$).
\[ \widehat{A} + \widehat{C} = 180^\circ \qquad \text{et} \qquad \widehat{B} + \widehat{D} = 180^\circ \]
\textbf{Justification :} les angles opposés interceptent des arcs complémentaires (le petit et le grand arc) dont les angles au centre totalisent $360^\circ$. Les angles inscrits valent la moitié : $360^\circ / 2 = 180^\circ$.
\end{propriete}

\begin{methode}[Angles inscrits de part et d'autre d'une corde]
Lorsque deux angles inscrits interceptent le même arc mais que leurs sommets sont de \textbf{part et d'autre} de la corde :
\begin{enumerate}
\item Chacun intercepte un arc différent (le petit et le grand arc $\overset{\frown}{AB}$), dont les angles au centre associés totalisent $360^\circ$.
\item Les deux angles inscrits sont alors \cle{supplémentaires} : leur somme vaut $180^\circ$.
\item Conséquence : dans un quadrilatère inscrit, les angles opposés sont supplémentaires.
\end{enumerate}
\textbf{Réflexe :} avant d'écrire « angles égaux », vérifie toujours la position des sommets par rapport à la corde !
\end{methode}

\begin{exoresolu}[Quadrilatère inscrit]
$ABCD$ est un quadrilatère inscrit dans un cercle $(\mathcal{C})$. On donne $\widehat{BAD} = 65^\circ$ et $\widehat{ABC} = 98^\circ$.
\begin{enumerate}
\item Calculer $\widehat{BCD}$.
\item Calculer $\widehat{ADC}$.
\end{enumerate}
\tcblower
\textbf{1.} Les angles $\widehat{BAD}$ et $\widehat{BCD}$ sont deux angles \textbf{opposés} du quadrilatère $ABCD$ inscrit dans le cercle $(\mathcal{C})$. Or, dans un quadrilatère inscrit, les angles opposés sont supplémentaires :
\[ \widehat{BAD} + \widehat{BCD} = 180^\circ \]
\[ \widehat{BCD} = 180^\circ - 65^\circ = 115^\circ. \]

\textbf{2.} De même, $\widehat{ABC}$ et $\widehat{ADC}$ sont opposés :
\[ \widehat{ABC} + \widehat{ADC} = 180^\circ \]
\[ \widehat{ADC} = 180^\circ - 98^\circ = 82^\circ. \]

\textbf{Vérification :} $65^\circ + 98^\circ + 115^\circ + 82^\circ = 360^\circ$, ce qui est bien la somme des angles d'un quadrilatère.

\textbf{Conclusion :} $\widehat{BCD} = 115^\circ$ et $\widehat{ADC} = 82^\circ$.
\end{exoresolu}

\coursec{Triangle inscrit dans un demi-cercle}

\begin{propriete}[Théorème de l'angle inscrit dans un demi-cercle (Réciproque de Thalès pour le cercle)]
\begin{itemize}
\item \textbf{Direct :} Si un triangle est inscrit dans un cercle et qu'un de ses côtés est un \cle{diamètre} du cercle, alors ce triangle est \textbf{rectangle} (l'angle opposé au diamètre vaut $90^\circ$).
\item \textbf{Réciproque :} Si un triangle rectangle est inscrit dans un cercle, alors son \textbf{hypoténuse est un diamètre} du cercle (le centre du cercle est le milieu de l'hypoténuse).
\end{itemize}
\textbf{Justification :} L'angle inscrit vaut la moitié de l'angle au centre associé. Ici l'angle au centre est un angle plat ($180^\circ$), d'où $\frac{180^\circ}{2} = 90^\circ$.
\end{propriete}

\begin{methode}[Utiliser le diamètre pour obtenir un angle droit]
\begin{enumerate}
\item Repérer le diamètre (segment passant par le centre $O$, ou dont le milieu est $O$).
\item Identifier l'angle opposé au diamètre : c'est l'angle droit.
\item Utiliser ensuite les propriétés du triangle rectangle (Pythagore, trigonométrie, angles complémentaires).
\end{enumerate}
\end{methode}

\begin{exoresolu}[Démontrer qu'un triangle est rectangle]
Soit $(\mathcal{C})$ un cercle de centre $O$ et de diamètre $[AB]$ tel que $AB = 10$ cm. $M$ est un point de $(\mathcal{C})$ tel que $AM = 6$ cm.
\begin{enumerate}
\item Démontrer que le triangle $AMB$ est rectangle en $M$.
\item Calculer la longueur $MB$.
\end{enumerate}
\tcblower
\textbf{1.} Le triangle $AMB$ est inscrit dans le cercle $(\mathcal{C})$ et le côté $[AB]$ est un \textbf{diamètre} de ce cercle. Or, tout triangle inscrit dans un cercle dont un côté est un diamètre est rectangle, l'angle droit étant opposé au diamètre.

Donc le triangle $AMB$ est \textbf{rectangle en $M$}, c'est-à-dire $\widehat{AMB} = 90^\circ$.

\textbf{2.} Dans le triangle $AMB$ rectangle en $M$, d'après le \textbf{théorème de Pythagore} :
\begin{align*}
AB^2 &= AM^2 + MB^2 \\
10^2 &= 6^2 + MB^2 \\
100 &= 36 + MB^2 \\
MB^2 &= 100 - 36 = 64 \\
MB &= \sqrt{64} = 8 \text{ cm}.
\end{align*}
\textbf{Conclusion :} le triangle $AMB$ est rectangle en $M$ et $MB = 8$ cm.
\end{exoresolu}

\coursec{Calculs d'angles, démonstrations et applications}

\begin{methode}[Construire un raisonnement en plusieurs étapes]
Pour les exercices de type BEPC mêlant plusieurs propriétés :
\begin{enumerate}
\item \textbf{Lire et coder la figure} : reporter toutes les données (centre, diamètre, angles connus).
\item \textbf{Nommer précisément} chaque angle utilisé et l'arc qu'il intercepte.
\item \textbf{Citer le théorème} à chaque étape : « angle au centre $=$ double de l'angle inscrit », « angles inscrits interceptant le même arc », « triangle inscrit dans un demi-cercle ».
\item \textbf{Enchaîner} les égalités ou calculs jusqu'à la conclusion demandée.
\item \textbf{Rédiger une phrase de conclusion} qui répond exactement à la question posée.
\end{enumerate}
\end{methode}

\begin{exoresolu}[Démonstration type BEPC]
Soit $(\mathcal{C})$ un cercle de centre $O$, $[AB]$ un diamètre de $(\mathcal{C})$ et $C$ un point de $(\mathcal{C})$ tel que $\widehat{BOC} = 56^\circ$.
\begin{enumerate}
\item Calculer $\widehat{BAC}$.
\item En déduire $\widehat{ABC}$.
\end{enumerate}
\tcblower
\textbf{1.} L'angle $\widehat{BOC}$ est un \textbf{angle au centre} (sommet $O$) interceptant l'arc $\overset{\frown}{BC}$. L'angle $\widehat{BAC}$ est un \textbf{angle inscrit} (sommet $A$ sur le cercle) interceptant le même arc $\overset{\frown}{BC}$.

D'après le théorème de l'angle inscrit et de l'angle au centre :
\[ \widehat{BAC} = \frac{\widehat{BOC}}{2} = \frac{56^\circ}{2} = 28^\circ. \]

\textbf{2.} Le triangle $ABC$ est inscrit dans le cercle $(\mathcal{C})$ et $[AB]$ est un \textbf{diamètre} : le triangle $ABC$ est donc \textbf{rectangle en $C$}, soit $\widehat{ACB} = 90^\circ$.

La somme des angles d'un triangle vaut $180^\circ$, donc :
\begin{align*}
\widehat{ABC} &= 180^\circ - \widehat{ACB} - \widehat{BAC} \\
&= 180^\circ - 90^\circ - 28^\circ \\
&= 62^\circ.
\end{align*}
\textbf{Conclusion :} $\widehat{BAC} = 28^\circ$ et $\widehat{ABC} = 62^\circ$.
\end{exoresolu}

\begin{methode}[Modéliser une situation de la vie courante]
\begin{enumerate}
\item \textbf{Traduire} l'énoncé en schéma : le cercle (stade, rond-point, terrain), le centre, les points sur le cercle.
\item Identifier dans la réalité ce qui joue le rôle de l'\textbf{angle inscrit} (point de vue d'une personne sur le cercle) ou de l'\textbf{angle au centre}.
\item Appliquer le théorème adapté.
\item \textbf{Interpréter} le résultat dans le contexte (phrase avec unité ou degrés).
\end{enumerate}
\end{methode}

\begin{exoresolu}[Le stade de football]
Lors d'un match au stade Ahmadou Ahidjo, on modélise la tribune circulaire par un cercle de centre $O$. Deux poteaux de corner sont placés aux points $A$ et $B$ du cercle. Un cameraman placé au centre $O$ filme les deux poteaux sous un angle $\widehat{AOB} = 84^\circ$. Un supporter assis au point $M$ de la tribune (sur le cercle, de l'autre côté de la corde $[AB]$) regarde ces mêmes poteaux.
\begin{enumerate}
\item Sous quel angle le supporter voit-il les deux poteaux ?
\item Un second supporter, placé en $N$ sur le même arc que $M$, voit-il les poteaux sous le même angle ? Justifier.
\end{enumerate}
\tcblower
\textbf{1.} L'angle $\widehat{AOB}$ est un angle au centre et l'angle $\widehat{AMB}$ est un angle inscrit ; ils interceptent le même arc $\overset{\frown}{AB}$ (celui qui ne contient pas $M$). D'après le théorème de l'angle inscrit et de l'angle au centre :
\[ \widehat{AMB} = \frac{\widehat{AOB}}{2} = \frac{84^\circ}{2} = 42^\circ. \]
Le supporter voit les deux poteaux sous un angle de $42^\circ$.

\textbf{2.} Oui. Les angles $\widehat{AMB}$ et $\widehat{ANB}$ sont deux angles inscrits dont les sommets $M$ et $N$ sont sur le même arc : ils interceptent le \textbf{même arc} $\overset{\frown}{AB}$. Or deux angles inscrits interceptant le même arc sont égaux, donc :
\[ \widehat{ANB} = \widehat{AMB} = 42^\circ. \]
\textbf{Conclusion :} tous les supporters placés sur le même arc de la tribune voient les poteaux sous le même angle de $42^\circ$ : c'est pourquoi les places d'une même section circulaire offrent une vue comparable.
\end{exoresolu}

\begin{savaistu}[Thalès, les pyramides et la navigation]
\emph{Thalès de Milet} (vers -625 à -547), l'un des sept sages de la Grèce antique, serait allé en Égypte et aurait mesuré la hauteur de la grande pyramide de Gizeh en utilisant l'ombre d'un bâton vertical (théorème de Thalès). Mais il a aussi énoncé le théorème de l'angle inscrit dans un demi-cercle : « Tout angle inscrit dans un demi-cercle est un angle droit ». Cette propriété est à la base de la \textbf{navigation astronomique} : en mesurant l'angle entre l'horizon et une étoile (angle inscrit), le navigateur se situe sur un cercle de position (cercle de hauteur) dont il connaît le centre géographique. Au Cameroun, les géomètres-topographes utilisent encore ces propriétés pour l'implantation des routes et des réseaux électriques (CAMTEL, ENEO) en terrain difficile.
\end{savaistu}

\coursec{Entraînement et erreurs à éviter}

\begin{attention}[Les 5 erreurs qui coûtent des points au BEPC]
\begin{enumerate}
\item \textbf{Confondre angle inscrit et angle au centre.} Avant d'écrire la relation, vérifie la position du sommet : sur le cercle (inscrit) ou au centre (au centre). Cite toujours le théorème par son nom dans la rédaction.
\item \textbf{Inverser le facteur 2.} L'angle au centre est le \textbf{double} de l'angle inscrit : $\widehat{AOB} = 2\,\widehat{AMB}$. Si l'angle au centre vaut $76^\circ$, l'angle inscrit vaut $38^\circ$, pas $152^\circ$ ! Demande-toi toujours : « lequel est le plus grand ? »
\item \textbf{Oublier la condition « même arc ».} Deux angles inscrits ne sont égaux que s'ils interceptent le \emph{même} arc, c'est-à-dire si leurs sommets sont du même côté de la corde. Si les sommets sont de part et d'autre, les angles sont \textbf{supplémentaires} (somme $180^\circ$), pas égaux.
\item \textbf{Oublier de justifier l'angle droit.} Écrire « le triangle est rectangle » sans dire que le côté est un \emph{diamètre} du cercle circonscrit fait perdre le point de justification. Rédige : « le triangle est inscrit dans le cercle et $[AB]$ est un diamètre, donc il est rectangle ».
\item \textbf{Négliger la conclusion et les unités.} Termine toujours par une phrase : « donc $\widehat{AMB} = 38^\circ$ ». Un calcul juste sans phrase de conclusion, ou avec une mesure sans le symbole de degré, est pénalisé.
\end{enumerate}
\end{attention}

\begin{exercice}[Entraînement BEPC]
\begin{enumerate}
\item Soit un cercle de centre $O$. $A$, $B$, $C$ sont trois points du cercle tels que $\widehat{AOB} = 110^\circ$ et $\widehat{BOC} = 80^\circ$. Calculer $\widehat{ACB}$ et $\widehat{BAC}$.
\item $ABCD$ est un quadrilatère inscrit dans un cercle de centre $O$. On sait que $\widehat{DAB} = 70^\circ$ et $\widehat{ABC} = 95^\circ$. Calculer $\widehat{BCD}$ et $\widehat{CDA}$. Le quadrilatère est-il un trapèze ? Justifier.
\item $[AB]$ est un diamètre d'un cercle de centre $O$. $M$ est un point du cercle tel que $\widehat{BAM} = 35^\circ$. Calculer $\widehat{ABM}$ et $\widehat{AMB}$. Quelle est la nature du triangle $AMB$ ?
\item \textbf{Problème concret :} Un rond-point circulaire a un centre $O$. Trois sorties $A$, $B$, $C$ sont sur le pourtour. Un panneau au centre $O$ indique la direction de $A$ et $B$ sous un angle de $100^\circ$. Un automobiliste arrive par la sortie $C$ (sur l'arc $AB$ ne contenant pas $C$). Sous quel angle voit-il les panneaux $A$ et $B$ ?
\end{enumerate}
\end{exercice}

\begin{corrige}[Corrigés des exercices d'entraînement]
\textbf{Exercice 1 :}
L'arc $\overset{\frown}{AB}$ correspond à $\widehat{AOB}=110^\circ$. L'angle inscrit $\widehat{ACB}$ intercepte cet arc (sommet $C$ sur l'autre arc).
$\widehat{ACB} = \frac{110^\circ}{2} = 55^\circ$.
L'arc $\overset{\frown}{BC}$ correspond à $\widehat{BOC}=80^\circ$. $\widehat{BAC}$ intercepte cet arc.
$\widehat{BAC} = \frac{80^\circ}{2} = 40^\circ$.

\textbf{Exercice 2 :}
Quadrilatère inscrit $\Rightarrow$ angles opposés supplémentaires.
$\widehat{BCD} = 180^\circ - \widehat{DAB} = 180^\circ - 70^\circ = 110^\circ$.
$\widehat{CDA} = 180^\circ - \widehat{ABC} = 180^\circ - 95^\circ = 85^\circ$.
Vérification : $70+95+110+85 = 360^\circ$.
Trapèze ? Non, car angles adjacents non supplémentaires ($70+95 \neq 180$, $95+110 \neq 180$, etc.) et pas de côtés parallèles évidents.

\textbf{Exercice 3 :}
$[AB]$ diamètre $\Rightarrow$ triangle $AMB$ rectangle en $M$ (angle inscrit interceptant demi-cercle).
$\widehat{AMB} = 90^\circ$.
Somme angles triangle $AMB$ : $\widehat{ABM} = 180^\circ - 90^\circ - 35^\circ = 55^\circ$.
Triangle rectangle en $M$.

\textbf{Exercice 4 :}
Angle au centre $\widehat{AOB} = 100^\circ$. Angle inscrit $\widehat{ACB}$ (automobiliste en $C$) interceptant même arc $\overset{\frown}{AB}$.
$\widehat{ACB} = \frac{100^\circ}{2} = 50^\circ$.
L'automobiliste voit les sorties sous un angle de $50^\circ$.
\end{corrige}

\newpage

\coursec{Exercices}

\courssub{Je vérifie mes connaissances}

\begin{exercice}[QCM : à chaque question, une seule bonne réponse]
Pour chaque question, entoure la bonne réponse.
\begin{enumerate}
\item Un angle inscrit dans un cercle est un angle dont :
\begin{enumerate}
\item le sommet est le centre du cercle ;
\item le sommet est sur le cercle et dont les côtés coupent le cercle ;
\item le sommet est à l'extérieur du cercle.
\end{enumerate}
\item Dans un cercle, un angle au centre mesure $70^{\circ}$. Un angle inscrit interceptant le même arc mesure :
\begin{enumerate}
\item $140^{\circ}$ ; \item $70^{\circ}$ ; \item $35^{\circ}$.
\end{enumerate}
\item Deux angles inscrits qui interceptent le même arc d'un cercle sont :
\begin{enumerate}
\item complémentaires ; \item supplémentaires ; \item de même mesure.
\end{enumerate}
\item Un triangle inscrit dans un cercle dont l'un des côtés est un diamètre est :
\begin{enumerate}
\item équilatéral ; \item rectangle ; \item isocèle.
\end{enumerate}
\end{enumerate}
\end{exercice}

\begin{exercice}[Vrai ou faux ? Justifier]
Réponds par VRAI ou FAUX, puis justifie ta réponse par une propriété du cours ou un contre-exemple.
\begin{enumerate}
\item L'angle au centre associé à un angle inscrit mesure le double de cet angle inscrit.
\item Deux angles inscrits d'un même cercle sont toujours égaux.
\item Si un angle inscrit mesure $45^{\circ}$, alors l'angle au centre interceptant le même arc mesure $90^{\circ}$.
\item Tout triangle inscrit dans un cercle est rectangle.
\end{enumerate}
\end{exercice}

\begin{exercice}[Définitions]
\begin{enumerate}
\item Définis un angle inscrit dans un cercle et illustre ta définition par une figure codée.
\item Définis un angle au centre. Qu'appelle-t-on \emph{angle inscrit et angle au centre associés} ?
\item Énonce la propriété des angles inscrits interceptant le même arc.
\end{enumerate}
\end{exercice}

\begin{exercice}[Calculs directs]
Dans un cercle de centre $O$, l'angle au centre $\widehat{AOB}$ mesure $110^{\circ}$.
\begin{enumerate}
\item Calcule la mesure de tout angle inscrit interceptant l'arc mineur $\overset{\frown}{AB}$.
\item Calcule la mesure d'un angle inscrit interceptant l'arc majeur $\overset{\frown}{AB}$.
\end{enumerate}
\end{exercice}

\courssub{Je m'entraîne}

\begin{exercice}[Même arc, mêmes angles]
Soit un cercle de centre $O$ et quatre points $A$, $B$, $M$, $N$ situés sur ce cercle, $M$ et $N$ étant sur le même arc d'extrémités $A$ et $B$. On donne $\widehat{AMB} = 42^{\circ}$.
\begin{enumerate}
\item Justifie que $\widehat{ANB} = 42^{\circ}$ en citant précisément la propriété utilisée.
\item Calcule la mesure de l'angle au centre $\widehat{AOB}$ associé à ces angles inscrits.
\end{enumerate}
\end{exercice}

\begin{exercice}[Triangle inscrit dans un demi-cercle]
Un triangle $ABC$ est inscrit dans un cercle de diamètre $[BC]$, avec $BC = 13$~cm et $AB = 5$~cm.
\begin{enumerate}
\item Montre que le triangle $ABC$ est rectangle, et précise en quel point.
\item Calcule la longueur $AC$.
\end{enumerate}
\end{exercice}

\begin{exercice}[Chasse aux angles]
Sur un cercle de centre $O$, on place quatre points $A$, $B$, $C$, $D$ dans cet ordre. On donne $\widehat{AOB} = 80^{\circ}$ et $\widehat{BOC} = 60^{\circ}$.
\begin{enumerate}
\item Calcule $\widehat{ADB}$ et $\widehat{ACB}$.
\item Calcule $\widehat{BAC}$ et $\widehat{BDC}$.
\item Que peux-tu dire des angles $\widehat{ADB}$ et $\widehat{ACB}$ ? Justifie.
\end{enumerate}
\end{exercice}

\begin{exercice}[Réciproque du demi-cercle]
Un triangle $MNP$ est rectangle en $M$, avec $MN = 6$~cm et $MP = 8$~cm.
\begin{enumerate}
\item Calcule la longueur $NP$.
\item Justifie que le cercle circonscrit au triangle $MNP$ a pour diamètre $[NP]$, puis donne son rayon.
\end{enumerate}
\end{exercice}

\courssub{Je me prépare au BEPC}

\begin{exercice}[Type BEPC : angles d'un triangle inscrit]
Sur la figure ci-dessous, $(\mathcal{C})$ est un cercle de centre $O$, $[AD]$ est un diamètre, et $[AB]$ et $[AC]$ sont des cordes. On donne $\widehat{BOD} = 80^{\circ}$ et $\widehat{COD} = 50^{\circ}$.
\begin{popfigure}
\begin{tikzpicture}[scale=1.1]
\draw[popblue, thick] (0,0) circle (2);
\coordinate (O) at (0,0);
\coordinate (A) at (2,0);
\coordinate (D) at (-2,0);
\coordinate (B) at (100:2);
\coordinate (C) at (230:2);
\draw[popdark] (A)--(D);
\draw[poporange, thick] (A)--(B);
\draw[popgreen, thick] (A)--(C);
\draw[popdark] (B)--(C);
\draw[popmuted] (O)--(B) (O)--(C);
\foreach \p/\pos in {A/right, D/left, B/above left, C/below left, O/below right}
  \fill[popink] (\p) circle (1.5pt) node[\pos] {$\p$};
\end{tikzpicture}
\legende{Cercle de centre $O$, diamètre $[AD]$, cordes $[AB]$ et $[AC]$.}
\end{popfigure}
\begin{enumerate}
\item Calcule la mesure de l'angle $\widehat{BAD}$, en justifiant par le théorème de l'angle inscrit.
\item Calcule la mesure de l'angle $\widehat{CAD}$.
\item Déduis-en la mesure de l'angle $\widehat{BAC}$.
\item Calcule la mesure de l'angle au centre $\widehat{BOC}$, puis celle de l'angle $\widehat{ABC}$.
\item Vérifie que la somme des angles du triangle $ABC$ vaut $180^{\circ}$.
\end{enumerate}
\end{exercice}

\begin{exercice}[Problème concret : la piste d'athlétisme]
Lors des championnats scolaires de Yaoundé, deux projecteurs placés au centre $O$ d'une piste d'athlétisme circulaire éclairent la ligne d'arrivée matérialisée par deux poteaux $A$ et $B$ plantés sur la piste. L'angle $\widehat{AOB}$ mesure $130^{\circ}$. Un spectateur est assis au point $S$ de la tribune circulaire, sur l'arc majeur $\overset{\frown}{AB}$.
\begin{enumerate}
\item Fais une figure codée de la situation.
\item Sous quel angle le spectateur voit-il la ligne d'arrivée $[AB]$ ? Justifie ta réponse en citant la propriété utilisée.
\item Un second spectateur est assis au point $T$, sur l'arc mineur $\overset{\frown}{AB}$. Sous quel angle voit-il la ligne d'arrivée ? Justifie.
\item Les deux spectateurs voient-ils la ligne d'arrivée sous le même angle ? Commente.
\end{enumerate}
\end{exercice}

\begin{exercice}[Démonstration et synthèse]
$M$, $N$, $P$, $Q$ sont quatre points d'un cercle $(\mathcal{C})$ de centre $O$, tels que les angles $\widehat{MPN}$ et $\widehat{MQN}$ interceptent le même arc $\overset{\frown}{MN}$. Les droites $(MP)$ et $(NQ)$ se coupent en un point $I$.
\begin{enumerate}
\item Montre que $\widehat{MPN} = \widehat{MQN}$, en citant la propriété utilisée.
\item On donne $\widehat{MON} = 96^{\circ}$. Calcule $\widehat{MPN}$ et $\widehat{MQN}$.
\item On suppose de plus que le triangle $MNP$ est inscrit dans un demi-cercle de diamètre $[NP]$. Déduis-en la nature du triangle $MNP$ et la mesure de l'angle $\widehat{NMP}$.
\item En utilisant l'égalité des angles démontrée à la question 1, montre que les triangles $IMP$ et $IQN$ ont deux angles respectivement égaux. Que peut-on en conclure sur leurs troisièmes angles ?
\end{enumerate}
\end{exercice}

\newpage

\coursec{Corrigés des exercices}

\begin{corrige}[Exercice 1 — QCM : à chaque question, une seule bonne réponse]
\begin{enumerate}
\item \textbf{Réponse b.} Un angle inscrit dans un cercle est un angle dont le sommet est \cle{sur le cercle} et dont les côtés coupent le cercle. (La réponse a décrivait un angle au centre, dont le sommet est le centre du cercle.)
\item \textbf{Réponse c : $35^{\circ}$.} D'après le théorème de l'angle inscrit, l'angle inscrit mesure la moitié de l'angle au centre interceptant le même arc : $\dfrac{70^{\circ}}{2} = 35^{\circ}$.
\item \textbf{Réponse c : de même mesure.} Deux angles inscrits qui interceptent le même arc d'un cercle sont égaux.
\item \textbf{Réponse b : rectangle.} Si un triangle est inscrit dans un cercle dont l'un des côtés est un diamètre, alors ce triangle est rectangle : l'angle opposé au diamètre intercepte un demi-cercle, donc mesure $\dfrac{180^{\circ}}{2} = 90^{\circ}$.
\end{enumerate}
\end{corrige}

\begin{corrige}[Exercice 2 — Vrai ou faux ? Justifier]
\begin{enumerate}
\item \textbf{VRAI.} C'est le théorème de l'angle inscrit : dans un cercle, l'angle au centre associé à un angle inscrit (c'est-à-dire interceptant le même arc) mesure le double de cet angle inscrit.
\item \textbf{FAUX.} Deux angles inscrits ne sont égaux que s'ils interceptent le \emph{même arc}. Contre-exemple : si $\widehat{AOB} = 60^{\circ}$ et $\widehat{COD} = 100^{\circ}$ sont deux angles au centre associés à des arcs distincts, les angles inscrits associés mesurent $30^{\circ}$ et $50^{\circ}$, qui ne sont pas égaux.
\item \textbf{VRAI.} Par le théorème de l'angle inscrit, l'angle au centre interceptant le même arc mesure $2 \times 45^{\circ} = 90^{\circ}$.
\item \textbf{FAUX.} Un triangle inscrit dans un cercle n'est rectangle que si l'un de ses côtés est un diamètre du cercle. Contre-exemple : un triangle équilatéral inscrit dans un cercle a trois angles de $60^{\circ}$ ; il n'est pas rectangle.
\end{enumerate}
\end{corrige}

\begin{corrige}[Exercice 3 — Définitions]
\begin{enumerate}
\item \textbf{Définition.} Un \cle{angle inscrit} dans un cercle est un angle dont le sommet appartient au cercle et dont les côtés coupent le cercle. On dit qu'il \emph{intercepte} l'arc compris entre ses deux côtés (l'arc ne contenant pas son sommet).
\begin{popfigure}
\begin{tikzpicture}[scale=1]
\draw[popblue, thick] (0,0) circle (2);
\coordinate (M) at (90:2);
\coordinate (A) at (200:2);
\coordinate (B) at (340:2);
\draw[poporange, thick] (M)--(A) (M)--(B);
\draw[popgreen, very thick] (200:2) arc (200:340:2);
\foreach \p/\pos in {M/above, A/below left, B/below right}
  \fill[popink] (\p) circle (1.5pt) node[\pos] {$\p$};
\draw[popdark] (M) ++(200:0.6) arc (200:340:0.6);
\node[popdark] at (270:1.3) {$\widehat{AMB}$};
\end{tikzpicture}
\legende{$\widehat{AMB}$ est un angle inscrit interceptant l'arc $\overset{\frown}{AB}$ (en vert).}
\end{popfigure}
\item \textbf{Définition.} Un \cle{angle au centre} est un angle dont le sommet est le centre du cercle. On dit qu'un angle inscrit $\widehat{AMB}$ et un angle au centre $\widehat{AOB}$ sont \cle{associés} lorsqu'ils interceptent le même arc $\overset{\frown}{AB}$.
\item \textbf{Propriété.} Dans un cercle, deux angles inscrits qui interceptent le même arc ont la même mesure.
\end{enumerate}
\end{corrige}

\begin{corrige}[Exercice 4 — Calculs directs]
\begin{enumerate}
\item L'angle au centre $\widehat{AOB} = 110^{\circ}$ intercepte l'arc mineur $\overset{\frown}{AB}$. D'après le théorème de l'angle inscrit, tout angle inscrit interceptant ce même arc mesure :
\[ \frac{\widehat{AOB}}{2} = \frac{110^{\circ}}{2} = 55^{\circ}. \]
\item L'arc majeur $\overset{\frown}{AB}$ correspond à l'angle au centre \emph{rentrant} : $360^{\circ} - 110^{\circ} = 250^{\circ}$. Un angle inscrit interceptant l'arc majeur mesure donc :
\[ \frac{250^{\circ}}{2} = 125^{\circ}. \]
\textbf{Vérification :} $55^{\circ} + 125^{\circ} = 180^{\circ}$ : les angles inscrits interceptant deux arcs complémentaires sont supplémentaires, c'est cohérent.
\end{enumerate}
\end{corrige}

\begin{corrige}[Exercice 5 — Même arc, mêmes angles]
\begin{enumerate}
\item Les angles $\widehat{AMB}$ et $\widehat{ANB}$ sont deux angles inscrits dans le même cercle, et leurs sommets $M$ et $N$ sont sur le même arc d'extrémités $A$ et $B$ : ils interceptent donc le \emph{même arc} $\overset{\frown}{AB}$. Or, \textbf{deux angles inscrits interceptant le même arc sont égaux}. Donc :
\[ \widehat{ANB} = \widehat{AMB} = 42^{\circ}. \]
\item L'angle au centre $\widehat{AOB}$ est associé à l'angle inscrit $\widehat{AMB}$ (ils interceptent le même arc $\overset{\frown}{AB}$). D'après le théorème de l'angle inscrit :
\[ \widehat{AOB} = 2 \times \widehat{AMB} = 2 \times 42^{\circ} = 84^{\circ}. \]
\end{enumerate}
\end{corrige}

\begin{corrige}[Exercice 6 — Triangle inscrit dans un demi-cercle]
\begin{enumerate}
\item Le triangle $ABC$ est inscrit dans un cercle dont le côté $[BC]$ est un \cle{diamètre}. Or, \textbf{si un triangle est inscrit dans un cercle dont l'un des côtés est un diamètre, alors ce triangle est rectangle}, et le diamètre est son hypoténuse. Donc $ABC$ est \textbf{rectangle en $A$}.
\item D'après le théorème de Pythagore dans le triangle $ABC$ rectangle en $A$ :
\begin{align*}
BC^{2} &= AB^{2} + AC^{2} \\
AC^{2} &= BC^{2} - AB^{2} = 13^{2} - 5^{2} = 169 - 25 = 144 \\
AC &= \sqrt{144} = 12~\text{cm}.
\end{align*}
\textbf{Vérification :} $5^{2} + 12^{2} = 25 + 144 = 169 = 13^{2}$. \checkmark
\end{enumerate}
\end{corrige}

\begin{corrige}[Exercice 7 — Chasse aux angles]
\begin{enumerate}
\item $\widehat{ADB}$ et $\widehat{ACB}$ sont des angles inscrits interceptant l'arc $\overset{\frown}{AB}$, associé à l'angle au centre $\widehat{AOB} = 80^{\circ}$. D'après le théorème de l'angle inscrit :
\[ \widehat{ADB} = \widehat{ACB} = \frac{80^{\circ}}{2} = 40^{\circ}. \]
\item $\widehat{BAC}$ et $\widehat{BDC}$ interceptent l'arc $\overset{\frown}{BC}$, associé à l'angle au centre $\widehat{BOC} = 60^{\circ}$. Donc :
\[ \widehat{BAC} = \widehat{BDC} = \frac{60^{\circ}}{2} = 30^{\circ}. \]
\item On constate que $\widehat{ADB} = \widehat{ACB} = 40^{\circ}$. C'est normal : ce sont deux angles inscrits qui interceptent le \emph{même arc} $\overset{\frown}{AB}$ (les points $C$ et $D$ sont sur le même arc d'extrémités $A$ et $B$), donc ils sont égaux d'après la propriété des angles inscrits interceptant le même arc.
\end{enumerate}
\end{corrige}

\begin{corrige}[Exercice 8 — Réciproque du demi-cercle]
\begin{enumerate}
\item Le triangle $MNP$ est rectangle en $M$. D'après le théorème de Pythagore :
\begin{align*}
NP^{2} &= MN^{2} + MP^{2} = 6^{2} + 8^{2} = 36 + 64 = 100 \\
NP &= \sqrt{100} = 10~\text{cm}.
\end{align*}
\item \textbf{Propriété (réciproque) :} si un triangle est rectangle, alors son cercle circonscrit a pour diamètre son hypoténuse. Ici $MNP$ est rectangle en $M$, donc son hypoténuse est $[NP]$ : le cercle circonscrit a pour diamètre $[NP]$. Son rayon est :
\[ r = \frac{NP}{2} = \frac{10}{2} = 5~\text{cm}. \]
\end{enumerate}
\end{corrige}

\begin{corrige}[Exercice 9 — Type BEPC : angles d'un triangle inscrit]
\begin{enumerate}
\item L'angle inscrit $\widehat{BAD}$ et l'angle au centre $\widehat{BOD}$ interceptent le même arc $\overset{\frown}{BD}$. D'après le théorème de l'angle inscrit :
\[ \widehat{BAD} = \frac{\widehat{BOD}}{2} = \frac{80^{\circ}}{2} = 40^{\circ}. \]
\item De même, $\widehat{CAD}$ et l'angle au centre $\widehat{COD}$ interceptent le même arc $\overset{\frown}{CD}$ :
\[ \widehat{CAD} = \frac{\widehat{COD}}{2} = \frac{50^{\circ}}{2} = 25^{\circ}. \]
\item Les points $B$ et $C$ étant du même côté du diamètre $[AD]$ avec $B$ entre $A$ et $C$ sur l'arc, la demi-droite $[AD)$ est intérieure à l'angle $\widehat{BAC}$, donc :
\[ \widehat{BAC} = \widehat{BAD} + \widehat{CAD} = 40^{\circ} + 25^{\circ} = 65^{\circ}. \]
\item Comme $[AD]$ est un diamètre, les points $A$, $O$, $D$ sont alignés, donc :
\[ \widehat{AOB} = 180^{\circ} - \widehat{BOD} = 180^{\circ} - 80^{\circ} = 100^{\circ} \quad \text{et} \quad \widehat{AOC} = 180^{\circ} - \widehat{COD} = 180^{\circ} - 50^{\circ} = 130^{\circ}. \]
On en déduit :
\[ \widehat{BOC} = \widehat{AOC} - \widehat{AOB} = 130^{\circ} - 100^{\circ} = 30^{\circ}. \]
\textbf{Vérification de cohérence :} $\widehat{AOB} + \widehat{BOC} + \widehat{COD} = 100^{\circ} + 30^{\circ} + 50^{\circ} = 180^{\circ}$, ce qui correspond bien au demi-cercle de diamètre $[AD]$. \checkmark

L'angle inscrit $\widehat{ABC}$ intercepte l'arc $\overset{\frown}{AC}$ ne contenant pas $B$, dont l'angle au centre associé est $\widehat{AOC} = 130^{\circ}$. Donc :
\[ \widehat{ABC} = \frac{\widehat{AOC}}{2} = \frac{130^{\circ}}{2} = 65^{\circ}. \]
\item L'angle inscrit $\widehat{ACB}$ intercepte l'arc $\overset{\frown}{AB}$, d'angle au centre $\widehat{AOB} = 100^{\circ}$, donc :
\[ \widehat{ACB} = \frac{\widehat{AOB}}{2} = \frac{100^{\circ}}{2} = 50^{\circ}. \]
Alors :
\[ \widehat{BAC} + \widehat{ABC} + \widehat{ACB} = 65^{\circ} + 65^{\circ} + 50^{\circ} = 180^{\circ}. \]
La somme des angles du triangle $ABC$ vaut bien $180^{\circ}$. \checkmark
\end{enumerate}
\textbf{Barème indicatif (sur 10) :} Q1 : 2 pts (1 pt résultat + 1 pt justification) ; Q2 : 1,5 pt ; Q3 : 1,5 pt ; Q4 : 2,5 pts ; Q5 : 2,5 pts.

\textbf{Points de vigilance :}
\begin{itemize}
\item Citer explicitement le théorème de l'angle inscrit à chaque utilisation (« l'angle inscrit mesure la moitié de l'angle au centre interceptant le même arc »).
\item Bien préciser quel arc est intercepté : c'est la condition d'application du théorème.
\item Ne pas confondre $\widehat{BOC}$ avec $\widehat{BAC}$ : l'un est un angle au centre, l'autre un angle inscrit.
\item Pour calculer $\widehat{BOC}$, utiliser $\widehat{AOC} - \widehat{AOB}$ (et non $180^{\circ} - \widehat{BOD} - \widehat{COD}$, qui ne correspond pas à la configuration).
\end{itemize}
\end{corrige}

\begin{corrige}[Exercice 10 — Problème concret : la piste d'athlétisme]
\begin{enumerate}
\item Figure codée :
\begin{popfigure}
\begin{tikzpicture}[scale=1]
\draw[popblue, thick] (0,0) circle (2);
\coordinate (O) at (0,0);
\coordinate (A) at (115:2);
\coordinate (B) at (245:2);
\coordinate (S) at (0:2);
\coordinate (T) at (180:2);
\draw[popdark] (O)--(A) (O)--(B);
\draw[poporange, thick] (S)--(A) (S)--(B);
\draw[popgreen, thick] (T)--(A) (T)--(B);
\draw[popdark] (115:0.5) arc (115:245:0.5);
\node[popdark] at (180:0.75) {\small $130^{\circ}$};
\foreach \p/\pos in {A/above left, B/below left, S/right, T/left, O/below right}
  \fill[popink] (\p) circle (1.5pt) node[\pos] {$\p$};
\end{tikzpicture}
\legende{Le spectateur $S$ est sur l'arc majeur, $T$ sur l'arc mineur $\overset{\frown}{AB}$.}
\end{popfigure}
\item L'angle $\widehat{ASB}$ est un angle inscrit interceptant l'arc mineur $\overset{\frown}{AB}$ (car $S$ est sur l'arc majeur), associé à l'angle au centre $\widehat{AOB} = 130^{\circ}$. D'après le théorème de l'angle inscrit :
\[ \widehat{ASB} = \frac{\widehat{AOB}}{2} = \frac{130^{\circ}}{2} = 65^{\circ}. \]
Le spectateur voit la ligne d'arrivée sous un angle de $65^{\circ}$.
\item Le spectateur en $T$ est sur l'arc mineur : l'angle inscrit $\widehat{ATB}$ intercepte l'arc \emph{majeur} $\overset{\frown}{AB}$, dont l'angle au centre mesure $360^{\circ} - 130^{\circ} = 230^{\circ}$. Donc :
\[ \widehat{ATB} = \frac{230^{\circ}}{2} = 115^{\circ}. \]
\item Non : $\widehat{ASB} = 65^{\circ} \neq \widehat{ATB} = 115^{\circ}$. Les deux angles sont \textbf{supplémentaires} ($65^{\circ} + 115^{\circ} = 180^{\circ}$). Deux spectateurs voient la ligne d'arrivée sous le même angle uniquement s'ils sont sur le même arc ; c'est pourquoi, sur une tribune circulaire située du même côté, tous les spectateurs ont la même ouverture de vue sur l'arrivée.
\end{enumerate}
\textbf{Points de vigilance :} bien identifier l'arc intercepté selon la position du spectateur ; penser à utiliser $360^{\circ} - 130^{\circ}$ pour l'arc majeur.
\end{corrige}

\begin{corrige}[Exercice 11 — Démonstration et synthèse]
\begin{enumerate}
\item Les angles $\widehat{MPN}$ et $\widehat{MQN}$ sont deux angles inscrits dans le cercle $(\mathcal{C})$ qui interceptent le même arc $\overset{\frown}{MN}$. Or, \textbf{deux angles inscrits interceptant le même arc sont égaux}. Donc :
\[ \widehat{MPN} = \widehat{MQN}. \]
\item L'angle au centre $\widehat{MON} = 96^{\circ}$ est associé aux angles inscrits interceptant l'arc $\overset{\frown}{MN}$. D'après le théorème de l'angle inscrit :
\[ \widehat{MPN} = \widehat{MQN} = \frac{\widehat{MON}}{2} = \frac{96^{\circ}}{2} = 48^{\circ}. \]
\item Le triangle $MNP$ est inscrit dans un demi-cercle de diamètre $[NP]$ : d'après la propriété du triangle inscrit dans un demi-cercle, $MNP$ est \textbf{rectangle en $M$}. Donc :
\[ \widehat{NMP} = 90^{\circ}. \]
\item Dans les triangles $IMP$ et $IQN$ :
\begin{itemize}
\item $\widehat{IPM} = \widehat{IQN}$, car ce sont les angles $\widehat{MPN}$ et $\widehat{MQN}$ (les points $I$, $M$, $Q$ d'une part et $I$, $P$, $N$ d'autre part sont alignés), et ils sont égaux d'après la question 1 ;
\item $\widehat{MIP} = \widehat{NIQ}$, car ce sont des \textbf{angles opposés par le sommet} $I$.
\end{itemize}
Les triangles $IMP$ et $IQN$ ont donc deux angles respectivement égaux. Comme la somme des angles d'un triangle vaut $180^{\circ}$, leurs troisièmes angles sont aussi égaux :
\[ \widehat{IMP} = 180^{\circ} - \widehat{IPM} - \widehat{MIP} = 180^{\circ} - \widehat{IQN} - \widehat{NIQ} = \widehat{INQ}. \]
\textbf{Conclusion :} les triangles $IMP$ et $IQN$ ont leurs trois angles respectivement égaux (ils ont la même forme).
\end{enumerate}
\textbf{Barème indicatif (sur 10) :} Q1 : 2 pts ; Q2 : 2 pts ; Q3 : 2 pts ; Q4 : 4 pts.

\textbf{Points de vigilance :}
\begin{itemize}
\item À la question 4, justifier l'égalité $\widehat{MIP} = \widehat{NIQ}$ par « angles opposés par le sommet », pas seulement par la figure.
\item Pour conclure sur les troisièmes angles, invoquer la somme des angles d'un triangle ($180^{\circ}$).
\end{itemize}
\end{corrige}

\newpage

\coursec{Fiche bilan}

\begin{popfigure}
\begin{tikzpicture}[mindmap, grow cyclic, every node/.style={concept, font=\small},
  root concept/.append style={concept color=popblue, font=\bfseries, minimum size=2.6cm},
  level 1/.append style={level distance=3.4cm, sibling angle=60},
  level 1 concept/.append style={font=\footnotesize, minimum size=2.1cm}]
\node[root concept]{Angles inscrits}
  child[concept color=poporange]{ node {Définitions : angle inscrit, angle au centre, arc intercepté} }
  child[concept color=popgreen]{ node {Théorème : angle au centre $=2\times$ angle inscrit} }
  child[concept color=poppurple]{ node {Même arc $\Rightarrow$ angles inscrits égaux} }
  child[concept color=poppink]{ node {Triangle inscrit dans un demi-cercle = rectangle} }
  child[concept color=popgold]{ node {Calculs d'angles et démonstrations} }
  child[concept color=popblue!60]{ node {Vie courante : stades, roues, navigation} };
\end{tikzpicture}
\legende{Carte mentale de la leçon : les angles inscrits.}
\end{popfigure}

\begin{bilan}
\textbf{Définitions clés}
\begin{itemize}
  \item Un \cle{angle inscrit} dans un cercle est un angle dont le sommet est \textbf{sur le cercle} et dont les côtés coupent le cercle : $\widehat{AMB}$ avec $A$, $M$, $B$ sur le cercle.
  \item L'\cle{angle au centre} associé a son sommet au \textbf{centre} $O$ du cercle : $\widehat{AOB}$.
  \item On dit que ces deux angles \cle{interceptent le même arc} $\overset{\frown}{AB}$ (celui qui ne contient pas $M$).
\end{itemize}

\textbf{Théorèmes à connaître par c\oe ur}
\begin{center}
\begin{tabularx}{\linewidth}{|l|X|}
\hline
\rowcolor{popblueL} \textbf{Propriété} & \textbf{Énoncé} \\
\hline
Angle inscrit / angle au centre & Si $\widehat{AMB}$ et $\widehat{AOB}$ interceptent le même arc, alors $\widehat{AOB}=2\times\widehat{AMB}$, c'est-à-dire $\widehat{AMB}=\dfrac{\widehat{AOB}}{2}$. \\
\hline
Angles inscrits interceptant le même arc & Deux angles inscrits qui interceptent le \textbf{même arc} sont \textbf{égaux} : $\widehat{AMB}=\widehat{ANB}$. \\
\hline
Triangle et demi-cercle & Si un triangle est inscrit dans un cercle et qu'un de ses côtés est un \textbf{diamètre}, alors ce triangle est \textbf{rectangle} (l'angle droit est opposé au diamètre). Réciproquement, l'hypoténuse d'un triangle rectangle est un diamètre de son cercle circonscrit. \\
\hline
\end{tabularx}
\end{center}

\textbf{Méthodes express}
\begin{itemize}
  \item Pour calculer un angle inscrit : repérer l'arc intercepté, trouver l'angle au centre associé, puis \textbf{diviser par 2}.
  \item Pour prouver que deux angles sont égaux : vérifier qu'ils sont inscrits dans le \textbf{même cercle} et interceptent le \textbf{même arc}.
  \item Pour prouver qu'un triangle est rectangle : montrer qu'il est inscrit dans un cercle dont un côté du triangle est un \textbf{diamètre}.
  \item Toujours nommer l'arc intercepté dans la rédaction : c'est lui qui justifie la propriété utilisée.
\end{itemize}
\end{bilan}

\begin{aretenir}[Checklist avant le BEPC]
\begin{itemize}
  \item[$\square$] Je sais reconnaître un angle inscrit et l'angle au centre associé sur une figure.
  \item[$\square$] Je sais nommer l'arc intercepté par un angle inscrit.
  \item[$\square$] Je connais par c\oe ur : $\widehat{AOB}=2\times\widehat{AMB}$ pour le même arc.
  \item[$\square$] Je sais calculer un angle inscrit à partir de l'angle au centre (et inversement).
  \item[$\square$] Je sais que deux angles inscrits interceptant le même arc sont égaux.
  \item[$\square$] Je sais repérer un triangle inscrit dans un demi-cercle et conclure qu'il est rectangle.
  \item[$\square$] Je sais que le cercle circonscrit à un triangle rectangle a pour diamètre l'hypoténuse.
  \item[$\square$] Je vérifie toujours que les points sont bien \textbf{sur le même cercle} avant d'appliquer les propriétés.
  \item[$\square$] Je rédige en citant la propriété utilisée et l'arc intercepté.
  \item[$\square$] Je fais une figure propre (cercle tracé au compas, points nommés) pour vérifier mes calculs.
\end{itemize}
\end{aretenir}

\findecours

\end{document}
